%-----------------------------------------------------------------------------
% Master Prefix Router for Randomized Algorithms Class Notes & Book
%
% Dispatches to:
%   - prefix_lwarp.tex when warping/HTML mode is active (lwarpmk html or \warpMode)
%   - prefix_4ht.tex when TeX4ht/make4ht defines \HCode
%   - prefix_latex.tex for standard LaTeX / XeLaTeX / LuaLaTeX PDF builds
% Each backend loads prefix_common.tex for shared notation and reference helpers.
%-----------------------------------------------------------------------------

% Include only once; close the conditional before ending this input.
\ifdefined\PrefixLoadedAlready\expandafter\endinput\fi

\def\PrefixLoadedAlready{}


\makeatletter
\def\input@path{{shared/}}
\makeatother


\makeatletter
\def\LWR@isWarp{0}%
\ifx\warpMode\undefined\else\def\LWR@isWarp{1}\fi
\expandafter\ifx\csname opt@lwarp.sty\endcsname\relax\else\def\LWR@isWarp{1}\fi

\if1\LWR@isWarp
  \def\warpMode{1}%
  \IfFileExists{shared/prefix_lwarp.tex}{%-----------------------------------------------------------------------------
% Master Prefix File for Lwarp (HTML Web Generation)
% Compatible with LuaLaTeX and Lwarp
%-----------------------------------------------------------------------------

\RequirePackage{iftex}
\ifLuaTeX
  % LuaLaTeX verified
\else
  \PackageError{lwarp-gate}{%
    prefix_lwarp requires LuaLaTeX! Please run with lualatex.%
  }{HTML generation via lwarp must be run with LuaLaTeX.}
  \stop
\fi

% Always use memoir
\documentclass[12pt,oneside,fleqn]{memoir}

% Load lwarp immediately after documentclass
\usepackage[mathjax]{lwarp}

% Lwarp settings: split at chapter level (0) with section names
\setcounter{FileDepth}{0}
\booltrue{FileSectionNames}
\setcounter{SideTOCDepth}{1}

\makeatletter
% Multi-page SVG resolving hook for Ipe drawings with page=N
\def\LWR@current@figpage{}
\define@key{igraph}{page}{\def\LWR@current@figpage{#1}}
% Perpage compatibility for auxiliary files
\providecommand{\pp@pagectr}[4]{}
\AtBeginDocument{
  \ifdef{\LWR@nullfonts}{%
    \appto\LWR@nullfonts{%
      \def\xspace{ }%
      \def\Turan{Turán}%
      \def\Lovasz{Lovász}%
      \def\Erdos{Erdős}%
      \def\Holder{Hölder}%
      \def\Matousek{Matoušek}%
    }%
  }{}%
  \ifdef{\LWR@replaceEPSSVG}{
    \appto\LWR@replaceEPSSVG{%
      \ifx\LWR@current@figpage\@empty\else
        \edef\LWR@tempone{\LWR@tempone_p\LWR@current@figpage}%
        \def\LWR@current@figpage{}%
      \fi
    }%
  }{}%
}
\makeatother

%-----------------------------------------------------------------------------
% Typography, Fonts & Math Packages
%-----------------------------------------------------------------------------
\usepackage{mathtools}
\PassOptionsToPackage{warnings-off={mathtools-overbracket,mathtools-colon}}{unicode-math}
\usepackage{fontspec}
\usepackage{unicode-math}

% New Computer Modern OpenType text and math fonts
\setmainfont{NewCM10-Book.otf}[
  BoldFont          = NewCM10-Bold.otf,
  ItalicFont        = NewCM10-BookItalic.otf,
  BoldItalicFont    = NewCM10-BoldItalic.otf,
  SlantedFont       = NewCM10-Book.otf,
  SlantedFeatures   = {FakeSlant = 0.25},
  BoldSlantedFont   = NewCM10-Bold.otf,
  BoldSlantedFeatures = {FakeSlant = 0.25},
  SmallCapsFeatures = {Letters = SmallCaps}
]
\setsansfont{NewCMSans10-Book.otf}[
  BoldFont          = NewCMSans10-Bold.otf,
  ItalicFont        = NewCMSans10-BookOblique.otf,
  BoldItalicFont    = NewCMSans10-BoldOblique.otf
]
\setmonofont{NewCMMono10-Book.otf}[
  BoldFont          = NewCMMono10-Bold.otf,
  ItalicFont        = NewCMMono10-BookItalic.otf,
  BoldItalicFont    = NewCMMono10-BoldOblique.otf
]

\setmathfont{NewCMMath-Book.otf}
\setmathfont{NewCMMath-Bold.otf}[version=bold]

% OpenType math aliases and fixes
\providecommand{\square}{\mdlgwhtsquare}
\def\Box{\mdlgwhtsquare}
\def\Diamond{\mdlgwhtdiamond}
\def\blacksquare{\mdlgblksquare}
\def\leadsto{\rightsquigarrow}
\providecommand{\llbracket}{\lBrack}
\providecommand{\rrbracket}{\rBrack}
\DeclareTextCommand{\textcircled}{TU}[1]{%
  \mbox{\ooalign{\hfil\raise0.1ex\hbox{#1}\hfil\crcr\hbox{\ensuremath{\bigcirc}}}}%
}

\usepackage{alphabeta}
\usepackage[normalem]{ulem} % Provides \sout
\usepackage{mleftright}
\usepackage{microtype}
\usepackage{xspace}
\usepackage{xurl}
\usepackage[dvipsnames,table]{xcolor}
\usepackage{graphicx}
\usepackage{caption}
\usepackage{booktabs,array,colortbl,hhline,multirow,tabularx}
\usepackage[inline]{enumitem}
\usepackage[thicklines]{cancel}
\renewcommand\CancelColor{\color{purple}}
\usepackage{euscript}
\usepackage{marvosym}
\let\Cross\undefined
\let\Square\undefined
\usepackage{bbding}
\usepackage{epigraph}
\usepackage{makeidx}
\usepackage{perpage}
\usepackage{tikz}

%-----------------------------------------------------------------------------
% Boondox / mathcalb font declaration
%-----------------------------------------------------------------------------
\DeclareFontFamily{U}{BOONDOX-calo}{\skewchar\font=45 }
\DeclareFontShape{U}{BOONDOX-calo}{m}{n}{
  <-> s*[1.05] BOONDOX-r-calo}{}
\DeclareFontShape{U}{BOONDOX-calo}{b}{n}{
  <-> s*[1.05] BOONDOX-b-calo}{}
\DeclareMathAlphabet{\mathcalb}{U}{BOONDOX-calo}{m}{n}
\SetMathAlphabet{\mathcalb}{bold}{U}{BOONDOX-calo}{b}{n}
\DeclareMathAlphabet{\mathbcalb}{U}{BOONDOX-calo}{b}{n}

%-----------------------------------------------------------------------------
% Algorithm Environment (formerly salgorithm.sty)
%-----------------------------------------------------------------------------
\newcommand{\ProgWordX}[1]{{\small\textbf{#1}}{\ }}
\newcommand{\Do}{\ProgWordX{do}\xspace}
\newcommand{\End}{\ProgWordX{end}\xspace}
\newcommand{\Return}{\ProgWordX{return}\xspace}%
\newcommand{\Procbegin}{\ProgWordX{begin}}
\newcommand{\If}{\ProgWordX{if}\xspace}
\newcommand{\Then}{\ProgWordX{then}\xspace}
\newcommand{\Else}{\ProgWordX{else}\xspace}
\newcommand{\For}{\ProgWordX{for}\xspace}
\newcommand{\While}{{\small\bfseries while}\ }

\newcommand{\xbeginlgox}{\begin{minipage}{1in}\begin{tabbing}
           \quad\=\qquad\=\qquad\=\qquad\=\qquad\=\qquad\=\qquad\=\kill}
\newcommand{\xendlgox}{\end{tabbing}\end{minipage}}

\newenvironment{algorithmEnv}{%
   \begin{tabular}{|l|}\hline\xbeginlgox\bigskip}%
   {\xendlgox\\\hline\end{tabular}}

\newenvironment{program}{%
   \begin{minipage}{4.0in}%
   \begin{tabbing}%
       \ \ \ \ \=\ \ \ \=\ \ \ \ \=\ \ \ \ \=\ \ \ \ \=%
      \ \ \ \ \=\ \ \ \ \=\ \ \ \ \=\ \ \ \ \=%
      \ \ \ \ \=\ \ \ \ \=\ \ \ \ \=\ \ \ \ \=\kill%
}{%
   \end{tabbing}%
   \end{minipage}%
}

\newcommand{\AlgorithmBox}[1]{%
   \fbox{\begin{program} #1 \end{program}}%
}

\definecolor{AlgorithmColor}{cmyk}{0.07,0.90,0,0.34}
\providecommand{\Algorithm}[1]{{\AlgorithmI{#1}\index{algorithm!#1@{\AlgorithmI{#1}}}}}
\newcommand{\CodeComment}[1]{\textcolor{blue}{\texttt{#1}}}

%-----------------------------------------------------------------------------
% Path & Directory Resolution
%-----------------------------------------------------------------------------
\newcommand{\RootDIR}{../}
\IfFileExists{../../../book.tex}{\renewcommand{\RootDIR}{../../../}}{}
\IfFileExists{../../book.tex}{\renewcommand{\RootDIR}{../../}}{}
\IfFileExists{../book.tex}{\renewcommand{\RootDIR}{../}}{}
\IfFileExists{./book.tex}{\renewcommand{\RootDIR}{}}{}

\providecommand{\File}[1]{#1}
\providecommand{\RootFile}[1]{\RootDIR #1}

%-----------------------------------------------------------------------------
% QED Duck Symbol
%-----------------------------------------------------------------------------
\newlength{\CharHeight}
\AtBeginDocument{\setlength{\CharHeight}{\fontcharht\font`X}}

\IfFileExists{shared/qed_duck.svg}{%
   \providecommand{\myqedsymbol}{%
      \reflectbox{\includegraphics[height=1.4\CharHeight]{shared/qed_duck}}%
   }%
}{%
   \providecommand{\myqedsymbol}{$\blacksquare$}%
}

%-----------------------------------------------------------------------------
% BibLaTeX Setup (Pure Biber, No BibTeX)
%-----------------------------------------------------------------------------
\usepackage[bibencoding=utf8,style=alphabetic,refsegment=chapter,
indexing=cite,
backend=biber]{biblatex}

\ExecuteBibliographyOptions{%
   maxnames = 40,%
   maxalphanames = 4,%
   maxcitenames = 3,%
   minalphanames = 3,%
   minnames = 1,%
   giveninits = true,%
   doi = false,%
   url = false,%
   isbn = false%
}

\renewbibmacro{in:}{}

\AtEveryBibitem{%
   \ifentrytype{inproceedings}{%
      \clearname{editor}%
      \clearlist{editor}%
      \clearfield{editor}%
      \clearlist{publisher}%
      \clearfield{series}%
      \clearlist{location}%
      \clearfield{publisher}%
   }{}%
   \ifentrytype{book}{%
      \clearname{publisher}%
   }{}%
}
\AtEveryCitekey{%
   \ifentrytype{inproceedings}{%
      \clearlist{publisher}%
   }{}%
   \clearname{editor}%
   \clearlist{location}%
}
\DeclareFieldFormat[inproceedings]{date}{}
\DeclareFieldFormat[inproceedings]{pages}{#1, \printfield{year}}
\DeclareFieldFormat[inproceedings]{booktitle}{\mkbibemph{#1,\nopunct}}

\DeclareFieldFormat{journaltitle}{\mkbibemph{#1}\isdot}

\DeclareFieldFormat[article]{number}{\mkbibparens{#1}}
\DeclareFieldFormat[article]{date}{}
\DeclareFieldFormat[article]{pages}{#1, \printfield{year}}
\DeclareFieldFormat[article]{issue}{\nopunct#1}
\renewbibmacro*{issue+date}{%
   \iffieldundef{issue}{}%
   {%
      \printtext{\unspace}%
      (\printfield{issue})%
   }%
   \iffieldundef{pages}%
   {\addcomma\printfield{year}}%
   {}%
}

\renewcommand{\bibpagespunct}{%
   \ifentrytype{article}{\addcolon}{\addcomma}\space%
}

\DeclareFieldFormat[article]{titlecase}{\MakeSentenceCase{#1}}
\DeclareFieldFormat[inproceedings]{titlecase}{\MakeSentenceCase{#1}}
\DeclareFieldFormat[incollection]{titlecase}{\MakeSentenceCase{#1}}
\DeclareFieldFormat[book]{titlecase}{\MakeSentenceCase{#1}}

\renewbibmacro*{journal}{%
   \iffieldundef{journaltitle} {} {\printtext[journaltitle]{%
         \printfield[myplain]{journaltitle}%
         \setunit{\subtitlepunct}%
         \printfield[myplain]{journalsubtitle}}}}
\DeclareFieldFormat{myplain}{#1}

\renewbibmacro*{booktitle}{%
   \ifboolexpr{ test {\iffieldundef{booktitle}} and test
      {\iffieldundef{booksubtitle}} } {} {\printtext[booktitle]{%
         \printfield[myplain]{booktitle}%
         \setunit{\subtitlepunct}%
         \printfield[myplain]{booksubtitle}}%
      \newunit}%
   \printfield{booktitleaddon}}

\renewbibmacro*{volume+number+eid}{%
   \setunit*{\addcomma\addspace}%
   \printfield{volume}%
   \printfield{number}%
   \setunit{\addcomma\space}%
   \printfield{eid}%
}

\newbibmacro{string+doi}[1]{%
   \iffieldundef{doi}{%
      \iffieldundef{url}{%
         #1%
      }{%
         \href{\thefield{url}}{#1}%
      }%
   }{%
      \href{http://dx.doi.org/\thefield{doi}}{#1}}%
}
\DeclareFieldFormat*{title}{\usebibmacro{string+doi}{\mkbibemph{#1}}}
\renewcommand*{\multicitedelim}{\addcomma\space}
\DeclareFieldFormat[inproceedings]{title}{\usebibmacro{string+doi}{{#1}}}
\DeclareFieldFormat[article]{title}{\usebibmacro{string+doi}{{#1}}}
\DeclareFieldFormat[incollection]{title}{\usebibmacro{string+doi}{{#1}}}
\DeclareFieldFormat[book]{titlecase}{{#1}}

\DeclareCiteCommand{\tcite}%
{\usebibmacro{prenote}}%
{\usebibmacro{citeindex}\usebibmacro{cite}}%
{\multicitedelim}%
{\usebibmacro{postnote}}

%-----------------------------------------------------------------------------
% Hyperref & Visual Formatting
%-----------------------------------------------------------------------------
\usepackage{hyperref}
\hypersetup{%
   breaklinks,%
   colorlinks=true,%
   linkcolor=[rgb]{0.45,0.0,0.0},%
   citecolor=[rgb]{0,0,0.45}%
}

\definecolor{DeepNavy}{HTML}{002366}
\definecolor{RoyalBlue}{HTML}{004E92}
\definecolor{SteelBlue}{HTML}{4682B4}
\definecolor{TealBlue}{HTML}{008080}
\definecolor{DarkSlate}{HTML}{2F4F4F}
\definecolor{DustyGrey}{HTML}{696969}

\makeatletter
\AtBeginDocument{\let\@elt\relax}
\makeatother

\headstyles{bringhurst}
\chapterstyle{dash}

\renewcommand{\chapnamefont}{\normalfont\huge\bfseries\color{DeepNavy}}
\renewcommand{\chapnumfont}{\normalfont\huge\bfseries\color{DeepNavy}}
\renewcommand{\chaptitlefont}{\normalfont\Huge\bfseries\color{DeepNavy}}

\counterwithin{footnote}{chapter}

\setsubparaheadstyle{\normalfont\normalsize\itshape\color{DustyGrey}}
\setlength{\beforechapskip}{-20pt} 
\setlength{\afterchapskip}{30pt}
\setsecnumformat{\csname the#1\endcsname.\quad}

%-----------------------------------------------------------------------------
% Standalone Chapter vs Master Book Modes
%-----------------------------------------------------------------------------
\newcommand{\Copyright}{This work is licensed under the Creative Commons Attribution-Noncommercial 3.0 License. To view a copy of this license, visit \url{http://creativecommons.org/licenses/by-nc/3.0/} or send a letter to Creative Commons, 171 Second Street, Suite 300, San Francisco, California, 94105, USA.}

\newcommand{\RealChapterEndInner}{%
   {%
      \begingroup%
      \emergencystretch=1em%
      \printbibliography[heading=subbibliography,segment=\therefsegment]{}%
      \endgroup%
   }%
}

\providecommand{\BookMode}[1]{}
\providecommand{\StandAloneMode}[1]{#1}

\ifx\bookMode\undefined
   % STANDALONE CHAPTER MODE
   \renewcommand{\StandAloneMode}[1]{#1}
   \renewcommand{\BookMode}[1]{}

   \newcommand{\ChapterNumPage}[2]{%
      \if\relax\detokenize{#1}\relax
         \setcounter{chapter}{0}%
      \else\ifnum#1>0
         \setcounter{chapter}{\numexpr#1-1\relax}%
      \else
         \setcounter{chapter}{0}%
      \fi\fi
      \if\relax\detokenize{#2}\relax
         \setcounter{page}{1}%
      \else\ifnum#2>0
         \setcounter{page}{#2}%
      \else
         \setcounter{page}{1}%
      \fi\fi
   }

   \DeclareRobustCommand{\setheadingcolor}{\color{RoyalBlue}}
   \setsecheadstyle{\setheadingcolor\normalfont\Large\bfseries}
   \setsubsecheadstyle{\normalfont\large\bfseries\color{SteelBlue}}
   \setsubsubsecheadstyle{\normalfont\normalsize\bfseries\color{TealBlue}}
   \setparaheadstyle{\normalfont\normalsize\bfseries\color{DarkSlate}}

   \DeclareDocumentCommand{\Chapter}{s o m}{%
      \IfBooleanTF{#1}{%
         \IfValueTF{#2}{\chapter*[]{#2}}{\chapter*[]{#3}}%
      }{%
         \IfValueTF{#2}{\chapter*[]{#2}}{\chapter[]{#3}}%
      }%
      \setcounter{footnote}{0}%
   }
   \newcommand{\ChapterPrefix}{}
   \newcommand{\ChapterGen}{\pagestyle{plain}\vspace*{-0.7cm}}
   \newcommand{\ChapterEnd}{%
      \RealChapterEndInner{}%
      \end{document}%
   }

   \renewcommand{\File}[1]{#1}
   \renewcommand{\RootFile}[1]{../#1}
\else
   % MASTER BOOK MODE
   \renewcommand{\StandAloneMode}[1]{}
   \renewcommand{\BookMode}[1]{#1}

   \newcommand{\ChapterNumPage}[2]{}

   \DeclareDocumentCommand{\Chapter}{s o m}{%
      \IfBooleanTF{#1}{%
         \IfValueTF{#2}{\chapter*[#2]{#3}}{\chapter*{#3}}%
         \ifx\Dir\undefined\else
            \typeout{CHAPTER_MAP: \Dir/\CurrentChapFile => chap: 0, page: \thepage}%
         \fi
      }{%
         \IfValueTF{#2}{\chapter[#2]{#3}}{\chapter{#3}}%
         \ifx\Dir\undefined\else
            \typeout{CHAPTER_MAP: \Dir/\CurrentChapFile => chap: \thechapter, page: \thepage}%
         \fi
      }%
   }
   \newcommand{\ChapterPrefix}{%
      \vspace{-1cm}%
      \subsection*{\small
         {598 - Class notes for Randomized Algorithms}\\
         {Sariel Har-Peled}\\
         {\today}}%
   }
   \newcommand{\ChapterGen}{}
   \newcommand{\ChapterEnd}{\RealChapterEndInner{}}

   \renewcommand{\File}[1]{\Dir/#1}
   \renewcommand{\RootFile}[1]{#1}

   \newcommand{\InputChap}[2]{%
      \def\Dir{#1}%
      \def\CurrentChapFile{#2}%
      \typeout{}%
      \typeout{------------------------------------------------------------}%
      \typeout{PROCESSING FILE: #1/#2}%
      \typeout{}%
      \input{#1/#2}%
   }
\fi

%-----------------------------------------------------------------------------
% Fragments & Quotes
%-----------------------------------------------------------------------------
\newcommand{\InputExt}[1]{%
   \IfFileExists{#1}%
   {\input{#1}}%
   {\IfFileExists{../#1}%
      {\input{../#1}}%
      {\IfFileExists{../../#1}%
         {\input{../../#1}}%
         {\typeout{ERROR: Unable to find #1}}%
      }%
   }%
}
\makeatletter
\def\FragNormKey#1{%
  \frag@norm@check#1def_\relax{#1}%
}
\def\frag@norm@check#1def_#2\relax#3{%
  \if\relax\detokenize{#1}\relax
    #3%
  \else
    def_#3%
  \fi
}
\makeatother

\newcommand{\FragmentName}[1]{shared/fragment/\FragNormKey{#1}.tex}

\newenvironment{fragment}[1]{}{}

\makeatletter
\AtBeginDocument{%
   \let\real@latex@label\label
   \def\label#1{%
      \global\@namedef{all@labels@seen@#1}{1}%
      \real@latex@label{#1}%
   }%
}

\DeclareDocumentCommand{\IncFragment}{m}{%
   \begingroup
      \def\label##1{%
         \@ifundefined{all@labels@seen@##1}{%
            \global\@namedef{all@labels@seen@##1}{1}%
            \real@latex@label{##1}%
         }{}%
      }%
      \IfFileExists{\FragmentName{#1}}%
         {\InputExt{\FragmentName{#1}}}%
         {\InputExt{shared/fragment/#1.tex}}%
   \endgroup
}
\makeatother

\newcommand{\Quote}[2]{%
   \SaveIndent
   \begin{center}%
      {\footnotesize
         \begin{quote}
             \RestoreIndent
             {#1}\\
             \hspace*{2cm}-- {#2.}
         \end{quote}%
      }%
      \index{quotation!#2}%
   \end{center}%
   \bigskip
}

\newcommand{\QuotePExt}[3]{%
   \bigskip%
   \hfill%
   \begin{minipage}{0.95\linewidth}%
       \footnotesize {#1}%

       \rule{\linewidth}{0.3mm}%
       \hfill #2, #3%
   \end{minipage}%
   \bigskip%
}

\epigraphfontsize{\small\itshape}
\setlength{\epigraphwidth}{0.6\textwidth}
\newcommand{\QuoteExt}[3]{%
   \vspace*{-1.6cm}%
   \epigraph{{\footnotesize #1}}{{#2}, {#3}}%
   \vspace*{-1.2cm}%
}

\providecommand{\delX}[1]{{\color{red}\sout{#1}}}
\providecommand{\newX}[1]{{\color{blue}#1}}
\providecommand{\chgY}[2]{{\color{red}\sout{#1}}{\color{blue}#2}}
\providecommand{\remX}[1]{{\color{purple}\textbf{[Remark: #1]}}}
\newcommand{\SpellIgnore}[1]{#1}
\newcommand{\EQ}{\mathrel{\;\;=\;\;}}
\newcommand{\LEQ}{\mathrel{\;\;\leq\;\;}}
\newcommand{\LT}{\mathrel{\;\;<\;\;}}
\newcommand{\GT}{\mathrel{\;\;<\;\;}}

\newcommand{\GEQ}{\mathrel{\;\;\geq\;\;}}

%-----------------------------------------------------------------------------
% Save a symbol that we know is going to get redefined.
\def\savesymbol#1{%
  \expandafter\let\expandafter\origsym\expandafter=\csname#1\endcsname
  \expandafter\let\csname orig#1\endcsname=\origsym
  \expandafter\let\csname#1\endcsname=\relax
}

% Restore a previously saved symbol, and rename the current one.
\def\restoresymbol#1#2{%
  \expandafter\let\expandafter\newsym\expandafter=\csname#2\endcsname
  \expandafter\global\expandafter\let\csname#1#2\endcsname=\newsym
  \expandafter\let\expandafter\origsym\expandafter=\csname orig#2\endcsname
  \expandafter\global\expandafter\let\csname#2\endcsname=\origsym
}

\usepackage{pictex}%
\savesymbol{triangle}%
\usepackage{tree}%
\restoresymbol{TreePKG}{triangle}%

% Definitions absorbed from sbook.sty
\providecommand{\MdPrc}{\nu}

\StartDefiningTabulars
\providecommand{\TwoFiguresBase}[4]{%
   \begin{tabular}{cc}
      \begin{minipage}{0.48\linewidth}
          \begin{center}
              {{{#1}}}
          \end{center}
      \end{minipage}
      &
      \begin{minipage}{0.48\linewidth}
          \begin{center}
              {{#3}}
          \end{center}
      \end{minipage}
      \\
      {#2} & {#4}
   \end{tabular}%
}

\providecommand{\TwoFigures}[6]{%
   \begin{tabular}{cc}
      \begin{minipage}{0.48\linewidth}
          \begin{center}
              {\IncludeGraphics[#1]{#2}}
          \end{center}
      \end{minipage}
      &
      \begin{minipage}{0.48\linewidth}
          \centerline{\IncludeGraphics[#4]{#5}}
      \end{minipage}
      \\
      {#3} & {#6}
   \end{tabular}%
}

\providecommand{\ThreeFigures}[6]{%
   \begin{tabular}{ccc}
        \begin{minipage}{0.30\linewidth}
            \IncludeGraphics[width=5.1cm]{#1}\\
        \end{minipage}
        &
        \begin{minipage}{0.30\linewidth}
            \IncludeGraphics[width=5.1cm]{#2}\\
        \end{minipage}
        &
        \begin{minipage}{0.30\linewidth}
            \IncludeGraphics[width=5.1cm]{#3}\\
        \end{minipage}
        \vspace{-1cm}
        \\
        {#4} & {#5}& {#6}
   \end{tabular}%
}
\StopDefiningTabulars

%-----------------------------------------------------------------------------
% Shared notation and reference helpers; backend variants follow below.
\IfFileExists{shared/prefix_common.tex}{%-----------------------------------------------------------------------------
% Shared notation and reference helpers for PDF, lwarp, and TeX4ht.
% Loaded by each backend after its package and layout setup, before ntheorem.
% Keep rendering-dependent definitions in the corresponding backend prefix.
% Add new backend-independent notation here so all outputs share it.
%-----------------------------------------------------------------------------

\providecommand{\remove}[1]{}
\providecommand{\NotAMSMode}[1]{#1}
\providecommand{\Solution}[1]{}
\providecommand{\Version}[1]{\hfill{Version: #1}}
\providecommand{\Distor}{\mathrm{dist}}
\providecommand{\Hardness}[1]{(h:#1)}
\providecommand{\QuoteOptExt}[3]{\QuoteExt{#1}{#2}{#3}}
\providecommand{\Hypercube}{\EuScript{H}}
\providecommand{\CoverTime}{\EuScript{C}}
\providecommand{\Nbr}{\Gamma}
\providecommand{\CommuteTime}[2]{\mathbf{CT}_{#1#2}}
\providecommand{\Resistance}{\mathbf{R}}
\providecommand{\EffectResist}[2]{\mathbf{R}_{#1#2}}
\providecommand{\NetworkE}{\EuScript{N}}
\providecommand{\trinom}[4]{\binom{#1}{#2 \;\; #3 \;\; #4}}
\providecommand{\Ranges}{\EuScript{R}}
\providecommand{\CClass}[1]{\ensuremath{\mathbf{#1}}}
\providecommand{\coCClass}[1]{\ensuremath{\mathrm{co-}}\CClass{#1}}
\providecommand{\CX}{\EuScript{C}}
\providecommand{\packet}{\mathbf{c}}
\providecommand{\packetA}{\mathbf{d}}
\providecommand{\mopt}{m_{\mathrm{opt}}}
\providecommand{\girth}{\mathop{\mathrm{girth}}}
\providecommand{\CR}{\mathop{\mathrm{cr}}}
\providecommand{\TRUE}{{\textcolor[named]{RedViolet}{\textsf{TRUE}}}\xspace}
\providecommand{\FALSE}{{\textcolor[named]{RedViolet}{\textsf{FALSE}}}\xspace}
\providecommand{\Part}[1]{\part{#1}}
\providecommand{\Section}[1]{\section{#1}}
\providecommand{\SubSection}[1]{\subsection{#1}}
\providecommand{\SubSubSection}[1]{\subsubsection{#1}}
\providecommand{\Exercise}[5]{%
\begin{exercise_h}[#1]%
   \exelab{#3}{}%
   #4%
\end{exercise_h}%
}
\providecommand{\si}[1]{#1}
\providecommand{\myParagraph}[1]{\noindent\textbf{#1}}
\providecommand{\DefName}[1]{#1}
\providecommand{\deflabel}[1]{#1}
\providecommand{\Description}[1]{}

%-----------------------------------------------------------------------------
% Mathematical Notations & Environments (from notations.sty)
%-----------------------------------------------------------------------------
%-----------------------------------------------------------------------------
% List environments (questions & compactenum)
%-----------------------------------------------------------------------------

\newlist{compactenumA}{enumerate}{5}%
\setlist[compactenumA]{topsep=0pt,itemsep=-1ex,partopsep=1ex,parsep=1ex,%
   label=(\Alph*)}%

\newlist{compactenumI}{enumerate}{5}%
\setlist[compactenumI]{topsep=0pt,itemsep=-1ex,partopsep=1ex,parsep=1ex,%
   label=(\Roman*)}%

\newlist{compactenuma}{enumerate}{5}%
\setlist[compactenuma]{topsep=0pt,itemsep=-1ex,partopsep=1ex,parsep=1ex,%
   label=(\alph*)}%

\newlist{compactitem}{enumerate}{5}%
\setlist[compactitem]{topsep=0pt,itemsep=-1ex,partopsep=1ex,parsep=1ex,%
   label=\textbullet}%

\newlist{compactenumi}{enumerate}{5}%
\setlist[compactenumi]{topsep=0pt,itemsep=-1ex,partopsep=1ex,parsep=1ex,%
   label=(\roman*)}%

\setcounter{secnumdepth}{5}

\newcommand{\Term}[1]{\textsf{#1}}
\newcommand{\TermI}[1]{\Term{#1}\index{#1@\Term{#1}}}

\DeclareMathAlphabet{\mathpzc}{OT1}{pzc}{m}{it}

\newcommand{\seg}{\mathsf{s}}

\newcommand{\Diam}{\mathrm{diam}}
\newcommand{\diam}{\mathrm{diam}}

\newcommand{\Copt}{C_{\mathrm{opt}}}
\newcommand{\CenterPnt}{\mathrm{cen}}

\newcommand{\parent}{\overline{\mathrm{p}}}
\newcommand{\parentX}[2][\!]{\parent\pth[#1]{#2}}

\newcommand{\FastCut}{\Algorithm{FastCut}\xspace}
\newcommand{\MinCut}{\Algorithm{MinCut}\xspace}
\newcommand{\MinCutRep}{\Algorithm{MinCutRep}\xspace}
\newcommand{\Contract}{\Algorithm{Contract}\xspace}

\newcommand{\Grid}{\mathrm{G}}
\newcommand{\ClosestPair}{\mathcal{CP}}
\newcommand{\SphereChar}{\ensuremath{\mathbb{S}}}
\newcommand{\Sphere}[1]{{\SphereChar^{(#1)}}}

\newcommand{\PntSet}{{\mathsf{P}}}%
\providecommand{\PntSetA}{{\mathsf{Q}}}%

\newcommand{\query}{\mathcalb{q}}

\newcommand{\pnt}{\mathsf{p}}
\newcommand{\pntA}{\mathsf{u}}
\newcommand{\pntB}{\mathsf{v}}
\newcommand{\pntC}{{\mathsf{x}}}
\newcommand{\pntD}{{\mathsf{y}}}
\newcommand{\pntE}{{\mathsf{z}}}

\newcommand{\pntCoord}{p}
\newcommand{\pntCoordA}{q}

%%%%%%%%%%%%%%%%%%%%%%%%%%%%%%%%%%%%%%%%%%%%%%%%%%%%%%%%%%%%%%%%%%%%% 
%%%%%%%%%%%%%%%%%%%%%%%%%%%%%%%%%%%%%%%%%%%%%%%%%%%%%%%%%%%%%%%%%%%%%
% Markov chain notations.

\newcommand{\MCSpc}{\mathtt{S}}
\newcommand{\MCPrb}{\mathtt{P}}
\newcommand{\MCPrbY}[2]{\MCPrb_{#1#2}}
\newcommand{\MCPrbM}{\mathbf{P}}%
\newcommand{\MCPrbT}[3]{\MCPrb_{#1#2}^{(#3)}}
\newcommand{\MCFPrbT}[3]{\mathtt{r}_{#1#2}^{(#3)}}
\newcommand{\MCVstProb}[2]{\mathtt{f}_{#1#2}}
\newcommand{\MCExSteps}[2]{\mathtt{h}_{#1#2}}
\newcommand{\MCPVec}[1]{\mathbf{q}^{(#1)}}
%%%%%%%%%%%%%%%%%%%%%%%%%%%%%%%%%%%%%%%%%%%%%%%%%%%%%%%%%%%%%%%%%%%%%

\newcommand{\Spread}{\Phi}
\newcommand{\VC}{\ensuremath{\mathsf{VC}}\xspace}
\newcommand{\MC}{\ensuremath{\mathsf{MC}}\xspace}
\newcommand{\hyperplane}{h}

\newcommand{\QuickSort}{\Algorithm{QuickSort}\xspace}
\newcommand{\Alg}{\Algorithm{Alg}\xspace}

\newcommand{\Graph}{\mathsf{G}}%
\newcommand{\GraphB}{\mathsf{H}}
\newcommand{\GraphC}{\mathsf{K}}
\newcommand{\GraphD}{\mathsf{F}}

\providecommand{\G}{}
\renewcommand{\G}{\Graph}

\newcommand{\RLP}{\ComplexityClass{RLP}\xspace}
\newcommand{\USTCON}{\ComplexityClass{USTCON}}
\newcommand{\STCON}{\overrightarrow{\ComplexityClass{STCON}}}

\newcommand{\Mtrx}{\mathsf{M}}
\newcommand{\TMtrx}{\mathsf{Q}}
\newcommand{\IMtrx}{\mathcal{I}}

\newcommand{\Expander}[3]{$\left[#1,#2,#3\right]$-expander}

\newcommand{\ChartX}[1]{\ensuremath{\chi_{#1}}}
\newcommand{\Field}{\mathbb{F}}
\newcommand{\LSpace}{\mathbb{G}}
\newcommand{\LD}[2]{\mathsf{L{D}}\pth{#1,#2}}
\newcommand{\SEVN}{\widehat{\lambda_2}}
\newcommand{\NX}[2]{#1\!\left[#2\right]}
\newcommand{\Entry}[2]{#1\!\left[#2\right]}
\newcommand{\NXExt}[3]{#1_{#2}\!\left[#3\right]}
\newcommand{\RangeI}[1]{\llbracket #1 \rrbracket}
\newcommand{\DialX}[1]{\gamma\pth{#1}}
\newcommand{\BiDialX}[1]{{\gamma}_2\pth{#1}}

\newcommand{\EV}[1]{\widehat{\lambda_{#1}}}
\newcommand{\TEV}[1]{\widehat{\lambda_{#1}}}
\newcommand{\EVal}{\widehat{\lambda}}
\newcommand{\Basis}{\mathcal{B}}

\newcommand{\ZProd}[2]{\ensuremath{#1 \, \text{\textcircled{z}} #2}}
\newcommand{\RProd}[2]{\ensuremath{#1 \, \text{\textcircled{r}} #2}}

\newcommand{\OneVec}[1]{\mathbf{1}^{#1}}
\newcommand{\DstrbV}[1]{\mathbf{\mathsf{q}}^{(#1)}}
\newcommand{\DstrbE}[2]{\mathsf{q}^{(#1)}_{#2}}

\newcommand{\DEdges}{\widetilde{\Edges}}

\newcommand{\VV}{\mathsf{V}}%
\newcommand{\EE}{\mathsf{E}}%

\newcommand{\Vertices}{\VV}%
\newcommand{\Edges}{\EE}%

\newcommand{\EdgesX}[1]{\Edges\pth{#1}}
\newcommand{\VerticesX}[1]{V\pth{#1}}

\newcommand{\StringX}[1]{\mathsf{s}\pth{#1}}

\newcommand{\BMtrx}{\mathsf{B}}
\newcommand{\WMtrx}{\mathsf{W}}
\newcommand{\NWMtrx}{\overline{\mathsf{W}}}
\newcommand{\Sequence}{\EuScript{S}}

\newcommand{\med}{\mathrm{med}}
\newcommand{\medX}[1]{\mathrm{med}\pth{#1}}

\newcommand{\medianC}{\nu}%

\newlength{\savedparindent}
\newcommand{\SaveIndent}{\setlength{\savedparindent}{\parindent}}
\newcommand{\RestoreIndent}{\setlength{\parindent}{\savedparindent}}

\newcommand{\NNN}{\mathbb{N}}%
\providecommand{\NN}{\mathbb{N}}%
\newcommand{\reals}{\Re}%
\newcommand{\ZZ}{\mathbb{Z}}%

\newcommand{\level}{\mathop{\mathrm{level}}}

\providecommand{\TPDF}[2]{\texorpdfstring{#1}{#2}}

\newcommand{\Hint}[1]{\emph{[#1]}}
\newcommand{\Arr}{\mathop{\mathrm{\EuScript{A}}}}

\newcommand{\Array}{\mathcal{T}}%

\newcommand{\dirEdge}[2]{\left({{#1} \rightarrow {#2}} \right)}

%%%%%%%%%%%%%%%%%%%%%%%%%%%%%%%%%%%%%%%%%%%%%%%%%%%%%%%%%%%%%%%%%%%%%%
% Emphasize stuff...
\definecolor{blue25}{rgb}{0, 0, 11}
\newcommand{\emphic}[2]{%
   \textcolor{blue25}{%
      \bfseries{\emph{#1}}}%
   \index{#2}}
\newcommand{\emphi}[1]{\emphic{#1}{#1}}
\newcommand{\emphOnly}[1]{\textcolor{blue25}{\bfseries{\emph{#1}}}}
\providecommand{\emphiC}[1]{\emphic{#1}{#1}}

\definecolor{almostblack}{rgb}{0, 0, 0.3}
\providecommand{\emphw}[1]{{\textcolor{almostblack}{\emph{#1}}}}%

%
%%%%%%%%%%%%%%%%%%%%%%%%%%%%%%%%%%%%%%%%%%%%%%%%%%%%%%%%%%%%%%%%%%%%%%

\newcommand{\volX}[1]{\mathrm{vol}\pth{#1}}

\newcommand{\Interval}{\EuScript{I}}
\newcommand{\IntervalA}{\EuScript{I}}

\newcommand{\Exp}{\mathrm{Exp}}

\let\Ball\undefined
\newcommand{\Ball}{\mathbf{b}}
\newcommand{\hdim}{d}
\newcommand{\Domain}{\mathcal{D}}
\newcommand{\points}[1]{[\bfseries{#1 points}]}
\newcommand{\Event}{\EuScript{E}}%
\newcommand{\EventB}{\EuScript{B}}%
\newcommand{\EventF}{\EuScript{F}}%
\newcommand{\EventG}{\EuScript{G}}%

\newcommand{\distCmd}[1]{\left\| {#1}  \right\|}
\newcommand{\distY}[2]{\distCmd{#1 - #2}}

\newcommand{\NormalC}{\mathbf{N}}
\newcommand{\NormalY}[2]{\NormalC\pth{#1, #2}}%
\newcommand{\NormalDistOneDim}{\NormalC}
\newcommand{\NormalDistX}[1]{\NormalC^{#1}}
\DeclareMathOperator{\lca}{lca}
\newcommand{\RangeSpace}{\mathsf{S}}%
\newcommand{\RangeSpaceA}{\mathsf{T}}

\newcommand{\range}{\mathbf{r}}
\newcommand{\rangeA}{\mathbf{r}'}

\newcommand{\BadProb}{\varphi}

\newcommand{\EntropyChar}{\mathbb{H}}
\newcommand{\Entropy}[1]{\EntropyChar\pth{#1}}
\newcommand{\EntropyX}[1]{\EntropyChar\pth[]{#1}}

\newcommand{\Matrix}{\mathsf{M}}
\newcommand{\MatrixB}{\mathsf{B}}
\newcommand{\MatrixC}{\mathsf{C}}
\newcommand{\MatrixD}{\mathsf{D}}

\renewcommand{\th}{th\xspace}

\newcommand{\Family}{\mathcal{F}}
\newcommand{\FamilyB}{\mathcal{G}}
\newcommand{\FamilyC}{\mathcal{H}}
\newcommand{\FMSC}{\EuScript{Y}\index{metric space}}

\newcommand{\median}{\overline{\mathsf{m}}}
\newcommand{\QuickSelect}{\Algorithm{Quick{}Select}\xspace}

\providecommand{\row}{\mathbf{r}}%
\newcommand{\rowc}{{r}}%

\newcommand{\vecb}{\mathbf{b}}

\newcommand{\vecA}{\mathbf{v}}
\newcommand{\vecB}{\mathbf{u}}
\newcommand{\vecRand}{\mathbf{r}}

\newcommand{\vecH}{\mathbf{h}}
\newcommand{\vecV}{\mathbf{v}}

\newcommand{\vecX}{\mathbf{x}}
\newcommand{\vecY}{\mathbf{y}}

\newcommand{\DistBinom}[2]{\mathrm{Bin}\pth{#1,#2}}

\newcommand{\DistGeom}[1]{\mathrm{Geom}\pth{#1}}

\newcommand{\Sample}{\mathsf{R}}

\newcommand{\lcm}{\mathop{\mathrm{lcm}}}
\newcommand{\bdiv}{\mathop{\mathrm{div}}}%
\newcommand{\divides}{\mid}%
\newcommand{\poly}{\mathrm{poly}}%

\newcommand{\Group}{\mathcal{G}} 
\newcommand{\GroupA}{\mathcal{H}} 

\newcommand{\algGCD}{\Algorithm{EuclidGCD}\xspace}%

\newcommand{\gOp}{\times}%
\newcommand{\gIdent}{\mathsf{i}}%

\newcommand{\ordX}[2][\!]{\mathop{\mathrm{o{r}d}}\pth[#1]{#2}}

\newcommand{\lSym}[2]{%
   \left(
       \ds {#1}\divides{#2}%
   \right) %
}
\newcommand{\jSym}[2]{\left\llbracket%
       \ds {#1}\divides{#2}%
       \right\rrbracket%
}
\newcommand{\Line}{\ell}

\newcommand{\LevelChar}{\mathbb{L}}
\newcommand{\AtMostKLevelChar}{\LevelChar_{\leq k}}
\newcommand{\AtMostKLevel}[1]{\LevelChar_{\leq k} \pth{ #1}}

\newcommand{\LineSet}{\mathsf{L}}%
\newcommand{\ArrX}[1]{\EuScript{A}\pth{#1}}%

\newcommand{\ProblemC}[1]{{\textsf{\textcolor{red}{#1}\index{problem!#1}}}}
\DefineNamedColor{named}{OliveGreen} {cmyk}{0.64,0,0.95,0.40}
\providecommand{\ComplexityClass}[1]{{{\textcolor[named]{OliveGreen}{%
            \textsc{#1}}}}}
\providecommand{\NPHard}{{\ComplexityClass{NP-Hard}}%
   \index{NP!hard}\xspace}
\providecommand{\NPComplete}{\ComplexityClass{NP-Complete}%
   \index{NP!complete}\xspace}
\providecommand{\APXHard}{\ComplexityClass{APX-Hard}%
   \index{APX!Hard}\xspace}
\providecommand{\NPC}{\ComplexityClass{NPC}%
   \index{NP!complete}\xspace}
\providecommand{\NP}{\ComplexityClass{NP}%
   \index{NP}\xspace}
\providecommand{\POLYT}{\ComplexityClass{P}\xspace}

\newcommand{\envFrame}[1]{\begin{center}\fbox{\begin{minipage}[c]{0.75\textwidth}{#1}\end{minipage}}\end{center}}

\newcommand{\mysqrt}[1]{\ts\!\!\sqrt{#1}}
\newcommand{\SegSet}{\mathsf{S}}
\newcommand{\ArrVD}[1]{\Arr^{|} \pth{#1}}
\newcommand{\RIC}{\textsf{RIC}\index{random incremental
      construction}\xspace}
\newcommand{\ArrB}{\EuScript{B}}

\newcommand{\TreeOpt}{\EuScript{T}_\mathrm{opt}}
\newcommand{\TreeOptA}{\EuScript{T}_\mathrm{opt}'}

\newcommand{\DepthExt}[2]{\mathrm{depth}_{#1}\pth{#2}}
\newcommand{\depthX}[1]{\mathrm{depth}\pth{#1}}

\newcommand{\TreeA}{\EuScript{U}}
\newcommand{\TreeB}{\EuScript{V}}

\newcommand{\Trap}{\sigma}
\newcommand{\TrapA}{\sigma'}
\newcommand{\TrapB}{\sigma''}

\newcommand{\SampleA}{\mathsf{R}'}

\newcommand{\AllRegions}{\mathcal{T}}
\newcommand{\RegionsSet}[1]{\mathcal{F}\pth{#1}}
\newcommand{\RegionsHeavy}[2]{\mathcal{F}_{\!\geq#1}\pth{#2}}
\newcommand{\NumOfRegionsZero}[1]{\ExChar \!f\pth{#1}}
\newcommand{\NumOfRegions}[2]{\ExChar\! f_{\!\geq #1}\pth[]{#2}}
\newcommand{\ObjSet}{\mathsf{S}}

\newcommand{\WeightX}[2][\!]{\omega \pth[#1]{#2}}
\newcommand{\Edge}{\mathsf{e}}

\newcommand{\Face}{\mathsf{f}}

\newcommand{\DefSet}[1]{D\pth{#1}}
\newcommand{\KillSet}[1]{K\pth{#1}}

\newcommand{\CSProbFull}[4]{
   \,\rho_{\!\!\stackrel{\rule[-.005cm]{0cm}{0.1cm}}{#1,
       #4}}\pth{\MakeSBig #2,#3}}

\newcommand{\disk}{D}
\newcommand{\diskA}{D'}

\newcommand{\IntervalSet}{\widehat{\mathcal{I}}}

\newcommand{\GroundSet}{\textsf{X}}
\newcommand{\FGroundSet}{\mathsf{x}}
\newcommand{\RangeSet}{{\mathcal{R}}}
\newcommand{\RangeSetA}{\mathcal{R}'}
\newcommand{\MeasureChar}{\overline{m}}
\newcommand{\Measure}[1]{\MeasureChar\pth{#1}}
\newcommand{\sMeasure}[1]{\overline{s}\pth{#1}}

\newcommand{\SetA}{S}%
\newcommand{\SetB}{T}%
\newcommand{\SetC}{N}%

\newcommand{\VCProj}[2]{#1_{|#2}}
\newcommand{\VCDim}{\mathrm{dim}_{\VC}}
\newcommand{\Dim}{\delta}
\newcommand{\DimA}{\delta'}
\newcommand{\PntSetB}{{{\mathsf{Q}}}}
\newcommand{\PntSetC}{{{\mathsf{R}}}}
\newcommand{\Growth}[2]{{\EuScript{G}}_{#1}\pth{#2}}
\newcommand{\ShatterFunc}[1]{\pi_{#1}}
\newcommand{\Distribution}{\EuScript{D}}
\newcommand{\Error}{\EuScript{E}}

\newcommand{\bd}{{\partial}}%
\newcommand{\eps}{{\varepsilon}}%

\newcommand{\clmlab}[1]{\label{claim:#1}}
\newcommand{\clmref}[1]{\HLink{Claim}{claim:#1}}

\newcommand{\CHX}[1]{{\mathcal{CH}}\pth{#1}}

\newcommand{\ts}{\hspace{0.6pt}}

\def\qedsymbol{\myqedsymbol}%

\newcommand{\qedsign}{\hfill{\hfill\myqedsymbol}}
\newcommand{\qedsignBig}{\hfill{\hfill\rule{4mm}{4mm}}}
\newcommand{\QED}{{\hfill{\hfill\qedsymbol}}}

\newcommand{\num}{\alpha}

\newcommand{\NEW}[1]{#1}

\newcommand{\IncludeGraphics}[2][]{%
   \typeout{}%
   \typeout{Graphics: #2}%
   \typeout{\ includegraphics[#1]{#2}}%
   \includegraphics[#1]{#2}
   \typeout{}%
}

%    -----------------------------------------------------------------
%    -----------------------------------------------------------------
%    -----------------------------------------------------------------
%    Handling references
%    -----------------------------------------------------------------
%    -----------------------------------------------------------------
%    -----------------------------------------------------------------
\newcommand{\HLinkShort}[2]{\hyperref[#2]{#1\ref*{#2}}}
\newcommand{\HLink}[2]{\hyperref[#2]{#1~\ref*{#2}}}
\newcommand{\HLinkPage}[2]{\hyperref[#2]{#1~\ref*{#2}%
      $_\text{p\pageref{#2}}$}}

\newcommand{\HLinkSuffix}[3]{\hyperref[#2]{#1\ref*{#2}{#3}}}
\newcommand{\HLinkPageSuffix}[3]{\hyperref[#2]{#1\ref*{#2}%
      #3$_\text{p\pageref{#2}}$}}

\newcommand{\figlab}[1]{\label{fig:#1}}
\newcommand{\figref}[1]{\HLink{Figure}{fig:#1}}

\newcommand{\exelab}[1]{\label{exer:#1}}

\providecommand{\deflab}[1]{\label{def:#1}}
\newcommand{\defref}[1]{\HLink{Definition}{def:#1}}
\newcommand{\defrefY}[2]{\hyperref[def:#2]{#1}}
\newcommand{\defrefpage}[1]{\HLinkPage{Definition}{def:#1}}

\newcommand{\etal}{\textit{et~al.}\xspace}

\newcommand{\obslab}[1]{\label{observation:#1}}
\newcommand{\obsref}[1]{\HLink{Observation}{observation:#1}}
\newcommand{\obsrefY}[2]{\hyperref[observation:#2]{#1}}

\newcommand{\exmlab}[1]{\label{example:#1}}
\newcommand{\exmref}[1]{\HLink{Example}{example:#1}}

\newcommand{\remlab}[1]{\label{rem:#1}}
\newcommand{\remref}[1]{\HLink{Remark}{rem:#1}}%

\newcommand{\chaplab}[1]{\label{chap:#1}}

\newcommand{\itemlab}[1]{\label{item:#1}}
\newcommand{\itemref}[1]{\HLinkSuffix{}{item:#1}{}}

\newcommand{\tedlab}[1]{\label{tedium:#1}}
\newcommand{\tedref}[1]{\HLink{Tedium}{tedium:#1}}%

% Deferred tedium: `by \TediumHere{...}, we have ...` inside a proof.
% Use \TediumHereY{...}{label} to allow later citations with \tedref{label}.
% Bodies may contain paragraphs and displays. They are typeset, in call order,
% after the enclosing proof, using ordinary tedium numbering and references.
% As with other forward references, compile twice to resolve the links.
% Define macros needed by a body outside the proof: its local scope has ended
% when the deferred body is typeset. Verbatim material is not supported.
\ExplSyntaxOn
\int_new:N \g__notes_tedium_proof_int
\int_new:N \g__notes_tedium_item_int
\seq_new:N \g__notes_tedium_stack_seq
\tl_new:N \l__notes_tedium_proof_tl

\cs_new_protected:Npn \__notes_tedium_begin:
  {
    \int_gincr:N \g__notes_tedium_proof_int
    \seq_new:c
      { g__notes_tedium_\int_use:N \g__notes_tedium_proof_int _seq }
    \seq_gput_right:Nx \g__notes_tedium_stack_seq
      { \int_use:N \g__notes_tedium_proof_int }
  }

\cs_new_protected:Npn \__notes_tedium_print:n #1
  {
    \seq_use:cn { g__notes_tedium_#1_seq } { }
    \cs_undefine:c { g__notes_tedium_#1_seq }
  }

\cs_new_protected:Npn \__notes_tedium_end:
  {
    \seq_gpop_right:NN \g__notes_tedium_stack_seq
      \l__notes_tedium_proof_tl
    \exp_args:NV \__notes_tedium_print:n \l__notes_tedium_proof_tl
  }

\cs_new_protected:Npn \__notes_tedium_queue:nn #1#2
  {
    \seq_if_empty:NTF \g__notes_tedium_stack_seq
      {
        \PackageError{notes}
          {Deferred~tedium~used~outside~a~proof}
          {Use~\string\TediumHere\space or~\string\TediumHereY
           \space inside~a~proof~environment.}
      }
      {
        \seq_get_right:NN \g__notes_tedium_stack_seq
          \l__notes_tedium_proof_tl
        \seq_gput_right:cx
          { g__notes_tedium_\l__notes_tedium_proof_tl _seq }
          {
            \exp_not:N \begin {tedium}
            \exp_not:N \tedlab
              { \exp_not:n {#2} }
            \exp_not:n {#1}
            \exp_not:N \end {tedium}
          }
        \tedref{#2}
      }
  }

\NewDocumentCommand{\TediumHere}{+m}
  {
    \int_gincr:N \g__notes_tedium_item_int
    \exp_args:Nnx \__notes_tedium_queue:nn {#1}
      { auto-td-\int_use:N \g__notes_tedium_item_int }
  }

\NewDocumentCommand{\TediumHereY}{+m m}
  { \__notes_tedium_queue:nn {#1}{#2} }

% Each proof has its own queue; the stack also keeps nested proofs separate.
\clist_map_inline:nn
  { proof,altproof,proofextin,proofnoqed,proofof,proofExt,proofnoproof }
  {
    \AddToHook{env/#1/before}[notes/tedium]{\__notes_tedium_begin:}
    \AddToHook{env/#1/after}[notes/tedium]{\__notes_tedium_end:}
  }
\ExplSyntaxOff

\newcommand{\factlab}[1]{\label{fact:#1}}%

\newcommand{\alglab}[1]{\label{algorithm:#1}}

\newcommand{\net}{\mathcal{N}}%

\newcommand{\lemlab}[1]{\label{lemma:#1}}
\newcommand{\lemref}[1]{\HLink{Lemma}{lemma:#1}}%
\newcommand{\lemrefY}[2]{\hyperref[lemma:#2]{#1}}
\newcommand{\lemrefpage}[1]{\HLinkPage{Lemma}{lemma:#1}}
\newcommand{\lemrefshort}[1]{\HLinkShort{L}{lemma:#1}}

\newcommand{\problab}[1]{\label{prob:#1}}%
\newcommand{\probref}[1]{\HLink{Problem}{prob:#1}}%

\newcommand{\eqlab}[1]{\label{equation:#1}}
\newcommand{\Eqrefbare}[1]{\HLinkSuffix{(}{equation:#1}{)}}
\newcommand{\Eqref}[1]{\HLinkSuffix{Eq.~(}{equation:#1}{)}}
\newcommand{\Eqrefpage}[1]{\HLinkPageSuffix{Eq.~(}{equation:#1}{)}}

\newcommand{\seclab}[1]{\label{sec:#1}}
\newcommand{\secref}[1]{\HLink{Section}{sec:#1}}
\newcommand{\secrefY}[2]{\hyperref[sec:#2]{#1}}
\newcommand{\secrefpage}[1]{\HLinkPage{Section}{sec:#1}}

\newcommand{\thmlab}[1]{{\label{theo:#1}}}
\newcommand{\thmref}[1]{\HLink{Theorem}{theo:#1}}
\newcommand{\thmrefY}[2]{\hyperref[theo:#2]{#1}}
\newcommand{\thmrefpage}[1]{\HLinkPage{Theorem}{theo:#1}}

\newcommand{\corlab}[1]{\label{cor:#1}}
\newcommand{\corref}[1]{\HLink{Corollary}{cor:#1}}%
\newcommand{\correfY}[2]{\hyperref[cor:#2]{#1}}

\newcommand{\Yend}{Y_{\mathrm{end}}}%

%    -----------------------------------------------------------------
%    -----------------------------------------------------------------
%    -----------------------------------------------------------------
%    Handling references
%    -----------------------------------------------------------------
%    -----------------------------------------------------------------
%    -----------------------------------------------------------------

\newcommand{\ceil}[1]{\left\lceil {#1} \right\rceil}
\newcommand{\floor}[1]{\left\lfloor {#1} \right\rfloor}
\newcommand{\brc}[1]{\left\{ {#1} \right\}}
\newcommand{\pth}[2][\!]{\mleft({#2}\mright)}%
\newcommand{\brk}[1]{\mleft[{#1}\mright]}%
\newcommand{\pbrc}[2][\!\!]{#1\left[ {#2} \bigr. \right]}
\newcommand{\pbrcS}[1]{\left[ {#1} \right]}

\newcommand{\dotprod}[2]{{#1}\cdot{#2}}
\newcommand{\DotProd}[2]{\permut{{#1},{#2}}}
\newcommand{\permut}[1]{\left\langle {#1} \right\rangle}

\newcommand{\MakeBig}{}
\newcommand{\MakesBig}{}
\newcommand{\MakeSBig}{}
\newcommand{\MakeVBig}{}

\newcommand{\DistChar}{\mathsf{d}}%
\newcommand{\DistSetY}[2]{\DistChar\pth{#1,#2}}
\newcommand{\distG}[2]{\DistChar_\Graph\pth{#1, #2}}
\newcommand{\DistMtrCharX}[1]{\DistChar_{#1}}

\newcommand{\dY}[2]{\DistChar\pth{#1, #2}}%
\newcommand{\dGY}[2]{\DistChar_\Graph\pth{#1, #2}}%

\newcommand{\cset}{\mathcal{C}}%

\newcommand{\MTR}{\mathcal{M}}
\newcommand{\DistMtr}[2]{\DistChar\pth{#1,#2}}

\newcommand{\dist}[1]{\left\| {#1} \right\|}
\newcommand{\norm}[1]{\left\| {#1} \right\|}

\providecommand{\ds}{\displaystyle}
\newcommand{\sep}[1]{\,\left|\, {#1} \bigr.\right.}%
\newcommand{\sepw}[1]{\left|\, {#1} \right.}%

\newcommand{\ProbExt}[2]{\mathop{\mathbf{Pr}}_{#1}\pbrcx{#2}}
\newcommand{\ProbLTR}{\mathbb{P}}%

\newcommand{\Prob}[1]{\mathop{\ProbLTR} \mleft[ #1 \mright]}%
\newcommand{\ProbCond}[2]{\mathop{\ProbLTR}\!\left[%
       #1 \;\middle\vert\; #2 \right]}
\newcommand{\ProbY}[2]{\ProbLTR^{}_{\!#1}\pbrc{#2}}

\newcommand{\ExChar}{\mathbb{E}}%
\newcommand{\ExSym}{\mathop{\ExChar}}%
\newcommand{\Ex}[2][\!]{\ExSym#1\pbrcx{#2}}
\newcommand{\ExExt}[2]{\ExSym_{#1}\!\pbrcx{#2}}
\newcommand{\ExExtS}[2]{\ExSym_{\substack{#1}}\!\pbrcx{#2}}
\newcommand{\ExCond}[2]{\ExSym\!\left[%
       #1 \;\middle\vert\; #2 \right]}

\newcommand{\VarSymb}{\mathop{\mathbb{V}}}
\newcommand{\Var}[1]{\mathop{\mathbb{V}}\!\pbrcx{#1}}

\newcommand{\pbrcx}[1]{\left[ {#1} \right]}%
\newcommand{\nfrac}[2]{{#1}/{#2}}
\newcommand{\cardin}[1]{\left| {#1} \right|}%

\newcommand\CCancel[2][black]{{\renewcommand\CancelColor{\color{#1}}\cancel{#2}}}

\newcommand{\IncGraphPage}[4][]{%
   \IncludeGraphics[page=#4,#1]{\File{#2/#3}}%
}

\newenvironment{enumEx}{
   \begin{compactenumA}
    }{\end{compactenumA}}

\newcommand{\cl}{{\mathrm{cl}}}

\newcommand{\HSetY}[2]{H\pth{#1, #2}}%
\newcommand{\CHSetY}[2]{C\pth{#1, #2}}%
\newcommand{\OldOm}{{\mathcal{O}}}
\newcommand{\NewOm}{{\mathcal{N}}}

%%%%%%%%%%%%%%%%%%%%%%%%%%%%%%%%%%%%%%%%%%%%%%%%%%%%%%%%%%%%%%%%%%%%%%%
%---------------------------------------------------------------
\newcommand{\IntRange}[1]{\left\llbracket #1 \right\rrbracket}
\newcommand{\IRX}[1]{\left\llbracket #1 \right\rrbracket}%
\newcommand{\IRY}[2]{\left\llbracket #1:#2 \right\rrbracket}

%---------------------------------------------------------------
%%%%%%%%%%%%%%%%%%%%%%%%%%%%%%%%%%%%%%%%%%%%%%%%%%%%%%%%%%%%%%%%%%%%%%%

\newcommand{\obj}{f}%
\newcommand{\Wg}{\omega}%
\newcommand{\TotalW}{W}%

\newcommand{\rankY}[2]{{#1}_{ \permut{#2}}}%%
\newcommand{\areaX}[1]{\mathrm{area}\pth{#1}}%%

\newcommand{\mybinom}[2]%
{\Bigl(\begin{array}{@{}c@{}}#1\\#2\end{array}\Bigr)}

\providecommand{\Set}[2]{\left\{ #1 \;\middle\vert\; #2 \right\}}
\newcommand*{\biga}[1]{\vcenter{\hbox{\scalebox{1.5}{\ensuremath#1}}}}
\newcommand{\CNF}{\TermI{CNF}\xspace}%
\newcommand{\Tree}{T}%

\DeclareSymbolFont{sfletters}{OML}{cmbrm}{m}{it}

\DeclareMathSymbol{\salpha}{\mathord}{sfletters}{"0B}

%%%%%%%%%%%%%%%%%%%%%%%%%%%%%%%%%%%%%%%%%%%%%%%%%%%%%%%%%%%%%%%%%%%%

\providecommand{\ProblemInd}[2]{\textsf{\textcolor{red}{#1}\index{Problem!#2}}}
\newcommand{\NPProblemInd}[4]{
   \begin{center}
      \vspace{0.6cm}
       \fbox{
          \begin{minipage}{0.9\linewidth}
              \vspace{-0.6cm} \hspace{-0.35cm}
              {\large \bfseries \ProblemInd{#1}{#2}}
              \vspace{+0.2cm}%

              \noindent{\textcolor{blue}{\bfseries \textsf{Instance}}:} {#3}
              
              \noindent{\textcolor{blue}{\bfseries \textsf{Question}:}} {#4}
          \end{minipage}
       }
   \end{center}
   \smallskip
   }

\newcommand{\NPProblemOpt}[4]{\NPProblemInd{#1}{#2}{#3}{#4}}

%%%%%%%%%%%%%%%%%%%%%%%%%%%%%%%%%%%%%%%%%%%%%%%%%%%%%%%%%%%%%%%%%%%%

\newcommand{\Hastad}{H{\aa}stad\xspace}

\newcommand{\ANNbr}{\Term{ANN}\index{approximate nearest neighbor}\xspace}
\newcommand{\LSH}{\TermI{L{S}H}\xspace}

\newcommand{\DS}{\EuScript{D}} % DataStructure
\newcommand{\DSC}{\mathsf{D}} % DataStructure

\newcommand{\distH}[2]{\DistChar_{H}\pth{#1, #2}}
\newcommand{\distHY}[2]{\DistChar_{H}\pth{#1, #2}}

\newcommand{\DSAprxNearNbr}{\ensuremath{\DataStructure{{\approx%
         \text{Near{}}}}}%
   \index{near neighbor!data-structure!approximate}%
}
\newcommand{\DataStructure}[1]{\ensuremath{\DS_{#1}}}
\newcommand{\DistNN}[2]{\DistChar\pth{#1,#2}}
\newcommand{\pntc}{{p}}%

\newcommand{\seq}{{M}}%
\newcommand{\cTimes}{\beta}%

\newcommand{\DSTimes}{L}%
\newcommand{\rad}{r}%

\newcommand{\HypercubeX}[1]{\Hypercube^{#1}}
\newcommand{\kFunctions}[2]{\mathrm{combine}\pth{#1,#2}}
\newcommand{\WEqualChar}{\Bumpeq}
\newcommand{\WEqual}[2]{{#1}\WEqualChar{#2}}
\newcommand{\FamilyAmp}{\FamilyC_{\WEqualChar}}
\newcommand{\PSPACE}{\ComplexityClass{PSPACE}\xspace}%

\newcommand{\PS}{P}%

\newcommand{\PSetFar}{\ensuremath{\PntSet^{}_{\geq}\xspace}}%
\newcommand{\DA}{D}%
\newcommand{\uniqX}[1]{\mathrm{uniq}\pth{#1}}%
\newcommand{\diffX}[1]{\mathop{}\!\mathrm{d}#1}

\newcommand{\Frq}{F}%
\newcommand{\frq}{f}%
\newcommand{\STRM}{\mathcal{S}}%
\newcommand{\strm}{{s}}%

\newcommand{\Calinescu}{C\u{a}linescu\xspace}

\newcommand{\fGrowX}[1]{G^{}_{\Dim}\pth{#1}}
\newcommand{\CompressFunc}{\mathtt{Compress}}
\newcommand{\RSA}{\Term{RSA}\xspace}
\newcommand{\MST}{\Term{MST}\xspace}
\providecommand{\CClassPP}{\CClass{P{}P}\xspace}
\providecommand{\ID}{\mathrm{id}}

\newcommand{\EncodeFunc}{\text{\ttfamily{Enc}}}
\newcommand{\DecodeFunc}{\text{\ttfamily{Dec}}}
\newcommand{\discX}[1]{\mathrm{disc}\pth{#1}}%

\newcommand{\propertyC}{\psi}%
\newcommand{\propertyX}[1]{\propertyC\pth{#1}}%
\newcommand{\numC}{\Gamma}%

\newcommand{\degreeX}[1]{\mathrm{degree}\pth{#1}}

\newcommand{\ind}{m}

\providecommand{\nequiv}{\not\equiv}
\providecommand{\sqcupplus}{\sqcup^+}

\newcommand{\AlgorithmI}[1]{{\protect\color{RedViolet}{#1}}}

\newcommand{\binsY}[2]{\sqcup_{\geq #1}\pth{#2}}

\newcommand{\BLX}[1]{{{\sqcupplus{}}_{\geq #1}}}

\newcommand{\pdf}{\textsf{pdf}\xspace}

\newcommand{\cdf}{\textsf{cdf}\xspace}

\newcommand{\costX}[1]{\mathrm{cost}\pth{#1}}

\newcommand{\codeX}[1]{\mathrm{code}\pth{#1}}

\newcommand{\lengthX}[1]{\mathrm{l{}e{}n}\pth{#1}}
}{\input{prefix_common}}%

\newcommand{\HST}{\Term{HST}\xspace}

\newcommand{\FMS}{\EuScript{X}}  

\newcommand{\FMSB}{\EuScript{Y}}  

\renewcommand{\Re}{\mathbb{R}}%

\newcommand{\ballsY}[2]{\mathcal{F}_{\geq #1}\pth{#2}}

\DeclareRobustCommand{\Matousek}{Matoušek\xspace}

\DeclareRobustCommand{\Erdos}{Erdős\xspace}

\DeclareRobustCommand{\Lovasz}{Lovász\xspace}

\StartDefiningTabulars
\newcommand{\ThreeFiguresExt}[9]{%
   \begin{tabular}{ccc}
     \begin{minipage}{0.31\linewidth}
         \begin{center}
             {\IncludeGraphics[#1]{#2}}
         \end{center}
     \end{minipage}
     &
       \begin{minipage}{0.31\linewidth}
           \centerline{\IncludeGraphics[#4]{#5}}
       \end{minipage}
     &
       \begin{minipage}{0.31\linewidth}
           \centerline{\IncludeGraphics[#7]{#8}}
       \end{minipage}
     \\
     ~\\[-0.2cm]
     {#3} & {#6} & {#9}
   \end{tabular}}
\StopDefiningTabulars

\newcommand{\exeref}[1]{Exercise~\ref{exer:#1}}

% Theorem and proof rendering
\usepackage[amsmath,thmmarks]{ntheorem}
\theoremseparator{.}%

\newtheorem{theorem:unnumbered}{Theorem}[section]
\newtheorem{theorem}{Theorem}[section]

\newtheorem{observation}[theorem]{Observation}
\newtheorem{conjecture}[theorem]{Conjecture}
\newtheorem{lemma}[theorem]{Lemma}
\newtheorem{corollary}[theorem]{Corollary}
\newtheorem{claim}[theorem]{Claim}
\newtheorem{assumption}[theorem]{Assumption}
\newtheorem{tedium}[theorem]{Tedium}

\theoremstyle{plain}%
\theoremheaderfont{\sffamily}
\theorembodyfont{\upshape}%

\newtheorem{defn}[theorem]{Definition}
\newtheorem{definition}[theorem]{Definition}

\newtheorem*{remarkunnumbered}{Remark}%
\newenvironment{remark_unnumbered}{\begin{remarkunnumbered}}{\end{remarkunnumbered}}
\newtheorem*{remarks}{Remarks}%
\newtheorem{remark}[theorem]{Remark}%
\newtheorem{fact}[theorem]{Fact}%
\newtheorem{example}[theorem]{Example}
\newtheorem{exercise}[theorem]{Exercise}
\newtheorem{problem}[theorem]{Problem}
\newtheorem{algorithm}[theorem]{Algorithm}%
\newtheorem{xca}[theorem]{Exercise}
\newtheorem{openproblem}[theorem]{Open Problem}
\newtheorem{exerciseh}[theorem]{Exercise}
\newenvironment{exercise_h}[1][]{\begin{exerciseh}[#1]}{\end{exerciseh}}

% Proof environments
\theoremheaderfont{\em}%
\theorembodyfont{\upshape}%
\theoremstyle{nonumberplain}
\theoremseparator{}
\theoremsymbol{\myqedsymbol}
\newtheorem{proof}{Proof:}
\newtheorem{proofextin}{Proof}
\newtheorem{altproof}{Alternative proof:}

\theoremsymbol{}
\newtheorem{proofnoqed}{Proof:}

\theoremsymbol{\myqedsymbol}
\theorempreskip{0pt}%
\theorempostskip{0pt}%
\newtheorem{proofof}{Proof of\!}%

\newenvironment{proofExt}[1]{{\emph{#1:}}}{\hfill{\hfill\myqedsymbol}}
\newenvironment{proofnoproof}{{}}{{\hfill{\hfill\myqedsymbol}}}

%%%%%%%%%%%%%%%%%%%%%%%%%%%%%%%%%%%%%%%%%%%%%%%%%%%%%%%%%%%%%%%%%%%%
%%%%%%%%%%%%%%%%%%%%%%%%%%%%%%%%%%%%%%%%%%%%%%%%%%%%%%%%%%%%%%%%%%%%
% Macros to embed figures in text

\newcommand{\EBRY}[2]{#1_{\langle#2\rangle}}%

\newcommand{\NiceFrac}[2]{%
   {\raisebox{0.2em}{\ensuremath{#1}}\!\!\!}%
   {\rotatebox[origin=c]{-20}{\Big/}}% Rotated /
   {\! \raisebox{-0.2em}{\ensuremath{#2}} }}%

% end - Macros to embed figures in text
%%%%%%%%%%%%%%%%%%%%%%%%%%%%%%%%%%%%%%%%%%%%%%%%%%%%%%%%%%%%%%%%%%%% 
%%%%%%%%%%%%%%%%%%%%%%%%%%%%%%%%%%%%%%%%%%%%%%%%%%%%%%%%%%%%%%%%%%%%

%%%%%%%%%%%%%%%%%%%%%%%%%%%%%%%%%%%%%%%%%%%%%%%%%%%%%%%%%%%%%%%%%%
% Color in math
\providecommand{\Mh}[1]{{#1}}%
\providecommand{\IfColorMath}[2]{#2}%

\IfFileExists{.latex_printer_friendly}{\def\GenPrinterVer{1}}{}%

\ifx\GenPrinterVer\undefined
   \IfFileExists{.latex_color}{\def\GenColorMath{1}}{}
\else
   \renewcommand{\IfPrinterVer}[2]{#1}%
\fi

\ifx\GenColorMath\undefined
\else
\renewcommand{\IfColorMath}[2]{#1}%
\renewcommand{\Mh}[1]{{\textcolor{red}{#1}}}%
\fi
% Color in math - end
%%%%%%%%%%%%%%%%%%%%%%%%%%%%%%%%%%%%%%%%%%%%%%%%%%%%%%%%%%%%%%%%%%

\DeclareRobustCommand{\Holder}{Hölder\xspace}

\DeclareRobustCommand{\Turan}{Turán\xspace}

\usepackage{marginnote}
\usepackage{ragged2e} 
% Short, subtle snark macro
\setlength{\marginparwidth}{1.5in}
\setlength{\marginparsep}{10pt}
\newcommand{\snark}[1]{%
   \marginnote{\raggedright\scriptsize\sffamily\color{gray}#1}%
}

\MakePerPage{footnote} % Resets footnote counter on every page

\allowdisplaybreaks % Allows page breaks inside multiline math environments

%-----------------------------------------------------------------------------
% Bibliography & Document Begin
%-----------------------------------------------------------------------------
\addbibresource{shared/book.bib}%
\makeindex
\begin{document}
}{\input{prefix_lwarp}}%
\else
  \ifdefined\HCode
    \IfFileExists{shared/prefix_4ht.tex}{%-----------------------------------------------------------------------------
% TeX4ht Prefix File for Randomized Algorithms Class Notes
% Kept separate from prefix_latex.tex so HTML-specific simplifications cannot
% affect the standard PDF or lwarp builds.
%-----------------------------------------------------------------------------

\RequirePackage{iftex}
\ifPDFTeX
  \PackageError{engine-guard}{This document requires LuaLaTeX or XeLaTeX}{%
    Please recompile using LuaLaTeX or XeLaTeX instead of pdfLaTeX.%
  }
  \stop
\fi

\ifdefined\HCode\else
  \PackageError{tex4ht-gate}{%
    Attempting to use prefix_4ht.tex outside TeX4ht.%
  }{Compile through make4ht, or use prefix_latex.tex for PDF generation.}
  \stop
\fi

% Always use memoir
\documentclass[12pt,oneside,fleqn]{memoir}

%-----------------------------------------------------------------------------
% Typography, Fonts & Math Packages
%-----------------------------------------------------------------------------
\usepackage{mathtools}
\PassOptionsToPackage{warnings-off={mathtools-overbracket,mathtools-colon}}{unicode-math}
\usepackage{fontspec}
\usepackage{unicode-math}

% New Computer Modern OpenType text and math fonts
\setmainfont{NewCM10-Book.otf}[
  BoldFont          = NewCM10-Bold.otf,
  ItalicFont        = NewCM10-BookItalic.otf,
  BoldItalicFont    = NewCM10-BoldItalic.otf,
  SlantedFont       = NewCM10-Book.otf,
  SlantedFeatures   = {FakeSlant = 0.25},
  BoldSlantedFont   = NewCM10-Bold.otf,
  BoldSlantedFeatures = {FakeSlant = 0.25},
  SmallCapsFeatures = {Letters = SmallCaps}
]
\setsansfont{NewCMSans10-Book.otf}[
  BoldFont          = NewCMSans10-Bold.otf,
  ItalicFont        = NewCMSans10-BookOblique.otf,
  BoldItalicFont    = NewCMSans10-BoldOblique.otf
]
\setmonofont{NewCMMono10-Book.otf}[
  BoldFont          = NewCMMono10-Bold.otf,
  ItalicFont        = NewCMMono10-BookItalic.otf,
  BoldItalicFont    = NewCMMono10-BoldOblique.otf
]

\setmathfont{NewCMMath-Book.otf}
\setmathfont{NewCMMath-Bold.otf}[version=bold]

% OpenType math aliases and fixes
\providecommand{\square}{\mdlgwhtsquare}
\def\Box{\mdlgwhtsquare}
\def\Diamond{\mdlgwhtdiamond}
\def\blacksquare{\mdlgblksquare}
\def\leadsto{\rightsquigarrow}
\providecommand{\llbracket}{\lBrack}
\providecommand{\rrbracket}{\rBrack}
\DeclareTextCommand{\textcircled}{TU}[1]{%
  \mbox{\ooalign{\hfil\raise0.1ex\hbox{#1}\hfil\crcr\hbox{\ensuremath{\bigcirc}}}}%
}

\usepackage{alphabeta}
\usepackage[normalem]{ulem} % Provides \sout
\usepackage{mleftright}
\usepackage{microtype}
\usepackage{xspace}
\usepackage{xurl}
\usepackage[dvipsnames,table]{xcolor}
\usepackage{graphicx}
\usepackage{caption}
\usepackage{booktabs,array,colortbl,hhline,multirow,tabularx}
\usepackage[inline]{enumitem}
\usepackage[thicklines]{cancel}
\renewcommand\CancelColor{\color{purple}}
\usepackage{euscript}
\usepackage{marvosym}
\let\Cross\undefined
\let\Square\undefined
\usepackage{bbding}
\usepackage{epigraph}
\usepackage{makeidx}
\usepackage{perpage}
\usepackage{tikz}

%-----------------------------------------------------------------------------
% Boondox / mathcalb font declaration
%-----------------------------------------------------------------------------
\DeclareFontFamily{U}{BOONDOX-calo}{\skewchar\font=45 }
\DeclareFontShape{U}{BOONDOX-calo}{m}{n}{
  <-> s*[1.05] BOONDOX-r-calo}{}
\DeclareFontShape{U}{BOONDOX-calo}{b}{n}{
  <-> s*[1.05] BOONDOX-b-calo}{}
\DeclareMathAlphabet{\mathcalb}{U}{BOONDOX-calo}{m}{n}
\SetMathAlphabet{\mathcalb}{bold}{U}{BOONDOX-calo}{b}{n}
\DeclareMathAlphabet{\mathbcalb}{U}{BOONDOX-calo}{b}{n}

%-----------------------------------------------------------------------------
% Algorithm Environment (formerly salgorithm.sty)
%-----------------------------------------------------------------------------
\newcommand{\ProgWordX}[1]{{\small\textbf{#1}}{\ }}
\newcommand{\Do}{\ProgWordX{do}\xspace}
\newcommand{\End}{\ProgWordX{end}\xspace}
\newcommand{\Return}{\ProgWordX{return}\xspace}%
\newcommand{\Procbegin}{\ProgWordX{begin}}
\newcommand{\If}{\ProgWordX{if}\xspace}
\newcommand{\Then}{\ProgWordX{then}\xspace}
\newcommand{\Else}{\ProgWordX{else}\xspace}
\newcommand{\For}{\ProgWordX{for}\xspace}
\newcommand{\While}{{\small\bfseries while}\ }

\newcommand{\xbeginlgox}{\begin{minipage}{1in}\begin{tabbing}
           \quad\=\qquad\=\qquad\=\qquad\=\qquad\=\qquad\=\qquad\=\kill}
\newcommand{\xendlgox}{\end{tabbing}\end{minipage}}

\newenvironment{algorithmEnv}{%
   \begin{tabular}{|l|}\hline\xbeginlgox\bigskip}%
   {\xendlgox\\\hline\end{tabular}}

\newenvironment{program}{%
   \begin{minipage}{4.0in}%
   \begin{tabbing}%
       \ \ \ \ \=\ \ \ \=\ \ \ \ \=\ \ \ \ \=\ \ \ \ \=%
      \ \ \ \ \=\ \ \ \ \=\ \ \ \ \=\ \ \ \ \=%
      \ \ \ \ \=\ \ \ \ \=\ \ \ \ \=\ \ \ \ \=\kill%
}{%
   \end{tabbing}%
   \end{minipage}%
}

\newcommand{\AlgorithmBox}[1]{%
   \fbox{\begin{program} #1 \end{program}}%
}

\definecolor{AlgorithmColor}{cmyk}{0.07,0.90,0,0.34}
\providecommand{\Algorithm}[1]{{\AlgorithmI{#1}\index{algorithm!#1@{\AlgorithmI{#1}}}}}
\newcommand{\CodeComment}[1]{\textcolor{blue}{\texttt{#1}}}

%-----------------------------------------------------------------------------
% Path & Directory Resolution
%-----------------------------------------------------------------------------
\newcommand{\RootDIR}{../}
\IfFileExists{../../../book.tex}{\renewcommand{\RootDIR}{../../../}}{}
\IfFileExists{../../book.tex}{\renewcommand{\RootDIR}{../../}}{}
\IfFileExists{../book.tex}{\renewcommand{\RootDIR}{../}}{}
\IfFileExists{./book.tex}{\renewcommand{\RootDIR}{}}{}

\providecommand{\File}[1]{#1}
\providecommand{\RootFile}[1]{\RootDIR #1}

%-----------------------------------------------------------------------------
% QED Duck Symbol
%-----------------------------------------------------------------------------
\IfFileExists{shared/qed_duck.svg}{%
   \providecommand{\myqedsymbol}{%
      \HCode{<img class="qed-duck" src="shared/qed_duck.svg" alt="QED" />}%
   }%
}{%
   \providecommand{\myqedsymbol}{\HCode{<span class="qed-fallback">&#x25A0;</span>}}%
}

%-----------------------------------------------------------------------------
% BibLaTeX Setup (Pure Biber, No BibTeX)
%-----------------------------------------------------------------------------
\usepackage[bibencoding=utf8,style=alphabetic,refsegment=chapter,
indexing=cite,
backend=biber]{biblatex}

\ExecuteBibliographyOptions{%
   maxnames = 40,%
   maxalphanames = 4,%
   maxcitenames = 3,%
   minalphanames = 3,%
   minnames = 1,%
   giveninits = true,%
   doi = false,%
   url = false,%
   isbn = false%
}

\renewbibmacro{in:}{}

\AtEveryBibitem{%
   \ifentrytype{inproceedings}{%
      \clearname{editor}%
      \clearlist{editor}%
      \clearfield{editor}%
      \clearlist{publisher}%
      \clearfield{series}%
      \clearlist{location}%
      \clearfield{publisher}%
   }{}%
   \ifentrytype{book}{%
      \clearname{publisher}%
   }{}%
}
\AtEveryCitekey{%
   \ifentrytype{inproceedings}{%
      \clearlist{publisher}%
   }{}%
   \clearname{editor}%
   \clearlist{location}%
}
\DeclareFieldFormat[inproceedings]{date}{}
\DeclareFieldFormat[inproceedings]{pages}{#1, \printfield{year}}
\DeclareFieldFormat[inproceedings]{booktitle}{\mkbibemph{#1,\nopunct}}

\DeclareFieldFormat{journaltitle}{\mkbibemph{#1}\isdot}

\DeclareFieldFormat[article]{number}{\mkbibparens{#1}}
\DeclareFieldFormat[article]{date}{}
\DeclareFieldFormat[article]{pages}{#1, \printfield{year}}
\DeclareFieldFormat[article]{issue}{\nopunct#1}
\renewbibmacro*{issue+date}{%
   \iffieldundef{issue}{}%
   {%
      \printtext{\unspace}%
      (\printfield{issue})%
   }%
   \iffieldundef{pages}%
   {\addcomma\printfield{year}}%
   {}%
}

\renewcommand{\bibpagespunct}{%
   \ifentrytype{article}{\addcolon}{\addcomma}\space%
}

\DeclareFieldFormat[article]{titlecase}{\MakeSentenceCase{#1}}
\DeclareFieldFormat[inproceedings]{titlecase}{\MakeSentenceCase{#1}}
\DeclareFieldFormat[incollection]{titlecase}{\MakeSentenceCase{#1}}
\DeclareFieldFormat[book]{titlecase}{\MakeSentenceCase{#1}}

\renewbibmacro*{journal}{%
   \iffieldundef{journaltitle} {} {\printtext[journaltitle]{%
         \printfield[myplain]{journaltitle}%
         \setunit{\subtitlepunct}%
         \printfield[myplain]{journalsubtitle}}}}
\DeclareFieldFormat{myplain}{#1}

\renewbibmacro*{booktitle}{%
   \ifboolexpr{ test {\iffieldundef{booktitle}} and test
      {\iffieldundef{booksubtitle}} } {} {\printtext[booktitle]{%
         \printfield[myplain]{booktitle}%
         \setunit{\subtitlepunct}%
         \printfield[myplain]{booksubtitle}}%
      \newunit}%
   \printfield{booktitleaddon}}

\renewbibmacro*{volume+number+eid}{%
   \setunit*{\addcomma\addspace}%
   \printfield{volume}%
   \printfield{number}%
   \setunit{\addcomma\space}%
   \printfield{eid}%
}

\newbibmacro{string+doi}[1]{%
   \iffieldundef{doi}{%
      \iffieldundef{url}{%
         #1%
      }{%
         \href{\thefield{url}}{#1}%
      }%
   }{%
      \href{http://dx.doi.org/\thefield{doi}}{#1}}%
}
\DeclareFieldFormat*{title}{\usebibmacro{string+doi}{\mkbibemph{#1}}}
\renewcommand*{\multicitedelim}{\addcomma\space}
\DeclareFieldFormat[inproceedings]{title}{\usebibmacro{string+doi}{{#1}}}
\DeclareFieldFormat[article]{title}{\usebibmacro{string+doi}{{#1}}}
\DeclareFieldFormat[incollection]{title}{\usebibmacro{string+doi}{{#1}}}
\DeclareFieldFormat[book]{titlecase}{{#1}}

\DeclareCiteCommand{\tcite}%
{\usebibmacro{prenote}}%
{\usebibmacro{citeindex}\usebibmacro{cite}}%
{\multicitedelim}%
{\usebibmacro{postnote}}

%-----------------------------------------------------------------------------
% Hyperref & Visual Formatting
%-----------------------------------------------------------------------------
\usepackage{hyperref}
\hypersetup{%
   breaklinks,%
   colorlinks=true,%
   linkcolor=[rgb]{0.45,0.0,0.0},%
   citecolor=[rgb]{0,0,0.45}%
}

\definecolor{DeepNavy}{HTML}{002366}
\definecolor{RoyalBlue}{HTML}{004E92}
\definecolor{SteelBlue}{HTML}{4682B4}
\definecolor{TealBlue}{HTML}{008080}
\definecolor{DarkSlate}{HTML}{2F4F4F}
\definecolor{DustyGrey}{HTML}{696969}

\makeatletter
\AtBeginDocument{\let\@elt\relax}
\makeatother

\headstyles{bringhurst}
\chapterstyle{dash}

\renewcommand{\chapnamefont}{\normalfont\huge\bfseries\color{DeepNavy}}
\renewcommand{\chapnumfont}{\normalfont\huge\bfseries\color{DeepNavy}}
\renewcommand{\chaptitlefont}{\normalfont\Huge\bfseries\color{DeepNavy}}

\counterwithin{footnote}{chapter}

\setsubparaheadstyle{\normalfont\normalsize\itshape\color{DustyGrey}}
\setlength{\beforechapskip}{-20pt} 
\setlength{\afterchapskip}{30pt}
\setsecnumformat{\csname the#1\endcsname.\quad}

%-----------------------------------------------------------------------------
% Standalone Chapter vs Master Book Modes
%-----------------------------------------------------------------------------
\newcommand{\Copyright}{This work is licensed under the Creative Commons Attribution-Noncommercial 3.0 License. To view a copy of this license, visit \url{http://creativecommons.org/licenses/by-nc/3.0/} or send a letter to Creative Commons, 171 Second Street, Suite 300, San Francisco, California, 94105, USA.}

\newcommand{\RealChapterEndInner}{%
   {%
      \begingroup%
      \emergencystretch=1em%
      \printbibliography[heading=subbibliography,segment=\therefsegment]{}%
      \endgroup%
   }%
}

\providecommand{\BookMode}[1]{}
\providecommand{\StandAloneMode}[1]{#1}

\ifx\bookMode\undefined
   % STANDALONE CHAPTER MODE
   \renewcommand{\StandAloneMode}[1]{#1}
   \renewcommand{\BookMode}[1]{}

   \newcommand{\ChapterNumPage}[2]{%
      \if\relax\detokenize{#1}\relax
         \setcounter{chapter}{0}%
      \else\ifnum#1>0
         \setcounter{chapter}{\numexpr#1-1\relax}%
      \else
         \setcounter{chapter}{0}%
      \fi\fi
      \if\relax\detokenize{#2}\relax
         \setcounter{page}{1}%
      \else\ifnum#2>0
         \setcounter{page}{#2}%
      \else
         \setcounter{page}{1}%
      \fi\fi
   }

   \DeclareRobustCommand{\setheadingcolor}{\color{RoyalBlue}}
   \setsecheadstyle{\setheadingcolor\normalfont\Large\bfseries}
   \setsubsecheadstyle{\normalfont\large\bfseries\color{SteelBlue}}
   \setsubsubsecheadstyle{\normalfont\normalsize\bfseries\color{TealBlue}}
   \setparaheadstyle{\normalfont\normalsize\bfseries\color{DarkSlate}}

   \DeclareDocumentCommand{\Chapter}{s o m}{%
      \IfBooleanTF{#1}{%
         \IfValueTF{#2}{\chapter*[]{#2}}{\chapter*[]{#3}}%
      }{%
         \IfValueTF{#2}{\chapter*[]{#2}}{\chapter[]{#3}}%
      }%
      \setcounter{footnote}{0}%
   }
   \newcommand{\ChapterPrefix}{}
   \newcommand{\ChapterGen}{\pagestyle{plain}\vspace*{-0.7cm}}
   \newcommand{\ChapterEnd}{%
      \RealChapterEndInner{}%
      \end{document}%
   }

   \renewcommand{\File}[1]{#1}
   \renewcommand{\RootFile}[1]{../#1}
\else
   % MASTER BOOK MODE
   \renewcommand{\StandAloneMode}[1]{}
   \renewcommand{\BookMode}[1]{#1}

   \newcommand{\ChapterNumPage}[2]{}

   \DeclareDocumentCommand{\Chapter}{s o m}{%
      \IfBooleanTF{#1}{%
         \IfValueTF{#2}{\chapter*[#2]{#3}}{\chapter*{#3}}%
         \ifx\Dir\undefined\else
            \typeout{CHAPTER_MAP: \Dir/\CurrentChapFile => chap: 0, page: \thepage}%
         \fi
      }{%
         \IfValueTF{#2}{\chapter[#2]{#3}}{\chapter{#3}}%
         \ifx\Dir\undefined\else
            \typeout{CHAPTER_MAP: \Dir/\CurrentChapFile => chap: \thechapter, page: \thepage}%
         \fi
      }%
   }
   \newcommand{\ChapterPrefix}{%
      \vspace{-1cm}%
      \subsection*{\small
         {CS 574: Randomized Algorithms}\\
         {Sariel Har-Peled}\\
         {\today}}%
   } \newcommand{\ChapterGen}{}
   \newcommand{\ChapterEnd}{\RealChapterEndInner{}}

   \renewcommand{\File}[1]{\Dir/#1}
   \renewcommand{\RootFile}[1]{#1}

   \newcommand{\InputChap}[2]{%
      \def\Dir{#1}%
      \def\CurrentChapFile{#2}%
      \typeout{}%
      \typeout{------------------------------------------------------------}%
      \typeout{PROCESSING FILE: #1/#2}%
      \typeout{}%
      \input{#1/#2}%
   }
\fi

%-----------------------------------------------------------------------------
% Fragments & Quotes
%-----------------------------------------------------------------------------
% The PDF build writes fragment bodies to shared files and later includes them
% in standalone appendices.  A chapter webpage already contains each fragment
% at its original location, so TeX4ht only needs a transparent wrapper and can
% omit the repeated standalone copies.
\newenvironment{fragment}[1]{}{}
\DeclareDocumentCommand{\IncFragment}{m}{}

\newcommand{\Quote}[2]{%
   \SaveIndent
   \begin{center}%
      {\footnotesize
         \begin{quote}
             \RestoreIndent
             {#1}\\
             \hspace*{2cm}-- {#2.}
         \end{quote}%
      }%
      \index{quotation!#2}%
   \end{center}%
   \bigskip
}

\newcommand{\WebChapterQuote}[3]{%
   \ifvmode\IgnorePar\fi\EndP%
   \HCode{<figure class="chapter-quote"><blockquote>}%
      #1%
   \HCode{</blockquote><hr /><figcaption><span class="quote-author">}%
      #2%
   \HCode{</span><span class="quote-work">}%
      #3%
   \HCode{</span></figcaption></figure>}%
   \par%
}

\newcommand{\QuotePExt}[3]{\WebChapterQuote{#1}{#2}{#3}}

\epigraphfontsize{\small\itshape}
\setlength{\epigraphwidth}{0.6\textwidth}
\newcommand{\QuoteExt}[3]{\WebChapterQuote{#1}{#2}{#3}}

\providecommand{\delX}[1]{{\color{red}\sout{#1}}}
\providecommand{\newX}[1]{{\color{blue}#1}}
\providecommand{\chgY}[2]{{\color{red}\sout{#1}}{\color{blue}#2}}
\providecommand{\remX}[1]{{\color{purple}\textbf{[Remark: #1]}}}
\newcommand{\SpellIgnore}[1]{#1}
\newcommand{\EQ}{\mathrel{\;\;=\;\;}}
\newcommand{\LEQ}{\mathrel{\;\;\leq\;\;}}
\newcommand{\LT}{\mathrel{\;\;<\;\;}}
\newcommand{\GT}{\mathrel{\;\;<\;\;}}

\newcommand{\GEQ}{\mathrel{\;\;\geq\;\;}}

%-----------------------------------------------------------------------------
% Save a symbol that we know is going to get redefined.
\def\savesymbol#1{%
  \expandafter\let\expandafter\origsym\expandafter=\csname#1\endcsname
  \expandafter\let\csname orig#1\endcsname=\origsym
  \expandafter\let\csname#1\endcsname=\relax
}

% Restore a previously saved symbol, and rename the current one.
\def\restoresymbol#1#2{%
  \expandafter\let\expandafter\newsym\expandafter=\csname#2\endcsname
  \expandafter\global\expandafter\let\csname#1#2\endcsname=\newsym
  \expandafter\let\expandafter\origsym\expandafter=\csname orig#2\endcsname
  \expandafter\global\expandafter\let\csname#2\endcsname=\origsym
}

\usepackage{pictex}%
\savesymbol{triangle}%
\usepackage{tree}%
\restoresymbol{TreePKG}{triangle}%

% PicTeX's tree macros perform box arithmetic that TeX4ht cannot convert and
% can terminate the HTML document at the first tree.  These lightweight web
% forms preserve the two children and optional subtree label as HTML/CSS.
\renewcommand{\tree}[2]{%
   \HCode{<span class="web-binary-tree"><span class="tree-node"></span><span class="tree-children"><span class="tree-child">}%
   #1%
   \HCode{</span><span class="tree-child">}%
   #2%
   \HCode{</span></span></span>}%
}
\renewcommand{\TreePKGtriangle}[3][]{%
   \HCode{<span class="web-subtree-triangle"><span class="subtree-label">}%
   #1%
   \HCode{</span></span>}%
}

% Definitions absorbed from sbook.sty
\providecommand{\MdPrc}{\nu}

\providecommand{\TwoFiguresBase}[4]{%
   \begin{tabular}{cc}
      \begin{minipage}{0.48\linewidth}
          \begin{center}
              {{{#1}}}
          \end{center}
      \end{minipage}
      &
      \begin{minipage}{0.48\linewidth}
          \begin{center}
              {{#3}}
          \end{center}
      \end{minipage}
      \\
      {#2} & {#4}
   \end{tabular}%
}

\providecommand{\TwoFigures}[6]{%
   \begin{tabular}{cc}
      \begin{minipage}{0.48\linewidth}
          \begin{center}
              {\IncludeGraphics[#1]{#2}}
          \end{center}
      \end{minipage}
      &
      \begin{minipage}{0.48\linewidth}
          \centerline{\IncludeGraphics[#4]{#5}}
      \end{minipage}
      \\
      {#3} & {#6}
   \end{tabular}%
}

\providecommand{\ThreeFigures}[6]{%
   \begin{tabular}{ccc}
        \begin{minipage}{0.30\linewidth}
            \IncludeGraphics[width=5.1cm]{#1}\\
        \end{minipage}
        &
        \begin{minipage}{0.30\linewidth}
            \IncludeGraphics[width=5.1cm]{#2}\\
        \end{minipage}
        &
        \begin{minipage}{0.30\linewidth}
            \IncludeGraphics[width=5.1cm]{#3}\\
        \end{minipage}
        \vspace{-1cm}
        \\
        {#4} & {#5}& {#6}
   \end{tabular}%
}

%-----------------------------------------------------------------------------
% Shared notation and reference helpers; backend variants follow below.
\IfFileExists{shared/prefix_common.tex}{%-----------------------------------------------------------------------------
% Shared notation and reference helpers for PDF, lwarp, and TeX4ht.
% Loaded by each backend after its package and layout setup, before ntheorem.
% Keep rendering-dependent definitions in the corresponding backend prefix.
% Add new backend-independent notation here so all outputs share it.
%-----------------------------------------------------------------------------

\providecommand{\remove}[1]{}
\providecommand{\NotAMSMode}[1]{#1}
\providecommand{\Solution}[1]{}
\providecommand{\Version}[1]{\hfill{Version: #1}}
\providecommand{\Distor}{\mathrm{dist}}
\providecommand{\Hardness}[1]{(h:#1)}
\providecommand{\QuoteOptExt}[3]{\QuoteExt{#1}{#2}{#3}}
\providecommand{\Hypercube}{\EuScript{H}}
\providecommand{\CoverTime}{\EuScript{C}}
\providecommand{\Nbr}{\Gamma}
\providecommand{\CommuteTime}[2]{\mathbf{CT}_{#1#2}}
\providecommand{\Resistance}{\mathbf{R}}
\providecommand{\EffectResist}[2]{\mathbf{R}_{#1#2}}
\providecommand{\NetworkE}{\EuScript{N}}
\providecommand{\trinom}[4]{\binom{#1}{#2 \;\; #3 \;\; #4}}
\providecommand{\Ranges}{\EuScript{R}}
\providecommand{\CClass}[1]{\ensuremath{\mathbf{#1}}}
\providecommand{\coCClass}[1]{\ensuremath{\mathrm{co-}}\CClass{#1}}
\providecommand{\CX}{\EuScript{C}}
\providecommand{\packet}{\mathbf{c}}
\providecommand{\packetA}{\mathbf{d}}
\providecommand{\mopt}{m_{\mathrm{opt}}}
\providecommand{\girth}{\mathop{\mathrm{girth}}}
\providecommand{\CR}{\mathop{\mathrm{cr}}}
\providecommand{\TRUE}{{\textcolor[named]{RedViolet}{\textsf{TRUE}}}\xspace}
\providecommand{\FALSE}{{\textcolor[named]{RedViolet}{\textsf{FALSE}}}\xspace}
\providecommand{\Part}[1]{\part{#1}}
\providecommand{\Section}[1]{\section{#1}}
\providecommand{\SubSection}[1]{\subsection{#1}}
\providecommand{\SubSubSection}[1]{\subsubsection{#1}}
\providecommand{\Exercise}[5]{%
\begin{exercise_h}[#1]%
   \exelab{#3}{}%
   #4%
\end{exercise_h}%
}
\providecommand{\si}[1]{#1}
\providecommand{\myParagraph}[1]{\noindent\textbf{#1}}
\providecommand{\DefName}[1]{#1}
\providecommand{\deflabel}[1]{#1}
\providecommand{\Description}[1]{}

%-----------------------------------------------------------------------------
% Mathematical Notations & Environments (from notations.sty)
%-----------------------------------------------------------------------------
%-----------------------------------------------------------------------------
% List environments (questions & compactenum)
%-----------------------------------------------------------------------------

\newlist{compactenumA}{enumerate}{5}%
\setlist[compactenumA]{topsep=0pt,itemsep=-1ex,partopsep=1ex,parsep=1ex,%
   label=(\Alph*)}%

\newlist{compactenumI}{enumerate}{5}%
\setlist[compactenumI]{topsep=0pt,itemsep=-1ex,partopsep=1ex,parsep=1ex,%
   label=(\Roman*)}%

\newlist{compactenuma}{enumerate}{5}%
\setlist[compactenuma]{topsep=0pt,itemsep=-1ex,partopsep=1ex,parsep=1ex,%
   label=(\alph*)}%

\newlist{compactitem}{enumerate}{5}%
\setlist[compactitem]{topsep=0pt,itemsep=-1ex,partopsep=1ex,parsep=1ex,%
   label=\textbullet}%

\newlist{compactenumi}{enumerate}{5}%
\setlist[compactenumi]{topsep=0pt,itemsep=-1ex,partopsep=1ex,parsep=1ex,%
   label=(\roman*)}%

\setcounter{secnumdepth}{5}

\newcommand{\Term}[1]{\textsf{#1}}
\newcommand{\TermI}[1]{\Term{#1}\index{#1@\Term{#1}}}

\DeclareMathAlphabet{\mathpzc}{OT1}{pzc}{m}{it}

\newcommand{\seg}{\mathsf{s}}

\newcommand{\Diam}{\mathrm{diam}}
\newcommand{\diam}{\mathrm{diam}}

\newcommand{\Copt}{C_{\mathrm{opt}}}
\newcommand{\CenterPnt}{\mathrm{cen}}

\newcommand{\parent}{\overline{\mathrm{p}}}
\newcommand{\parentX}[2][\!]{\parent\pth[#1]{#2}}

\newcommand{\FastCut}{\Algorithm{FastCut}\xspace}
\newcommand{\MinCut}{\Algorithm{MinCut}\xspace}
\newcommand{\MinCutRep}{\Algorithm{MinCutRep}\xspace}
\newcommand{\Contract}{\Algorithm{Contract}\xspace}

\newcommand{\Grid}{\mathrm{G}}
\newcommand{\ClosestPair}{\mathcal{CP}}
\newcommand{\SphereChar}{\ensuremath{\mathbb{S}}}
\newcommand{\Sphere}[1]{{\SphereChar^{(#1)}}}

\newcommand{\PntSet}{{\mathsf{P}}}%
\providecommand{\PntSetA}{{\mathsf{Q}}}%

\newcommand{\query}{\mathcalb{q}}

\newcommand{\pnt}{\mathsf{p}}
\newcommand{\pntA}{\mathsf{u}}
\newcommand{\pntB}{\mathsf{v}}
\newcommand{\pntC}{{\mathsf{x}}}
\newcommand{\pntD}{{\mathsf{y}}}
\newcommand{\pntE}{{\mathsf{z}}}

\newcommand{\pntCoord}{p}
\newcommand{\pntCoordA}{q}

%%%%%%%%%%%%%%%%%%%%%%%%%%%%%%%%%%%%%%%%%%%%%%%%%%%%%%%%%%%%%%%%%%%%% 
%%%%%%%%%%%%%%%%%%%%%%%%%%%%%%%%%%%%%%%%%%%%%%%%%%%%%%%%%%%%%%%%%%%%%
% Markov chain notations.

\newcommand{\MCSpc}{\mathtt{S}}
\newcommand{\MCPrb}{\mathtt{P}}
\newcommand{\MCPrbY}[2]{\MCPrb_{#1#2}}
\newcommand{\MCPrbM}{\mathbf{P}}%
\newcommand{\MCPrbT}[3]{\MCPrb_{#1#2}^{(#3)}}
\newcommand{\MCFPrbT}[3]{\mathtt{r}_{#1#2}^{(#3)}}
\newcommand{\MCVstProb}[2]{\mathtt{f}_{#1#2}}
\newcommand{\MCExSteps}[2]{\mathtt{h}_{#1#2}}
\newcommand{\MCPVec}[1]{\mathbf{q}^{(#1)}}
%%%%%%%%%%%%%%%%%%%%%%%%%%%%%%%%%%%%%%%%%%%%%%%%%%%%%%%%%%%%%%%%%%%%%

\newcommand{\Spread}{\Phi}
\newcommand{\VC}{\ensuremath{\mathsf{VC}}\xspace}
\newcommand{\MC}{\ensuremath{\mathsf{MC}}\xspace}
\newcommand{\hyperplane}{h}

\newcommand{\QuickSort}{\Algorithm{QuickSort}\xspace}
\newcommand{\Alg}{\Algorithm{Alg}\xspace}

\newcommand{\Graph}{\mathsf{G}}%
\newcommand{\GraphB}{\mathsf{H}}
\newcommand{\GraphC}{\mathsf{K}}
\newcommand{\GraphD}{\mathsf{F}}

\providecommand{\G}{}
\renewcommand{\G}{\Graph}

\newcommand{\RLP}{\ComplexityClass{RLP}\xspace}
\newcommand{\USTCON}{\ComplexityClass{USTCON}}
\newcommand{\STCON}{\overrightarrow{\ComplexityClass{STCON}}}

\newcommand{\Mtrx}{\mathsf{M}}
\newcommand{\TMtrx}{\mathsf{Q}}
\newcommand{\IMtrx}{\mathcal{I}}

\newcommand{\Expander}[3]{$\left[#1,#2,#3\right]$-expander}

\newcommand{\ChartX}[1]{\ensuremath{\chi_{#1}}}
\newcommand{\Field}{\mathbb{F}}
\newcommand{\LSpace}{\mathbb{G}}
\newcommand{\LD}[2]{\mathsf{L{D}}\pth{#1,#2}}
\newcommand{\SEVN}{\widehat{\lambda_2}}
\newcommand{\NX}[2]{#1\!\left[#2\right]}
\newcommand{\Entry}[2]{#1\!\left[#2\right]}
\newcommand{\NXExt}[3]{#1_{#2}\!\left[#3\right]}
\newcommand{\RangeI}[1]{\llbracket #1 \rrbracket}
\newcommand{\DialX}[1]{\gamma\pth{#1}}
\newcommand{\BiDialX}[1]{{\gamma}_2\pth{#1}}

\newcommand{\EV}[1]{\widehat{\lambda_{#1}}}
\newcommand{\TEV}[1]{\widehat{\lambda_{#1}}}
\newcommand{\EVal}{\widehat{\lambda}}
\newcommand{\Basis}{\mathcal{B}}

\newcommand{\ZProd}[2]{\ensuremath{#1 \, \text{\textcircled{z}} #2}}
\newcommand{\RProd}[2]{\ensuremath{#1 \, \text{\textcircled{r}} #2}}

\newcommand{\OneVec}[1]{\mathbf{1}^{#1}}
\newcommand{\DstrbV}[1]{\mathbf{\mathsf{q}}^{(#1)}}
\newcommand{\DstrbE}[2]{\mathsf{q}^{(#1)}_{#2}}

\newcommand{\DEdges}{\widetilde{\Edges}}

\newcommand{\VV}{\mathsf{V}}%
\newcommand{\EE}{\mathsf{E}}%

\newcommand{\Vertices}{\VV}%
\newcommand{\Edges}{\EE}%

\newcommand{\EdgesX}[1]{\Edges\pth{#1}}
\newcommand{\VerticesX}[1]{V\pth{#1}}

\newcommand{\StringX}[1]{\mathsf{s}\pth{#1}}

\newcommand{\BMtrx}{\mathsf{B}}
\newcommand{\WMtrx}{\mathsf{W}}
\newcommand{\NWMtrx}{\overline{\mathsf{W}}}
\newcommand{\Sequence}{\EuScript{S}}

\newcommand{\med}{\mathrm{med}}
\newcommand{\medX}[1]{\mathrm{med}\pth{#1}}

\newcommand{\medianC}{\nu}%

\newlength{\savedparindent}
\newcommand{\SaveIndent}{\setlength{\savedparindent}{\parindent}}
\newcommand{\RestoreIndent}{\setlength{\parindent}{\savedparindent}}

\newcommand{\NNN}{\mathbb{N}}%
\providecommand{\NN}{\mathbb{N}}%
\newcommand{\reals}{\Re}%
\newcommand{\ZZ}{\mathbb{Z}}%

\newcommand{\level}{\mathop{\mathrm{level}}}

\providecommand{\TPDF}[2]{\texorpdfstring{#1}{#2}}

\newcommand{\Hint}[1]{\emph{[#1]}}
\newcommand{\Arr}{\mathop{\mathrm{\EuScript{A}}}}

\newcommand{\Array}{\mathcal{T}}%

\newcommand{\dirEdge}[2]{\left({{#1} \rightarrow {#2}} \right)}

%%%%%%%%%%%%%%%%%%%%%%%%%%%%%%%%%%%%%%%%%%%%%%%%%%%%%%%%%%%%%%%%%%%%%%
% Emphasize stuff...
\definecolor{blue25}{rgb}{0, 0, 11}
\newcommand{\emphic}[2]{%
   \textcolor{blue25}{%
      \bfseries{\emph{#1}}}%
   \index{#2}}
\newcommand{\emphi}[1]{\emphic{#1}{#1}}
\newcommand{\emphOnly}[1]{\textcolor{blue25}{\bfseries{\emph{#1}}}}
\providecommand{\emphiC}[1]{\emphic{#1}{#1}}

\definecolor{almostblack}{rgb}{0, 0, 0.3}
\providecommand{\emphw}[1]{{\textcolor{almostblack}{\emph{#1}}}}%

%
%%%%%%%%%%%%%%%%%%%%%%%%%%%%%%%%%%%%%%%%%%%%%%%%%%%%%%%%%%%%%%%%%%%%%%

\newcommand{\volX}[1]{\mathrm{vol}\pth{#1}}

\newcommand{\Interval}{\EuScript{I}}
\newcommand{\IntervalA}{\EuScript{I}}

\newcommand{\Exp}{\mathrm{Exp}}

\let\Ball\undefined
\newcommand{\Ball}{\mathbf{b}}
\newcommand{\hdim}{d}
\newcommand{\Domain}{\mathcal{D}}
\newcommand{\points}[1]{[\bfseries{#1 points}]}
\newcommand{\Event}{\EuScript{E}}%
\newcommand{\EventB}{\EuScript{B}}%
\newcommand{\EventF}{\EuScript{F}}%
\newcommand{\EventG}{\EuScript{G}}%

\newcommand{\distCmd}[1]{\left\| {#1}  \right\|}
\newcommand{\distY}[2]{\distCmd{#1 - #2}}

\newcommand{\NormalC}{\mathbf{N}}
\newcommand{\NormalY}[2]{\NormalC\pth{#1, #2}}%
\newcommand{\NormalDistOneDim}{\NormalC}
\newcommand{\NormalDistX}[1]{\NormalC^{#1}}
\DeclareMathOperator{\lca}{lca}
\newcommand{\RangeSpace}{\mathsf{S}}%
\newcommand{\RangeSpaceA}{\mathsf{T}}

\newcommand{\range}{\mathbf{r}}
\newcommand{\rangeA}{\mathbf{r}'}

\newcommand{\BadProb}{\varphi}

\newcommand{\EntropyChar}{\mathbb{H}}
\newcommand{\Entropy}[1]{\EntropyChar\pth{#1}}
\newcommand{\EntropyX}[1]{\EntropyChar\pth[]{#1}}

\newcommand{\Matrix}{\mathsf{M}}
\newcommand{\MatrixB}{\mathsf{B}}
\newcommand{\MatrixC}{\mathsf{C}}
\newcommand{\MatrixD}{\mathsf{D}}

\renewcommand{\th}{th\xspace}

\newcommand{\Family}{\mathcal{F}}
\newcommand{\FamilyB}{\mathcal{G}}
\newcommand{\FamilyC}{\mathcal{H}}
\newcommand{\FMSC}{\EuScript{Y}\index{metric space}}

\newcommand{\median}{\overline{\mathsf{m}}}
\newcommand{\QuickSelect}{\Algorithm{Quick{}Select}\xspace}

\providecommand{\row}{\mathbf{r}}%
\newcommand{\rowc}{{r}}%

\newcommand{\vecb}{\mathbf{b}}

\newcommand{\vecA}{\mathbf{v}}
\newcommand{\vecB}{\mathbf{u}}
\newcommand{\vecRand}{\mathbf{r}}

\newcommand{\vecH}{\mathbf{h}}
\newcommand{\vecV}{\mathbf{v}}

\newcommand{\vecX}{\mathbf{x}}
\newcommand{\vecY}{\mathbf{y}}

\newcommand{\DistBinom}[2]{\mathrm{Bin}\pth{#1,#2}}

\newcommand{\DistGeom}[1]{\mathrm{Geom}\pth{#1}}

\newcommand{\Sample}{\mathsf{R}}

\newcommand{\lcm}{\mathop{\mathrm{lcm}}}
\newcommand{\bdiv}{\mathop{\mathrm{div}}}%
\newcommand{\divides}{\mid}%
\newcommand{\poly}{\mathrm{poly}}%

\newcommand{\Group}{\mathcal{G}} 
\newcommand{\GroupA}{\mathcal{H}} 

\newcommand{\algGCD}{\Algorithm{EuclidGCD}\xspace}%

\newcommand{\gOp}{\times}%
\newcommand{\gIdent}{\mathsf{i}}%

\newcommand{\ordX}[2][\!]{\mathop{\mathrm{o{r}d}}\pth[#1]{#2}}

\newcommand{\lSym}[2]{%
   \left(
       \ds {#1}\divides{#2}%
   \right) %
}
\newcommand{\jSym}[2]{\left\llbracket%
       \ds {#1}\divides{#2}%
       \right\rrbracket%
}
\newcommand{\Line}{\ell}

\newcommand{\LevelChar}{\mathbb{L}}
\newcommand{\AtMostKLevelChar}{\LevelChar_{\leq k}}
\newcommand{\AtMostKLevel}[1]{\LevelChar_{\leq k} \pth{ #1}}

\newcommand{\LineSet}{\mathsf{L}}%
\newcommand{\ArrX}[1]{\EuScript{A}\pth{#1}}%

\newcommand{\ProblemC}[1]{{\textsf{\textcolor{red}{#1}\index{problem!#1}}}}
\DefineNamedColor{named}{OliveGreen} {cmyk}{0.64,0,0.95,0.40}
\providecommand{\ComplexityClass}[1]{{{\textcolor[named]{OliveGreen}{%
            \textsc{#1}}}}}
\providecommand{\NPHard}{{\ComplexityClass{NP-Hard}}%
   \index{NP!hard}\xspace}
\providecommand{\NPComplete}{\ComplexityClass{NP-Complete}%
   \index{NP!complete}\xspace}
\providecommand{\APXHard}{\ComplexityClass{APX-Hard}%
   \index{APX!Hard}\xspace}
\providecommand{\NPC}{\ComplexityClass{NPC}%
   \index{NP!complete}\xspace}
\providecommand{\NP}{\ComplexityClass{NP}%
   \index{NP}\xspace}
\providecommand{\POLYT}{\ComplexityClass{P}\xspace}

\newcommand{\envFrame}[1]{\begin{center}\fbox{\begin{minipage}[c]{0.75\textwidth}{#1}\end{minipage}}\end{center}}

\newcommand{\mysqrt}[1]{\ts\!\!\sqrt{#1}}
\newcommand{\SegSet}{\mathsf{S}}
\newcommand{\ArrVD}[1]{\Arr^{|} \pth{#1}}
\newcommand{\RIC}{\textsf{RIC}\index{random incremental
      construction}\xspace}
\newcommand{\ArrB}{\EuScript{B}}

\newcommand{\TreeOpt}{\EuScript{T}_\mathrm{opt}}
\newcommand{\TreeOptA}{\EuScript{T}_\mathrm{opt}'}

\newcommand{\DepthExt}[2]{\mathrm{depth}_{#1}\pth{#2}}
\newcommand{\depthX}[1]{\mathrm{depth}\pth{#1}}

\newcommand{\TreeA}{\EuScript{U}}
\newcommand{\TreeB}{\EuScript{V}}

\newcommand{\Trap}{\sigma}
\newcommand{\TrapA}{\sigma'}
\newcommand{\TrapB}{\sigma''}

\newcommand{\SampleA}{\mathsf{R}'}

\newcommand{\AllRegions}{\mathcal{T}}
\newcommand{\RegionsSet}[1]{\mathcal{F}\pth{#1}}
\newcommand{\RegionsHeavy}[2]{\mathcal{F}_{\!\geq#1}\pth{#2}}
\newcommand{\NumOfRegionsZero}[1]{\ExChar \!f\pth{#1}}
\newcommand{\NumOfRegions}[2]{\ExChar\! f_{\!\geq #1}\pth[]{#2}}
\newcommand{\ObjSet}{\mathsf{S}}

\newcommand{\WeightX}[2][\!]{\omega \pth[#1]{#2}}
\newcommand{\Edge}{\mathsf{e}}

\newcommand{\Face}{\mathsf{f}}

\newcommand{\DefSet}[1]{D\pth{#1}}
\newcommand{\KillSet}[1]{K\pth{#1}}

\newcommand{\CSProbFull}[4]{
   \,\rho_{\!\!\stackrel{\rule[-.005cm]{0cm}{0.1cm}}{#1,
       #4}}\pth{\MakeSBig #2,#3}}

\newcommand{\disk}{D}
\newcommand{\diskA}{D'}

\newcommand{\IntervalSet}{\widehat{\mathcal{I}}}

\newcommand{\GroundSet}{\textsf{X}}
\newcommand{\FGroundSet}{\mathsf{x}}
\newcommand{\RangeSet}{{\mathcal{R}}}
\newcommand{\RangeSetA}{\mathcal{R}'}
\newcommand{\MeasureChar}{\overline{m}}
\newcommand{\Measure}[1]{\MeasureChar\pth{#1}}
\newcommand{\sMeasure}[1]{\overline{s}\pth{#1}}

\newcommand{\SetA}{S}%
\newcommand{\SetB}{T}%
\newcommand{\SetC}{N}%

\newcommand{\VCProj}[2]{#1_{|#2}}
\newcommand{\VCDim}{\mathrm{dim}_{\VC}}
\newcommand{\Dim}{\delta}
\newcommand{\DimA}{\delta'}
\newcommand{\PntSetB}{{{\mathsf{Q}}}}
\newcommand{\PntSetC}{{{\mathsf{R}}}}
\newcommand{\Growth}[2]{{\EuScript{G}}_{#1}\pth{#2}}
\newcommand{\ShatterFunc}[1]{\pi_{#1}}
\newcommand{\Distribution}{\EuScript{D}}
\newcommand{\Error}{\EuScript{E}}

\newcommand{\bd}{{\partial}}%
\newcommand{\eps}{{\varepsilon}}%

\newcommand{\clmlab}[1]{\label{claim:#1}}
\newcommand{\clmref}[1]{\HLink{Claim}{claim:#1}}

\newcommand{\CHX}[1]{{\mathcal{CH}}\pth{#1}}

\newcommand{\ts}{\hspace{0.6pt}}

\def\qedsymbol{\myqedsymbol}%

\newcommand{\qedsign}{\hfill{\hfill\myqedsymbol}}
\newcommand{\qedsignBig}{\hfill{\hfill\rule{4mm}{4mm}}}
\newcommand{\QED}{{\hfill{\hfill\qedsymbol}}}

\newcommand{\num}{\alpha}

\newcommand{\NEW}[1]{#1}

\newcommand{\IncludeGraphics}[2][]{%
   \typeout{}%
   \typeout{Graphics: #2}%
   \typeout{\ includegraphics[#1]{#2}}%
   \includegraphics[#1]{#2}
   \typeout{}%
}

%    -----------------------------------------------------------------
%    -----------------------------------------------------------------
%    -----------------------------------------------------------------
%    Handling references
%    -----------------------------------------------------------------
%    -----------------------------------------------------------------
%    -----------------------------------------------------------------
\newcommand{\HLinkShort}[2]{\hyperref[#2]{#1\ref*{#2}}}
\newcommand{\HLink}[2]{\hyperref[#2]{#1~\ref*{#2}}}
\newcommand{\HLinkPage}[2]{\hyperref[#2]{#1~\ref*{#2}%
      $_\text{p\pageref{#2}}$}}

\newcommand{\HLinkSuffix}[3]{\hyperref[#2]{#1\ref*{#2}{#3}}}
\newcommand{\HLinkPageSuffix}[3]{\hyperref[#2]{#1\ref*{#2}%
      #3$_\text{p\pageref{#2}}$}}

\newcommand{\figlab}[1]{\label{fig:#1}}
\newcommand{\figref}[1]{\HLink{Figure}{fig:#1}}

\newcommand{\exelab}[1]{\label{exer:#1}}

\providecommand{\deflab}[1]{\label{def:#1}}
\newcommand{\defref}[1]{\HLink{Definition}{def:#1}}
\newcommand{\defrefY}[2]{\hyperref[def:#2]{#1}}
\newcommand{\defrefpage}[1]{\HLinkPage{Definition}{def:#1}}

\newcommand{\etal}{\textit{et~al.}\xspace}

\newcommand{\obslab}[1]{\label{observation:#1}}
\newcommand{\obsref}[1]{\HLink{Observation}{observation:#1}}
\newcommand{\obsrefY}[2]{\hyperref[observation:#2]{#1}}

\newcommand{\exmlab}[1]{\label{example:#1}}
\newcommand{\exmref}[1]{\HLink{Example}{example:#1}}

\newcommand{\remlab}[1]{\label{rem:#1}}
\newcommand{\remref}[1]{\HLink{Remark}{rem:#1}}%

\newcommand{\chaplab}[1]{\label{chap:#1}}

\newcommand{\itemlab}[1]{\label{item:#1}}
\newcommand{\itemref}[1]{\HLinkSuffix{}{item:#1}{}}

\newcommand{\tedlab}[1]{\label{tedium:#1}}
\newcommand{\tedref}[1]{\HLink{Tedium}{tedium:#1}}%

% Deferred tedium: `by \TediumHere{...}, we have ...` inside a proof.
% Use \TediumHereY{...}{label} to allow later citations with \tedref{label}.
% Bodies may contain paragraphs and displays. They are typeset, in call order,
% after the enclosing proof, using ordinary tedium numbering and references.
% As with other forward references, compile twice to resolve the links.
% Define macros needed by a body outside the proof: its local scope has ended
% when the deferred body is typeset. Verbatim material is not supported.
\ExplSyntaxOn
\int_new:N \g__notes_tedium_proof_int
\int_new:N \g__notes_tedium_item_int
\seq_new:N \g__notes_tedium_stack_seq
\tl_new:N \l__notes_tedium_proof_tl

\cs_new_protected:Npn \__notes_tedium_begin:
  {
    \int_gincr:N \g__notes_tedium_proof_int
    \seq_new:c
      { g__notes_tedium_\int_use:N \g__notes_tedium_proof_int _seq }
    \seq_gput_right:Nx \g__notes_tedium_stack_seq
      { \int_use:N \g__notes_tedium_proof_int }
  }

\cs_new_protected:Npn \__notes_tedium_print:n #1
  {
    \seq_use:cn { g__notes_tedium_#1_seq } { }
    \cs_undefine:c { g__notes_tedium_#1_seq }
  }

\cs_new_protected:Npn \__notes_tedium_end:
  {
    \seq_gpop_right:NN \g__notes_tedium_stack_seq
      \l__notes_tedium_proof_tl
    \exp_args:NV \__notes_tedium_print:n \l__notes_tedium_proof_tl
  }

\cs_new_protected:Npn \__notes_tedium_queue:nn #1#2
  {
    \seq_if_empty:NTF \g__notes_tedium_stack_seq
      {
        \PackageError{notes}
          {Deferred~tedium~used~outside~a~proof}
          {Use~\string\TediumHere\space or~\string\TediumHereY
           \space inside~a~proof~environment.}
      }
      {
        \seq_get_right:NN \g__notes_tedium_stack_seq
          \l__notes_tedium_proof_tl
        \seq_gput_right:cx
          { g__notes_tedium_\l__notes_tedium_proof_tl _seq }
          {
            \exp_not:N \begin {tedium}
            \exp_not:N \tedlab
              { \exp_not:n {#2} }
            \exp_not:n {#1}
            \exp_not:N \end {tedium}
          }
        \tedref{#2}
      }
  }

\NewDocumentCommand{\TediumHere}{+m}
  {
    \int_gincr:N \g__notes_tedium_item_int
    \exp_args:Nnx \__notes_tedium_queue:nn {#1}
      { auto-td-\int_use:N \g__notes_tedium_item_int }
  }

\NewDocumentCommand{\TediumHereY}{+m m}
  { \__notes_tedium_queue:nn {#1}{#2} }

% Each proof has its own queue; the stack also keeps nested proofs separate.
\clist_map_inline:nn
  { proof,altproof,proofextin,proofnoqed,proofof,proofExt,proofnoproof }
  {
    \AddToHook{env/#1/before}[notes/tedium]{\__notes_tedium_begin:}
    \AddToHook{env/#1/after}[notes/tedium]{\__notes_tedium_end:}
  }
\ExplSyntaxOff

\newcommand{\factlab}[1]{\label{fact:#1}}%

\newcommand{\alglab}[1]{\label{algorithm:#1}}

\newcommand{\net}{\mathcal{N}}%

\newcommand{\lemlab}[1]{\label{lemma:#1}}
\newcommand{\lemref}[1]{\HLink{Lemma}{lemma:#1}}%
\newcommand{\lemrefY}[2]{\hyperref[lemma:#2]{#1}}
\newcommand{\lemrefpage}[1]{\HLinkPage{Lemma}{lemma:#1}}
\newcommand{\lemrefshort}[1]{\HLinkShort{L}{lemma:#1}}

\newcommand{\problab}[1]{\label{prob:#1}}%
\newcommand{\probref}[1]{\HLink{Problem}{prob:#1}}%

\newcommand{\eqlab}[1]{\label{equation:#1}}
\newcommand{\Eqrefbare}[1]{\HLinkSuffix{(}{equation:#1}{)}}
\newcommand{\Eqref}[1]{\HLinkSuffix{Eq.~(}{equation:#1}{)}}
\newcommand{\Eqrefpage}[1]{\HLinkPageSuffix{Eq.~(}{equation:#1}{)}}

\newcommand{\seclab}[1]{\label{sec:#1}}
\newcommand{\secref}[1]{\HLink{Section}{sec:#1}}
\newcommand{\secrefY}[2]{\hyperref[sec:#2]{#1}}
\newcommand{\secrefpage}[1]{\HLinkPage{Section}{sec:#1}}

\newcommand{\thmlab}[1]{{\label{theo:#1}}}
\newcommand{\thmref}[1]{\HLink{Theorem}{theo:#1}}
\newcommand{\thmrefY}[2]{\hyperref[theo:#2]{#1}}
\newcommand{\thmrefpage}[1]{\HLinkPage{Theorem}{theo:#1}}

\newcommand{\corlab}[1]{\label{cor:#1}}
\newcommand{\corref}[1]{\HLink{Corollary}{cor:#1}}%
\newcommand{\correfY}[2]{\hyperref[cor:#2]{#1}}

\newcommand{\Yend}{Y_{\mathrm{end}}}%

%    -----------------------------------------------------------------
%    -----------------------------------------------------------------
%    -----------------------------------------------------------------
%    Handling references
%    -----------------------------------------------------------------
%    -----------------------------------------------------------------
%    -----------------------------------------------------------------

\newcommand{\ceil}[1]{\left\lceil {#1} \right\rceil}
\newcommand{\floor}[1]{\left\lfloor {#1} \right\rfloor}
\newcommand{\brc}[1]{\left\{ {#1} \right\}}
\newcommand{\pth}[2][\!]{\mleft({#2}\mright)}%
\newcommand{\brk}[1]{\mleft[{#1}\mright]}%
\newcommand{\pbrc}[2][\!\!]{#1\left[ {#2} \bigr. \right]}
\newcommand{\pbrcS}[1]{\left[ {#1} \right]}

\newcommand{\dotprod}[2]{{#1}\cdot{#2}}
\newcommand{\DotProd}[2]{\permut{{#1},{#2}}}
\newcommand{\permut}[1]{\left\langle {#1} \right\rangle}

\newcommand{\MakeBig}{}
\newcommand{\MakesBig}{}
\newcommand{\MakeSBig}{}
\newcommand{\MakeVBig}{}

\newcommand{\DistChar}{\mathsf{d}}%
\newcommand{\DistSetY}[2]{\DistChar\pth{#1,#2}}
\newcommand{\distG}[2]{\DistChar_\Graph\pth{#1, #2}}
\newcommand{\DistMtrCharX}[1]{\DistChar_{#1}}

\newcommand{\dY}[2]{\DistChar\pth{#1, #2}}%
\newcommand{\dGY}[2]{\DistChar_\Graph\pth{#1, #2}}%

\newcommand{\cset}{\mathcal{C}}%

\newcommand{\MTR}{\mathcal{M}}
\newcommand{\DistMtr}[2]{\DistChar\pth{#1,#2}}

\newcommand{\dist}[1]{\left\| {#1} \right\|}
\newcommand{\norm}[1]{\left\| {#1} \right\|}

\providecommand{\ds}{\displaystyle}
\newcommand{\sep}[1]{\,\left|\, {#1} \bigr.\right.}%
\newcommand{\sepw}[1]{\left|\, {#1} \right.}%

\newcommand{\ProbExt}[2]{\mathop{\mathbf{Pr}}_{#1}\pbrcx{#2}}
\newcommand{\ProbLTR}{\mathbb{P}}%

\newcommand{\Prob}[1]{\mathop{\ProbLTR} \mleft[ #1 \mright]}%
\newcommand{\ProbCond}[2]{\mathop{\ProbLTR}\!\left[%
       #1 \;\middle\vert\; #2 \right]}
\newcommand{\ProbY}[2]{\ProbLTR^{}_{\!#1}\pbrc{#2}}

\newcommand{\ExChar}{\mathbb{E}}%
\newcommand{\ExSym}{\mathop{\ExChar}}%
\newcommand{\Ex}[2][\!]{\ExSym#1\pbrcx{#2}}
\newcommand{\ExExt}[2]{\ExSym_{#1}\!\pbrcx{#2}}
\newcommand{\ExExtS}[2]{\ExSym_{\substack{#1}}\!\pbrcx{#2}}
\newcommand{\ExCond}[2]{\ExSym\!\left[%
       #1 \;\middle\vert\; #2 \right]}

\newcommand{\VarSymb}{\mathop{\mathbb{V}}}
\newcommand{\Var}[1]{\mathop{\mathbb{V}}\!\pbrcx{#1}}

\newcommand{\pbrcx}[1]{\left[ {#1} \right]}%
\newcommand{\nfrac}[2]{{#1}/{#2}}
\newcommand{\cardin}[1]{\left| {#1} \right|}%

\newcommand\CCancel[2][black]{{\renewcommand\CancelColor{\color{#1}}\cancel{#2}}}

\newcommand{\IncGraphPage}[4][]{%
   \IncludeGraphics[page=#4,#1]{\File{#2/#3}}%
}

\newenvironment{enumEx}{
   \begin{compactenumA}
    }{\end{compactenumA}}

\newcommand{\cl}{{\mathrm{cl}}}

\newcommand{\HSetY}[2]{H\pth{#1, #2}}%
\newcommand{\CHSetY}[2]{C\pth{#1, #2}}%
\newcommand{\OldOm}{{\mathcal{O}}}
\newcommand{\NewOm}{{\mathcal{N}}}

%%%%%%%%%%%%%%%%%%%%%%%%%%%%%%%%%%%%%%%%%%%%%%%%%%%%%%%%%%%%%%%%%%%%%%%
%---------------------------------------------------------------
\newcommand{\IntRange}[1]{\left\llbracket #1 \right\rrbracket}
\newcommand{\IRX}[1]{\left\llbracket #1 \right\rrbracket}%
\newcommand{\IRY}[2]{\left\llbracket #1:#2 \right\rrbracket}

%---------------------------------------------------------------
%%%%%%%%%%%%%%%%%%%%%%%%%%%%%%%%%%%%%%%%%%%%%%%%%%%%%%%%%%%%%%%%%%%%%%%

\newcommand{\obj}{f}%
\newcommand{\Wg}{\omega}%
\newcommand{\TotalW}{W}%

\newcommand{\rankY}[2]{{#1}_{ \permut{#2}}}%%
\newcommand{\areaX}[1]{\mathrm{area}\pth{#1}}%%

\newcommand{\mybinom}[2]%
{\Bigl(\begin{array}{@{}c@{}}#1\\#2\end{array}\Bigr)}

\providecommand{\Set}[2]{\left\{ #1 \;\middle\vert\; #2 \right\}}
\newcommand*{\biga}[1]{\vcenter{\hbox{\scalebox{1.5}{\ensuremath#1}}}}
\newcommand{\CNF}{\TermI{CNF}\xspace}%
\newcommand{\Tree}{T}%

\DeclareSymbolFont{sfletters}{OML}{cmbrm}{m}{it}

\DeclareMathSymbol{\salpha}{\mathord}{sfletters}{"0B}

%%%%%%%%%%%%%%%%%%%%%%%%%%%%%%%%%%%%%%%%%%%%%%%%%%%%%%%%%%%%%%%%%%%%

\providecommand{\ProblemInd}[2]{\textsf{\textcolor{red}{#1}\index{Problem!#2}}}
\newcommand{\NPProblemInd}[4]{
   \begin{center}
      \vspace{0.6cm}
       \fbox{
          \begin{minipage}{0.9\linewidth}
              \vspace{-0.6cm} \hspace{-0.35cm}
              {\large \bfseries \ProblemInd{#1}{#2}}
              \vspace{+0.2cm}%

              \noindent{\textcolor{blue}{\bfseries \textsf{Instance}}:} {#3}
              
              \noindent{\textcolor{blue}{\bfseries \textsf{Question}:}} {#4}
          \end{minipage}
       }
   \end{center}
   \smallskip
   }

\newcommand{\NPProblemOpt}[4]{\NPProblemInd{#1}{#2}{#3}{#4}}

%%%%%%%%%%%%%%%%%%%%%%%%%%%%%%%%%%%%%%%%%%%%%%%%%%%%%%%%%%%%%%%%%%%%

\newcommand{\Hastad}{H{\aa}stad\xspace}

\newcommand{\ANNbr}{\Term{ANN}\index{approximate nearest neighbor}\xspace}
\newcommand{\LSH}{\TermI{L{S}H}\xspace}

\newcommand{\DS}{\EuScript{D}} % DataStructure
\newcommand{\DSC}{\mathsf{D}} % DataStructure

\newcommand{\distH}[2]{\DistChar_{H}\pth{#1, #2}}
\newcommand{\distHY}[2]{\DistChar_{H}\pth{#1, #2}}

\newcommand{\DSAprxNearNbr}{\ensuremath{\DataStructure{{\approx%
         \text{Near{}}}}}%
   \index{near neighbor!data-structure!approximate}%
}
\newcommand{\DataStructure}[1]{\ensuremath{\DS_{#1}}}
\newcommand{\DistNN}[2]{\DistChar\pth{#1,#2}}
\newcommand{\pntc}{{p}}%

\newcommand{\seq}{{M}}%
\newcommand{\cTimes}{\beta}%

\newcommand{\DSTimes}{L}%
\newcommand{\rad}{r}%

\newcommand{\HypercubeX}[1]{\Hypercube^{#1}}
\newcommand{\kFunctions}[2]{\mathrm{combine}\pth{#1,#2}}
\newcommand{\WEqualChar}{\Bumpeq}
\newcommand{\WEqual}[2]{{#1}\WEqualChar{#2}}
\newcommand{\FamilyAmp}{\FamilyC_{\WEqualChar}}
\newcommand{\PSPACE}{\ComplexityClass{PSPACE}\xspace}%

\newcommand{\PS}{P}%

\newcommand{\PSetFar}{\ensuremath{\PntSet^{}_{\geq}\xspace}}%
\newcommand{\DA}{D}%
\newcommand{\uniqX}[1]{\mathrm{uniq}\pth{#1}}%
\newcommand{\diffX}[1]{\mathop{}\!\mathrm{d}#1}

\newcommand{\Frq}{F}%
\newcommand{\frq}{f}%
\newcommand{\STRM}{\mathcal{S}}%
\newcommand{\strm}{{s}}%

\newcommand{\Calinescu}{C\u{a}linescu\xspace}

\newcommand{\fGrowX}[1]{G^{}_{\Dim}\pth{#1}}
\newcommand{\CompressFunc}{\mathtt{Compress}}
\newcommand{\RSA}{\Term{RSA}\xspace}
\newcommand{\MST}{\Term{MST}\xspace}
\providecommand{\CClassPP}{\CClass{P{}P}\xspace}
\providecommand{\ID}{\mathrm{id}}

\newcommand{\EncodeFunc}{\text{\ttfamily{Enc}}}
\newcommand{\DecodeFunc}{\text{\ttfamily{Dec}}}
\newcommand{\discX}[1]{\mathrm{disc}\pth{#1}}%

\newcommand{\propertyC}{\psi}%
\newcommand{\propertyX}[1]{\propertyC\pth{#1}}%
\newcommand{\numC}{\Gamma}%

\newcommand{\degreeX}[1]{\mathrm{degree}\pth{#1}}

\newcommand{\ind}{m}

\providecommand{\nequiv}{\not\equiv}
\providecommand{\sqcupplus}{\sqcup^+}

\newcommand{\AlgorithmI}[1]{{\protect\color{RedViolet}{#1}}}

\newcommand{\binsY}[2]{\sqcup_{\geq #1}\pth{#2}}

\newcommand{\BLX}[1]{{{\sqcupplus{}}_{\geq #1}}}

\newcommand{\pdf}{\textsf{pdf}\xspace}

\newcommand{\cdf}{\textsf{cdf}\xspace}

\newcommand{\costX}[1]{\mathrm{cost}\pth{#1}}

\newcommand{\codeX}[1]{\mathrm{code}\pth{#1}}

\newcommand{\lengthX}[1]{\mathrm{l{}e{}n}\pth{#1}}
}{\input{prefix_common}}%

\newcommand{\HST}{\TermI{HST}\xspace}

\newcommand{\FMS}{\EuScript{X}\index{metric space}}  

\newcommand{\FMSB}{\EuScript{Y}\index{metric space}}  

\renewcommand{\Re}{\mathbb{R}}%

\providecommand{\Matousek}{Matou{\v s}ek\xspace}

\providecommand{\Erdos}{Erd{\H o}s\xspace}

\providecommand{\Lovasz}{Lov\'asz\xspace}

\newcommand{\ThreeFiguresExt}[9]{%
   \begin{tabular}{ccc}
     \begin{minipage}{0.31\linewidth}
         \begin{center}
             {\IncludeGraphics[#1]{#2}}
         \end{center}
     \end{minipage}
     &
       \begin{minipage}{0.31\linewidth}
           \centerline{\IncludeGraphics[#4]{#5}}
       \end{minipage}
     &
       \begin{minipage}{0.31\linewidth}
           \centerline{\IncludeGraphics[#7]{#8}}
       \end{minipage}
     \\
     ~\\[-0.2cm]
     {#3} & {#6} & {#9}
   \end{tabular}}

\newcommand{\exeref}[1]{Exercise~\ref{exer:#1}}

% Theorem and proof rendering
\usepackage[amsmath,thmmarks]{ntheorem}
\theoremseparator{.}%

\newtheorem{theorem:unnumbered}{Theorem}[section]
\newtheorem{theorem}{Theorem}[section]

\newtheorem{observation}[theorem]{Observation}
\newtheorem{conjecture}[theorem]{Conjecture}
\newtheorem{lemma}[theorem]{Lemma}
\newtheorem{corollary}[theorem]{Corollary}
\newtheorem{claim}[theorem]{Claim}
\newtheorem{assumption}[theorem]{Assumption}
\newtheorem{tedium}[theorem]{Tedium}

\theoremstyle{plain}%
\theoremheaderfont{\sffamily}
\theorembodyfont{\upshape}%

\newtheorem{defn}[theorem]{Definition}
\newtheorem{definition}[theorem]{Definition}

\newtheorem*{remark_unnumbered}[theorem:unnumbered]{Remark}%
\newtheorem*{remarks}[theorem]{Remarks}%
\newtheorem{remark}[theorem]{Remark}%
\newtheorem{fact}[theorem]{Fact}%
\newtheorem{example}[theorem]{Example}
\newtheorem{exercise}[theorem]{Exercise}
\newtheorem{problem}[theorem]{Problem}
\newtheorem{algorithm}[theorem]{Algorithm}%
\newtheorem{xca}[theorem]{Exercise}
\newtheorem{openproblem}[theorem]{Open Problem}
\newtheorem{exercise_h}[theorem]{Exercise}

% Proof environments
\theoremheaderfont{\em}%
\theorembodyfont{\upshape}%
\theoremstyle{nonumberplain}
\theoremseparator{}
% TeX4ht's ntheorem support consumes theorem end symbols.  The web-only
% wrappers below place the duck before ntheorem closes its HTML block.
\theoremsymbol{}
\newtheorem{proof}{Proof:}
\newtheorem{proofextin}{Proof}
\newtheorem{altproof}{Alternative proof:}

\theoremsymbol{}
\newtheorem{proofnoqed}{Proof:}

\theoremsymbol{}
\theorempreskip{0pt}%
\theorempostskip{0pt}%
\newtheorem{proofof}{Proof of\!}%

\makeatletter
\let\proof@theorem\proof
\let\endproof@theorem\endproof
\renewenvironment{proof}{\proof@theorem}{\unskip\myqedsymbol\endproof@theorem}
\let\proofextin@theorem\proofextin
\let\endproofextin@theorem\endproofextin
\renewenvironment{proofextin}{\proofextin@theorem}{\unskip\myqedsymbol\endproofextin@theorem}
\let\altproof@theorem\altproof
\let\endaltproof@theorem\endaltproof
\renewenvironment{altproof}{\altproof@theorem}{\unskip\myqedsymbol\endaltproof@theorem}
\let\proofof@theorem\proofof
\let\endproofof@theorem\endproofof
\renewenvironment{proofof}{\proofof@theorem}{\unskip\myqedsymbol\endproofof@theorem}
\makeatother

\newenvironment{proofExt}[1]{{\emph{#1:}}}{\hfill{\hfill\myqedsymbol}}
\newenvironment{proofnoproof}{{}}{{\hfill{\hfill\myqedsymbol}}}

%%%%%%%%%%%%%%%%%%%%%%%%%%%%%%%%%%%%%%%%%%%%%%%%%%%%%%%%%%%%%%%%%%%%
%%%%%%%%%%%%%%%%%%%%%%%%%%%%%%%%%%%%%%%%%%%%%%%%%%%%%%%%%%%%%%%%%%%%
% Macros to embed figures in text

\newcommand{\EBRY}[2]{#1_{\langle#2\rangle}}%

\newcommand{\NiceFrac}[2]{%
   {\raisebox{0.2em}{\ensuremath{#1}}\!\!\!}%
   {\rotatebox[origin=c]{-20}{\Big/}}% Rotated /
   {\! \raisebox{-0.2em}{\ensuremath{#2}} }}%

% end - Macros to embed figures in text
%%%%%%%%%%%%%%%%%%%%%%%%%%%%%%%%%%%%%%%%%%%%%%%%%%%%%%%%%%%%%%%%%%%% 
%%%%%%%%%%%%%%%%%%%%%%%%%%%%%%%%%%%%%%%%%%%%%%%%%%%%%%%%%%%%%%%%%%%%

%%%%%%%%%%%%%%%%%%%%%%%%%%%%%%%%%%%%%%%%%%%%%%%%%%%%%%%%%%%%%%%%%%
% Color in math
%\providecommand{\Mh}[1]{{#1}}%
\providecommand{\IfColorMath}[2]{#2}%

\IfFileExists{.latex_printer_friendly}{\def\GenPrinterVer{1}}{}%

\ifx\GenPrinterVer\undefined
   \IfFileExists{.latex_color}{\def\GenColorMath{1}}{}
\else
   \renewcommand{\IfPrinterVer}[2]{#1}%
\fi

\ifx\GenColorMath\undefined
\else
\renewcommand{\IfColorMath}[2]{#1}%
%\renewcommand{\Mh}[1]{{\textcolor{red}{#1}}}%
\fi
% Color in math - end
%%%%%%%%%%%%%%%%%%%%%%%%%%%%%%%%%%%%%%%%%%%%%%%%%%%%%%%%%%%%%%%%%%

\newcommand{\Holder}{H\"older\xspace}

\providecommand{\Turan}{Tur\'an\xspace}

\usepackage{marginnote}
\usepackage{ragged2e} 
% Short, subtle snark macro
\setlength{\marginparwidth}{1.5in}
\setlength{\marginparsep}{10pt}
\newcommand{\snark}[1]{%
   \marginnote{\raggedright\scriptsize\sffamily\color{gray}#1}%
}

\newcommand{\ballsY}[2]{\text{\Football}_{\geq #1}\pth{#2}}

% Page-breaking hints have no meaning in HTML.  Display bodies are preserved
% verbatim for MathJax, so the builder also removes occurrences found there.
\renewcommand{\allowdisplaybreaks}[1][]{}

\let\Mh\undefined

% -----------------------------------------------------------------------------
% Bibliography & Document Begin
%-----------------------------------------------------------------------------
\addbibresource{shared/book.bib}%
\makeindex
\begin{document}
}{\input{prefix_4ht}}%
  \else
    \IfFileExists{shared/prefix_latex.tex}{%-----------------------------------------------------------------------------
% Master Prefix File for Randomized Algorithms Class Notes & Book
% Unified self-contained configuration for LuaLaTeX and XeLaTeX with Memoir
%-----------------------------------------------------------------------------

\RequirePackage{iftex}
\ifPDFTeX
  \PackageError{engine-guard}{This document requires LuaLaTeX or XeLaTeX}{%
    Please recompile using LuaLaTeX or XeLaTeX instead of pdfLaTeX.%
  }
  \stop
\fi

\ifdefined\LWR@started
  \PackageError{latex-gate}{%
    Attempting to run lwarp with prefix_latex.tex.%
  }{Please use prefix_lwarp.tex (or set \string\warpMode) for HTML generation.}
  \stop
\fi

% Always use memoir
\documentclass[12pt,oneside,fleqn]{memoir}

%-----------------------------------------------------------------------------
% Typography, Fonts & Math Packages
%-----------------------------------------------------------------------------
\usepackage{amssymb}
\usepackage{mathtools}
\PassOptionsToPackage{warnings-off={mathtools-overbracket,mathtools-colon}}{unicode-math}
\usepackage{fontspec}
\usepackage{unicode-math}

% New Computer Modern OpenType text and math fonts
\setmainfont{NewCM10-Book.otf}[
  BoldFont          = NewCM10-Bold.otf,
  ItalicFont        = NewCM10-BookItalic.otf,
  BoldItalicFont    = NewCM10-BoldItalic.otf,
  SlantedFont       = NewCM10-Book.otf,
  SlantedFeatures   = {FakeSlant = 0.25},
  BoldSlantedFont   = NewCM10-Bold.otf,
  BoldSlantedFeatures = {FakeSlant = 0.25},
  SmallCapsFeatures = {Letters = SmallCaps}
]
\setsansfont{NewCMSans10-Book.otf}[
  BoldFont          = NewCMSans10-Bold.otf,
  ItalicFont        = NewCMSans10-BookOblique.otf,
  BoldItalicFont    = NewCMSans10-BoldOblique.otf
]
\setmonofont{NewCMMono10-Book.otf}[
  BoldFont          = NewCMMono10-Bold.otf,
  ItalicFont        = NewCMMono10-BookItalic.otf,
  BoldItalicFont    = NewCMMono10-BoldOblique.otf
]

\setmathfont{NewCMMath-Book.otf}
\setmathfont{NewCMMath-Bold.otf}[version=bold]

% OpenType math aliases and fixes
\providecommand{\square}{\mdlgwhtsquare}
\def\Box{\mdlgwhtsquare}
\def\Diamond{\mdlgwhtdiamond}
\def\blacksquare{\mdlgblksquare}
\def\leadsto{\rightsquigarrow}
\providecommand{\llbracket}{\lBrack}
\providecommand{\rrbracket}{\rBrack}
\DeclareTextCommand{\textcircled}{TU}[1]{%
  \mbox{\ooalign{\hfil\raise0.1ex\hbox{#1}\hfil\crcr\hbox{\ensuremath{\bigcirc}}}}%
}

\usepackage{alphabeta}
\usepackage[normalem]{ulem} % Provides \sout
\usepackage{mleftright}
\usepackage{microtype}
\usepackage{xspace}
\usepackage{xurl}
\usepackage[dvipsnames,table]{xcolor}
\usepackage{graphicx}
\usepackage{caption}
\usepackage{booktabs,array,colortbl,hhline,multirow,tabularx}
\usepackage[inline]{enumitem}
\usepackage[thicklines]{cancel}
\renewcommand\CancelColor{\color{purple}}
\usepackage{euscript}
\usepackage{marvosym}
\let\Cross\undefined
\let\Square\undefined
\usepackage{bbding}
\usepackage{epigraph}
\usepackage{makeidx}
\usepackage{perpage}
\usepackage{tikz}

%-----------------------------------------------------------------------------
% Boondox / mathcalb font declaration
%-----------------------------------------------------------------------------
\DeclareFontFamily{U}{BOONDOX-calo}{\skewchar\font=45 }
\DeclareFontShape{U}{BOONDOX-calo}{m}{n}{
  <-> s*[1.05] BOONDOX-r-calo}{}
\DeclareFontShape{U}{BOONDOX-calo}{b}{n}{
  <-> s*[1.05] BOONDOX-b-calo}{}
\DeclareMathAlphabet{\mathcalb}{U}{BOONDOX-calo}{m}{n}
\SetMathAlphabet{\mathcalb}{bold}{U}{BOONDOX-calo}{b}{n}
\DeclareMathAlphabet{\mathbcalb}{U}{BOONDOX-calo}{b}{n}

%-----------------------------------------------------------------------------
% Algorithm Environment (formerly salgorithm.sty)
%-----------------------------------------------------------------------------
\newcommand{\ProgWordX}[1]{{\small\textbf{#1}}{\ }}
\newcommand{\Do}{\ProgWordX{do}\xspace}
\newcommand{\End}{\ProgWordX{end}\xspace}
\newcommand{\Return}{\ProgWordX{return}\xspace}%
\newcommand{\Procbegin}{\ProgWordX{begin}}
\newcommand{\If}{\ProgWordX{if}\xspace}
\newcommand{\Then}{\ProgWordX{then}\xspace}
\newcommand{\Else}{\ProgWordX{else}\xspace}
\newcommand{\For}{\ProgWordX{for}\xspace}
\newcommand{\While}{{\small\bfseries while}\ }

\newcommand{\xbeginlgox}{\begin{minipage}{1in}\begin{tabbing}
           \quad\=\qquad\=\qquad\=\qquad\=\qquad\=\qquad\=\qquad\=\kill}
\newcommand{\xendlgox}{\end{tabbing}\end{minipage}}

\newenvironment{algorithmEnv}{%
   \begin{tabular}{|l|}\hline\xbeginlgox\bigskip}%
   {\xendlgox\\\hline\end{tabular}}

\newenvironment{program}{%
   \begin{minipage}{4.0in}%
   \begin{tabbing}%
       \ \ \ \ \=\ \ \ \=\ \ \ \ \=\ \ \ \ \=\ \ \ \ \=%
      \ \ \ \ \=\ \ \ \ \=\ \ \ \ \=\ \ \ \ \=%
      \ \ \ \ \=\ \ \ \ \=\ \ \ \ \=\ \ \ \ \=\kill%
}{%
   \end{tabbing}%
   \end{minipage}%
}

\newcommand{\AlgorithmBox}[1]{%
   \fbox{\begin{program} #1 \end{program}}%
}

\definecolor{AlgorithmColor}{cmyk}{0.07,0.90,0,0.34}
\providecommand{\Algorithm}[1]{{\AlgorithmI{#1}\index{algorithm!#1@{\AlgorithmI{#1}}}}}
\newcommand{\CodeComment}[1]{\textcolor{blue}{\texttt{#1}}}

%-----------------------------------------------------------------------------
% Path & Directory Resolution
%-----------------------------------------------------------------------------
\newcommand{\RootDIR}{../}
\IfFileExists{../../../book.tex}{\renewcommand{\RootDIR}{../../../}}{}
\IfFileExists{../../book.tex}{\renewcommand{\RootDIR}{../../}}{}
\IfFileExists{../book.tex}{\renewcommand{\RootDIR}{../}}{}
\IfFileExists{./book.tex}{\renewcommand{\RootDIR}{}}{}

\providecommand{\File}[1]{#1}
\providecommand{\RootFile}[1]{\RootDIR #1}

%-----------------------------------------------------------------------------
% QED Duck Symbol
%-----------------------------------------------------------------------------
\newlength{\CharHeight}
\AtBeginDocument{\setlength{\CharHeight}{\fontcharht\font`X}}

\IfFileExists{shared/qed_duck.pdf}{%
   \providecommand{\myqedsymbol}{%
      \reflectbox{\includegraphics[height=1.4\CharHeight]{shared/qed_duck.pdf}}%
   }%
}{%
   \providecommand{\myqedsymbol}{$\blacksquare$}%
}

%-----------------------------------------------------------------------------
% BibLaTeX Setup (Pure Biber, No BibTeX)
%-----------------------------------------------------------------------------
\usepackage[bibencoding=utf8,style=alphabetic,refsegment=chapter,
indexing=cite,
backend=biber]{biblatex}

\ExecuteBibliographyOptions{%
   maxnames = 40,%
   maxalphanames = 4,%
   maxcitenames = 3,%
   minalphanames = 3,%
   minnames = 1,%
   giveninits = true,%
   doi = false,%
   url = false,%
   isbn = false%
}

\renewbibmacro{in:}{}

\AtEveryBibitem{%
   \ifentrytype{inproceedings}{%
      \clearname{editor}%
      \clearlist{editor}%
      \clearfield{editor}%
      \clearlist{publisher}%
      \clearfield{series}%
      \clearlist{location}%
      \clearfield{publisher}%
   }{}%
   \ifentrytype{book}{%
      \clearname{publisher}%
   }{}%
}
\AtEveryCitekey{%
   \ifentrytype{inproceedings}{%
      \clearlist{publisher}%
   }{}%
   \clearname{editor}%
   \clearlist{location}%
}
\DeclareFieldFormat[inproceedings]{date}{}
\DeclareFieldFormat[inproceedings]{pages}{#1, \printfield{year}}
\DeclareFieldFormat[inproceedings]{booktitle}{\mkbibemph{#1,\nopunct}}

\DeclareFieldFormat{journaltitle}{\mkbibemph{#1}\isdot}

\DeclareFieldFormat[article]{number}{\mkbibparens{#1}}
\DeclareFieldFormat[article]{date}{}
\DeclareFieldFormat[article]{pages}{#1, \printfield{year}}
\DeclareFieldFormat[article]{issue}{\nopunct#1}
\renewbibmacro*{issue+date}{%
   \iffieldundef{issue}{}%
   {%
      \printtext{\unspace}%
      (\printfield{issue})%
   }%
   \iffieldundef{pages}%
   {\addcomma\printfield{year}}%
   {}%
}

\renewcommand{\bibpagespunct}{%
   \ifentrytype{article}{\addcolon}{\addcomma}\space%
}

\DeclareFieldFormat[article]{titlecase}{\MakeSentenceCase{#1}}
\DeclareFieldFormat[inproceedings]{titlecase}{\MakeSentenceCase{#1}}
\DeclareFieldFormat[incollection]{titlecase}{\MakeSentenceCase{#1}}
\DeclareFieldFormat[book]{titlecase}{\MakeSentenceCase{#1}}

\renewbibmacro*{journal}{%
   \iffieldundef{journaltitle} {} {\printtext[journaltitle]{%
         \printfield[myplain]{journaltitle}%
         \setunit{\subtitlepunct}%
         \printfield[myplain]{journalsubtitle}}}}
\DeclareFieldFormat{myplain}{#1}

\renewbibmacro*{booktitle}{%
   \ifboolexpr{ test {\iffieldundef{booktitle}} and test
      {\iffieldundef{booksubtitle}} } {} {\printtext[booktitle]{%
         \printfield[myplain]{booktitle}%
         \setunit{\subtitlepunct}%
         \printfield[myplain]{booksubtitle}}%
      \newunit}%
   \printfield{booktitleaddon}}

\renewbibmacro*{volume+number+eid}{%
   \setunit*{\addcomma\addspace}%
   \printfield{volume}%
   \printfield{number}%
   \setunit{\addcomma\space}%
   \printfield{eid}%
}

\newbibmacro{string+doi}[1]{%
   \iffieldundef{doi}{%
      \iffieldundef{url}{%
         #1%
      }{%
         \href{\thefield{url}}{#1}%
      }%
   }{%
      \href{http://dx.doi.org/\thefield{doi}}{#1}}%
}
\DeclareFieldFormat*{title}{\usebibmacro{string+doi}{\mkbibemph{#1}}}
\renewcommand*{\multicitedelim}{\addcomma\space}
\DeclareFieldFormat[inproceedings]{title}{\usebibmacro{string+doi}{{#1}}}
\DeclareFieldFormat[article]{title}{\usebibmacro{string+doi}{{#1}}}
\DeclareFieldFormat[incollection]{title}{\usebibmacro{string+doi}{{#1}}}
\DeclareFieldFormat[book]{titlecase}{{#1}}

\DeclareCiteCommand{\tcite}%
{\usebibmacro{prenote}}%
{\usebibmacro{citeindex}\usebibmacro{cite}}%
{\multicitedelim}%
{\usebibmacro{postnote}}

%-----------------------------------------------------------------------------
% Hyperref & Visual Formatting
%-----------------------------------------------------------------------------
\usepackage{hyperref}
\hypersetup{%
   breaklinks,%
   colorlinks=true,%
   linkcolor=[rgb]{0.45,0.0,0.0},%
   citecolor=[rgb]{0,0,0.45}%
}

\definecolor{DeepNavy}{HTML}{002366}
\definecolor{RoyalBlue}{HTML}{004E92}
\definecolor{SteelBlue}{HTML}{4682B4}
\definecolor{TealBlue}{HTML}{008080}
\definecolor{DarkSlate}{HTML}{2F4F4F}
\definecolor{DustyGrey}{HTML}{696969}

\makeatletter
\AtBeginDocument{\let\@elt\relax}
\makeatother

\headstyles{bringhurst}
\chapterstyle{dash}

\renewcommand{\chapnamefont}{\normalfont\huge\bfseries\color{DeepNavy}}
\renewcommand{\chapnumfont}{\normalfont\huge\bfseries\color{DeepNavy}}
\renewcommand{\chaptitlefont}{\normalfont\Huge\bfseries\color{DeepNavy}}

\counterwithin{footnote}{chapter}

\setsubparaheadstyle{\normalfont\normalsize\itshape\color{DustyGrey}}
\setlength{\beforechapskip}{-20pt} 
\setlength{\afterchapskip}{30pt}
\setsecnumformat{\csname the#1\endcsname.\quad}

%-----------------------------------------------------------------------------
% Standalone Chapter vs Master Book Modes
%-----------------------------------------------------------------------------
\newcommand{\Copyright}{This work is licensed under the Creative Commons Attribution-Noncommercial 3.0 License. To view a copy of this license, visit \url{http://creativecommons.org/licenses/by-nc/3.0/} or send a letter to Creative Commons, 171 Second Street, Suite 300, San Francisco, California, 94105, USA.}

\newcommand{\RealChapterEndInner}{%
   {%
      \begingroup%
      \emergencystretch=1em%
      \printbibliography[heading=subbibliography,segment=\therefsegment]{}%
      \endgroup%
   }%
}

\providecommand{\BookMode}[1]{}
\providecommand{\StandAloneMode}[1]{#1}

\ifx\bookMode\undefined
   % STANDALONE CHAPTER MODE
   \renewcommand{\StandAloneMode}[1]{#1}
   \renewcommand{\BookMode}[1]{}

   \newcommand{\ChapterNumPage}[2]{%
      \if\relax\detokenize{#1}\relax
         \setcounter{chapter}{0}%
      \else\ifnum#1>0
         \setcounter{chapter}{\numexpr#1-1\relax}%
      \else
         \setcounter{chapter}{0}%
      \fi\fi
      \if\relax\detokenize{#2}\relax
         \setcounter{page}{1}%
      \else\ifnum#2>0
         \setcounter{page}{#2}%
      \else
         \setcounter{page}{1}%
      \fi\fi
   }

   \DeclareRobustCommand{\setheadingcolor}{\color{RoyalBlue}}
   \setsecheadstyle{\setheadingcolor\normalfont\Large\bfseries}
   \setsubsecheadstyle{\normalfont\large\bfseries\color{SteelBlue}}
   \setsubsubsecheadstyle{\normalfont\normalsize\bfseries\color{TealBlue}}
   \setparaheadstyle{\normalfont\normalsize\bfseries\color{DarkSlate}}

   \DeclareDocumentCommand{\Chapter}{s o m}{%
      \IfBooleanTF{#1}{%
         \IfValueTF{#2}{\chapter*[]{#2}}{\chapter*[]{#3}}%
      }{%
         \IfValueTF{#2}{\chapter*[]{#2}}{\chapter[]{#3}}%
      }%
      \setcounter{footnote}{0}%
   }
   \newcommand{\ChapterPrefix}{}
   \newcommand{\ChapterGen}{\pagestyle{plain}\vspace*{-0.7cm}}
   \newcommand{\ChapterEnd}{%
      \RealChapterEndInner{}%
      \end{document}%
   }

   \renewcommand{\File}[1]{#1}
   \renewcommand{\RootFile}[1]{../#1}
\else
   % MASTER BOOK MODE
   \renewcommand{\StandAloneMode}[1]{}
   \renewcommand{\BookMode}[1]{#1}

   \newcommand{\ChapterNumPage}[2]{}

   \DeclareDocumentCommand{\Chapter}{s o m}{%
      \IfBooleanTF{#1}{%
         \IfValueTF{#2}{\chapter*[#2]{#3}}{\chapter*{#3}}%
         \ifx\Dir\undefined\else
            \typeout{CHAPTER_MAP: \Dir/\CurrentChapFile => chap: 0, page: \thepage}%
         \fi
      }{%
         \IfValueTF{#2}{\chapter[#2]{#3}}{\chapter{#3}}%
         \ifx\Dir\undefined\else
            \typeout{CHAPTER_MAP: \Dir/\CurrentChapFile => chap: \thechapter, page: \thepage}%
         \fi
      }%
   }
   \newcommand{\ChapterPrefix}{%
      \vspace{-1cm}%
      \subsection*{\small
         {CS 574: Randomized Algorithms}\\
         {Sariel Har-Peled}\\
         {\today}}%
   } \newcommand{\ChapterGen}{}
   \newcommand{\ChapterEnd}{\RealChapterEndInner{}}

   \renewcommand{\File}[1]{\Dir/#1}
   \renewcommand{\RootFile}[1]{#1}

   \newcommand{\InputChap}[2]{%
      \def\Dir{#1}%
      \def\CurrentChapFile{#2}%
      \typeout{}%
      \typeout{------------------------------------------------------------}%
      \typeout{PROCESSING FILE: #1/#2}%
      \typeout{}%
      \input{#1/#2}%
   }
\fi

%-----------------------------------------------------------------------------
% Fragments & Quotes
%-----------------------------------------------------------------------------
\newcommand{\InputExt}[1]{%
   \IfFileExists{#1}%
   {\input{#1}}%
   {\IfFileExists{../#1}%
      {\input{../#1}}%
      {\IfFileExists{../../#1}%
         {\input{../../#1}}%
         {\typeout{ERROR: Unable to find #1}}%
      }%
   }%
}
\makeatletter
\def\FragNormKey#1{%
  \frag@norm@check#1def_\relax{#1}%
}
\def\frag@norm@check#1def_#2\relax#3{%
  \if\relax\detokenize{#1}\relax
    #3%
  \else
    def_#3%
  \fi
}
\makeatother

\newcommand{\FragmentName}[1]{shared/fragment/\FragNormKey{#1}.tex}

\newenvironment{fragment}[1]{}{}

\makeatletter
\AtBeginDocument{%
   \let\real@latex@label\label
   \def\label#1{%
      \global\@namedef{all@labels@seen@#1}{1}%
      \real@latex@label{#1}%
   }%
}

\DeclareDocumentCommand{\IncFragment}{m}{%
   \begingroup
      \def\label##1{%
         \@ifundefined{all@labels@seen@##1}{%
            \global\@namedef{all@labels@seen@##1}{1}%
            \real@latex@label{##1}%
         }{}%
      }%
      \IfFileExists{\FragmentName{#1}}%
         {\InputExt{\FragmentName{#1}}}%
         {\InputExt{shared/fragment/#1.tex}}%
   \endgroup
}
\makeatother

\newcommand{\Quote}[2]{%
   \SaveIndent
   \begin{center}%
      {\footnotesize
         \begin{quote}
             \RestoreIndent
             {#1}\\
             \hspace*{2cm}-- {#2.}
         \end{quote}%
      }%
      \index{quotation!#2}%
   \end{center}%
   \bigskip
}

\newcommand{\QuotePExt}[3]{%
   \bigskip%
   \hfill%
   \begin{minipage}{0.95\linewidth}%
       \footnotesize {#1}%

       \rule{\linewidth}{0.3mm}%
       \hfill #2, #3%
   \end{minipage}%
   \bigskip%
}

\epigraphfontsize{\small\itshape}
\setlength{\epigraphwidth}{0.6\textwidth}
\newcommand{\QuoteExt}[3]{%
   \vspace*{-1.6cm}%
   \epigraph{{\footnotesize #1}}{{#2}, {#3}}%
   \vspace*{-1.2cm}%
}

\providecommand{\delX}[1]{{\color{red}\sout{#1}}}
\providecommand{\newX}[1]{{\color{blue}#1}}
\providecommand{\chgY}[2]{{\color{red}\sout{#1}}{\color{blue}#2}}
\providecommand{\remX}[1]{{\color{purple}\textbf{[Remark: #1]}}}
\newcommand{\SpellIgnore}[1]{#1}
\newcommand{\EQ}{\mathrel{\;\;=\;\;}}
\newcommand{\LEQ}{\mathrel{\;\;\leq\;\;}}
\newcommand{\LT}{\mathrel{\;\;<\;\;}}
\newcommand{\GT}{\mathrel{\;\;<\;\;}}

\newcommand{\GEQ}{\mathrel{\;\;\geq\;\;}}

%-----------------------------------------------------------------------------
% Save a symbol that we know is going to get redefined.
\def\savesymbol#1{%
  \expandafter\let\expandafter\origsym\expandafter=\csname#1\endcsname
  \expandafter\let\csname orig#1\endcsname=\origsym
  \expandafter\let\csname#1\endcsname=\relax
}

% Restore a previously saved symbol, and rename the current one.
\def\restoresymbol#1#2{%
  \expandafter\let\expandafter\newsym\expandafter=\csname#2\endcsname
  \expandafter\global\expandafter\let\csname#1#2\endcsname=\newsym
  \expandafter\let\expandafter\origsym\expandafter=\csname orig#2\endcsname
  \expandafter\global\expandafter\let\csname#2\endcsname=\origsym
}

\usepackage{pictex}%
\savesymbol{triangle}%
\usepackage{tree}%
\restoresymbol{TreePKG}{triangle}%

% Definitions absorbed from sbook.sty
\providecommand{\MdPrc}{\nu}

\providecommand{\TwoFiguresBase}[4]{%
   \begin{tabular}{cc}
      \begin{minipage}{0.47\linewidth}
          \begin{center}
              {{{#1}}}
          \end{center}
      \end{minipage}
      &
      \begin{minipage}{0.47\linewidth}
          \begin{center}
              {{#3}}
          \end{center}
      \end{minipage}
      \\
      {#2} & {#4}
   \end{tabular}%
}

\providecommand{\TwoFigures}[6]{%
   \begin{tabular}{cc}
      \begin{minipage}{0.48\linewidth}
          \begin{center}
              {\IncludeGraphics[#1]{#2}}
          \end{center}
      \end{minipage}
      &
      \begin{minipage}{0.48\linewidth}
          \centerline{\IncludeGraphics[#4]{#5}}
      \end{minipage}
      \\
      {#3} & {#6}
   \end{tabular}%
}

\providecommand{\ThreeFigures}[6]{%
   \begin{tabular}{ccc}
        \begin{minipage}{0.30\linewidth}
            \IncludeGraphics[width=5.1cm]{#1}\\
        \end{minipage}
        &
        \begin{minipage}{0.30\linewidth}
            \IncludeGraphics[width=5.1cm]{#2}\\
        \end{minipage}
        &
        \begin{minipage}{0.30\linewidth}
            \IncludeGraphics[width=5.1cm]{#3}\\
        \end{minipage}
        \vspace{-1cm}
        \\
        {#4} & {#5}& {#6}
   \end{tabular}%
}

%-----------------------------------------------------------------------------
\providecommand{\remove}[1]{}
\providecommand{\NotAMSMode}[1]{#1}
\providecommand{\Solution}[1]{}
\providecommand{\Version}[1]{\hfill{Version: #1}}
\providecommand{\Distor}{\mathrm{dist}}
\providecommand{\Hardness}[1]{(h:#1)}
\providecommand{\QuoteOptExt}[3]{\QuoteExt{#1}{#2}{#3}}
\providecommand{\Hypercube}{\EuScript{H}}
\providecommand{\CoverTime}{\EuScript{C}}
\providecommand{\Nbr}{\Gamma}
\providecommand{\CommuteTime}[2]{\mathbf{CT}_{#1#2}}
\providecommand{\Resistance}{\mathbf{R}}
\providecommand{\EffectResist}[2]{\mathbf{R}_{#1#2}}
\providecommand{\NetworkE}{\EuScript{N}}
\providecommand{\trinom}[4]{\binom{#1}{#2 \;\; #3 \;\; #4}}
\providecommand{\Ranges}{\EuScript{R}}
\providecommand{\CClass}[1]{\ensuremath{\mathbf{#1}}}
\providecommand{\coCClass}[1]{\ensuremath{\mathrm{co-}}\CClass{#1}}
\providecommand{\CX}{\EuScript{C}}
\providecommand{\packet}{\mathbf{c}}
\providecommand{\packetA}{\mathbf{d}}
\providecommand{\mopt}{m_{\mathrm{opt}}}
\providecommand{\girth}{\mathop{\mathrm{girth}}}
\providecommand{\CR}{\mathop{\mathrm{cr}}}
\providecommand{\TRUE}{{\textcolor[named]{RedViolet}{\textsf{TRUE}}}\xspace}
\providecommand{\FALSE}{{\textcolor[named]{RedViolet}{\textsf{FALSE}}}\xspace}
\providecommand{\Part}[1]{\part{#1}}
\providecommand{\Section}[1]{\section{#1}}
\providecommand{\SubSection}[1]{\subsection{#1}}
\providecommand{\SubSubSection}[1]{\subsubsection{#1}}
\providecommand{\Exercise}[5]{%
\begin{exercise_h}[#1]%
   \exelab{#3}{}%
   #4%
\end{exercise_h}%
}
\providecommand{\si}[1]{#1}
\providecommand{\myParagraph}[1]{\noindent\textbf{#1}}
\providecommand{\DefName}[1]{#1}
\providecommand{\deflabel}[1]{#1}
\providecommand{\Description}[1]{}

%-----------------------------------------------------------------------------
% Mathematical Notations & Environments (from notations.sty)
%-----------------------------------------------------------------------------
%-----------------------------------------------------------------------------
% List environments (questions & compactenum)
%-----------------------------------------------------------------------------

\newlist{compactenumA}{enumerate}{5}%
\setlist[compactenumA]{topsep=0pt,itemsep=-1ex,partopsep=1ex,parsep=1ex,%
   label=(\Alph*)}%

\newlist{compactenumI}{enumerate}{5}%
\setlist[compactenumI]{topsep=0pt,itemsep=-1ex,partopsep=1ex,parsep=1ex,%
   label=(\Roman*)}%

\newlist{compactenuma}{enumerate}{5}%
\setlist[compactenuma]{topsep=0pt,itemsep=-1ex,partopsep=1ex,parsep=1ex,%
   label=(\alph*)}%

\newlist{compactitem}{enumerate}{5}%
\setlist[compactitem]{topsep=0pt,itemsep=-1ex,partopsep=1ex,parsep=1ex,%
   label=\textbullet}%

\newlist{compactenumi}{enumerate}{5}%
\setlist[compactenumi]{topsep=0pt,itemsep=-1ex,partopsep=1ex,parsep=1ex,%
   label=(\roman*)}%

\setcounter{secnumdepth}{5}

\newcommand{\HST}{\TermI{HST}\xspace}
\newcommand{\Term}[1]{\textsf{#1}}
\newcommand{\TermI}[1]{\Term{#1}\index{#1@\Term{#1}}}

\DeclareMathAlphabet{\mathpzc}{OT1}{pzc}{m}{it}

\newcommand{\seg}{\mathsf{s}}

\newcommand{\Diam}{\mathrm{diam}}
\newcommand{\diam}{\mathrm{diam}}

\newcommand{\Copt}{C_{\mathrm{opt}}}
\newcommand{\CenterPnt}{\mathrm{cen}}

\newcommand{\parent}{\overline{\mathrm{p}}}
\newcommand{\parentX}[2][\!]{\parent\pth[#1]{#2}}

\newcommand{\FMS}{\EuScript{X}\index{metric space}}  
\newcommand{\FMSB}{\EuScript{Y}\index{metric space}}  
\newcommand{\FastCut}{\Algorithm{FastCut}\xspace}
\newcommand{\MinCut}{\Algorithm{MinCut}\xspace}
\newcommand{\MinCutRep}{\Algorithm{MinCutRep}\xspace}
\newcommand{\Contract}{\Algorithm{Contract}\xspace}

\newcommand{\Grid}{\mathrm{G}}
\newcommand{\ClosestPair}{\mathcal{CP}}
\newcommand{\SphereChar}{\ensuremath{\mathbb{S}}}
\newcommand{\Sphere}[1]{{\SphereChar^{(#1)}}}

\newcommand{\PntSet}{{\mathsf{P}}}%
\providecommand{\PntSetA}{{\mathsf{Q}}}%

\newcommand{\query}{\mathcalb{q}}

\newcommand{\pnt}{\mathsf{p}}
\newcommand{\pntA}{\mathsf{u}}
\newcommand{\pntB}{\mathsf{v}}
\newcommand{\pntC}{{\mathsf{x}}}
\newcommand{\pntD}{{\mathsf{y}}}
\newcommand{\pntE}{{\mathsf{z}}}

\newcommand{\pntCoord}{p}
\newcommand{\pntCoordA}{q}

%%%%%%%%%%%%%%%%%%%%%%%%%%%%%%%%%%%%%%%%%%%%%%%%%%%%%%%%%%%%%%%%%%%%% 
%%%%%%%%%%%%%%%%%%%%%%%%%%%%%%%%%%%%%%%%%%%%%%%%%%%%%%%%%%%%%%%%%%%%%
% Markov chain notations.

\newcommand{\MCSpc}{\mathtt{S}}
\newcommand{\MCPrb}{\mathtt{P}}
\newcommand{\MCPrbY}[2]{\MCPrb_{#1#2}}
\newcommand{\MCPrbM}{\mathbf{P}}%
\newcommand{\MCPrbT}[3]{\MCPrb_{#1#2}^{(#3)}}
\newcommand{\MCFPrbT}[3]{\mathtt{r}_{#1#2}^{(#3)}}
\newcommand{\MCVstProb}[2]{\mathtt{f}_{#1#2}}
\newcommand{\MCExSteps}[2]{\mathtt{h}_{#1#2}}
\newcommand{\MCPVec}[1]{\mathbf{q}^{(#1)}}
%%%%%%%%%%%%%%%%%%%%%%%%%%%%%%%%%%%%%%%%%%%%%%%%%%%%%%%%%%%%%%%%%%%%%

\newcommand{\Spread}{\Phi}
\newcommand{\VC}{\ensuremath{\mathsf{VC}}\xspace}
\newcommand{\MC}{\ensuremath{\mathsf{MC}}\xspace}
\newcommand{\hyperplane}{h}

\newcommand{\QuickSort}{\Algorithm{QuickSort}\xspace}
\newcommand{\Alg}{\Algorithm{Alg}\xspace}

\newcommand{\Graph}{\mathsf{G}}%
\newcommand{\GraphB}{\mathsf{H}}
\newcommand{\GraphC}{\mathsf{K}}
\newcommand{\GraphD}{\mathsf{F}}

\providecommand{\G}{}
\renewcommand{\G}{\Graph}

\newcommand{\RLP}{\ComplexityClass{RLP}\xspace}
\newcommand{\USTCON}{\ComplexityClass{USTCON}}
\newcommand{\STCON}{\overrightarrow{\ComplexityClass{STCON}}}

\newcommand{\Mtrx}{\mathsf{M}}
\newcommand{\TMtrx}{\mathsf{Q}}
\newcommand{\IMtrx}{\mathcal{I}}

\newcommand{\Expander}[3]{$\left[#1,#2,#3\right]$-expander}

\newcommand{\ChartX}[1]{\ensuremath{\chi_{#1}}}
\newcommand{\Field}{\mathbb{F}}
\newcommand{\LSpace}{\mathbb{G}}
\newcommand{\LD}[2]{\mathsf{L{D}}\pth{#1,#2}}
\newcommand{\SEVN}{\widehat{\lambda_2}}
\newcommand{\NX}[2]{#1\!\left[#2\right]}
\newcommand{\Entry}[2]{#1\!\left[#2\right]}
\newcommand{\NXExt}[3]{#1_{#2}\!\left[#3\right]}
\newcommand{\RangeI}[1]{\llbracket #1 \rrbracket}
\newcommand{\DialX}[1]{\gamma\pth{#1}}
\newcommand{\BiDialX}[1]{{\gamma}_2\pth{#1}}

\newcommand{\EV}[1]{\widehat{\lambda_{#1}}}
\newcommand{\TEV}[1]{\widehat{\lambda_{#1}}}
\newcommand{\EVal}{\widehat{\lambda}}
\newcommand{\Basis}{\mathcal{B}}

\newcommand{\ZProd}[2]{\ensuremath{#1 \, \text{\textcircled{z}} #2}}
\newcommand{\RProd}[2]{\ensuremath{#1 \, \text{\textcircled{r}} #2}}

\newcommand{\OneVec}[1]{\mathbf{1}^{#1}}
\newcommand{\DstrbV}[1]{\mathbf{\mathsf{q}}^{(#1)}}
\newcommand{\DstrbE}[2]{\mathsf{q}^{(#1)}_{#2}}

\newcommand{\DEdges}{\widetilde{\Edges}}

\newcommand{\VV}{\mathsf{V}}%
\newcommand{\EE}{\mathsf{E}}%

\newcommand{\Vertices}{\VV}%
\newcommand{\Edges}{\EE}%

\newcommand{\EdgesX}[1]{\Edges\pth{#1}}
\newcommand{\VerticesX}[1]{V\pth{#1}}

\newcommand{\StringX}[1]{\mathsf{s}\pth{#1}}

\newcommand{\BMtrx}{\mathsf{B}}
\newcommand{\WMtrx}{\mathsf{W}}
\newcommand{\NWMtrx}{\overline{\mathsf{W}}}
\newcommand{\Sequence}{\EuScript{S}}

\newcommand{\med}{\mathrm{med}}
\newcommand{\medX}[1]{\mathrm{med}\pth{#1}}

\newcommand{\medianC}{\nu}%

\newlength{\savedparindent}
\newcommand{\SaveIndent}{\setlength{\savedparindent}{\parindent}}
\newcommand{\RestoreIndent}{\setlength{\parindent}{\savedparindent}}

\newcommand{\NNN}{\mathbb{N}}%
\providecommand{\NN}{\mathbb{N}}%
\AtBeginDocument{%
   \renewcommand{\Re}{\mathbb{R}}%
}
%\shogi
\newcommand{\reals}{\Re}%
\newcommand{\ZZ}{\mathbb{Z}}%

\newcommand{\level}{\mathop{\mathrm{level}}}

\providecommand{\TPDF}[2]{\texorpdfstring{#1}{#2}}

\newcommand{\Hint}[1]{\emph{[#1]}}
\newcommand{\Arr}{\mathop{\mathrm{\EuScript{A}}}}

\newcommand{\Array}{\mathcal{T}}%

\newcommand{\dirEdge}[2]{\left({{#1} \rightarrow {#2}} \right)}

%%%%%%%%%%%%%%%%%%%%%%%%%%%%%%%%%%%%%%%%%%%%%%%%%%%%%%%%%%%%%%%%%%%%%%
% Emphasize stuff...
\definecolor{blue25}{rgb}{0, 0, 11}
\newcommand{\emphic}[2]{%
%   \BADVERYBAD%
   \textcolor{blue25}{%
      \bfseries{\emph{#1}}}%
   \index{#2}}
\newcommand{\emphi}[1]{\emphic{#1}{#1}}
\newcommand{\emphOnly}[1]{\textcolor{blue25}{\bfseries{\emph{#1}}}}
\providecommand{\emphiC}[1]{\emphic{#1}{#1}}


\definecolor{almostblack}{rgb}{0, 0, 0.3}
\providecommand{\emphw}[1]{{\textcolor{almostblack}{\emph{#1}}}}%

%
%%%%%%%%%%%%%%%%%%%%%%%%%%%%%%%%%%%%%%%%%%%%%%%%%%%%%%%%%%%%%%%%%%%%%%

\newcommand{\volX}[1]{\mathrm{vol}\pth{#1}}

\newcommand{\Interval}{\EuScript{I}}
\newcommand{\IntervalA}{\EuScript{I}}

\newcommand{\Exp}{\mathrm{Exp}}

\let\Ball\undefined
\newcommand{\Ball}{\mathbf{b}}
\newcommand{\hdim}{d}
\newcommand{\Domain}{\mathcal{D}}
\newcommand{\points}[1]{[\bfseries{#1 points}]}
\newcommand{\Event}{\EuScript{E}}%
\newcommand{\EventB}{\EuScript{B}}%
\newcommand{\EventF}{\EuScript{F}}%
\newcommand{\EventG}{\EuScript{G}}%

\newcommand{\distCmd}[1]{\left\| {#1}  \right\|}
\newcommand{\distY}[2]{\distCmd{#1 - #2}}

\newcommand{\NormalC}{\mathbf{N}}
\newcommand{\NormalY}[2]{\NormalC\pth{#1, #2}}%
\newcommand{\NormalDistOneDim}{\NormalC}
\newcommand{\NormalDistX}[1]{\NormalC^{#1}}
\DeclareMathOperator{\lca}{lca}
\newcommand{\RangeSpace}{\mathsf{S}}%
\newcommand{\RangeSpaceA}{\mathsf{T}}

\newcommand{\range}{\mathbf{r}}
\newcommand{\rangeA}{\mathbf{r}'}

\newcommand{\BadProb}{\varphi}

\newcommand{\EntropyChar}{\mathbb{H}}
\newcommand{\Entropy}[1]{\EntropyChar\pth{#1}}
\newcommand{\EntropyX}[1]{\EntropyChar\pth[]{#1}}

\newcommand{\Matrix}{\mathsf{M}}
\newcommand{\MatrixB}{\mathsf{B}}
\newcommand{\MatrixC}{\mathsf{C}}
\newcommand{\MatrixD}{\mathsf{D}}

\renewcommand{\th}{th\xspace}

\newcommand{\Family}{\mathcal{F}}
\newcommand{\FamilyB}{\mathcal{G}}
\newcommand{\FamilyC}{\mathcal{H}}
\newcommand{\FMSC}{\EuScript{Y}\index{metric space}}

\newcommand{\median}{\overline{\mathsf{m}}}
\newcommand{\QuickSelect}{\Algorithm{Quick{}Select}\xspace}

\providecommand{\row}{\mathbf{r}}%
\newcommand{\rowc}{{r}}%

\newcommand{\vecb}{\mathbf{b}}

\newcommand{\vecA}{\mathbf{v}}
\newcommand{\vecB}{\mathbf{u}}
\newcommand{\vecRand}{\mathbf{r}}

\newcommand{\vecH}{\mathbf{h}}
\newcommand{\vecV}{\mathbf{v}}

\newcommand{\vecX}{\mathbf{x}}
\newcommand{\vecY}{\mathbf{y}}

\newcommand{\DistBinom}[2]{\mathrm{Bin}\pth{#1,#2}}

\newcommand{\DistGeom}[1]{\mathrm{Geom}\pth{#1}}

\newcommand{\Sample}{\mathsf{R}}

\newcommand{\lcm}{\mathop{\mathrm{lcm}}}
\newcommand{\bdiv}{\mathop{\mathrm{div}}}%
\newcommand{\divides}{\mid}%
\newcommand{\poly}{\mathrm{poly}}%

\newcommand{\Group}{\mathcal{G}} 
\newcommand{\GroupA}{\mathcal{H}} 

\newcommand{\algGCD}{\Algorithm{EuclidGCD}\xspace}%

\newcommand{\gOp}{\times}%
\newcommand{\gIdent}{\mathsf{i}}%

\newcommand{\ordX}[2][\!]{\mathop{\mathrm{o{r}d}}\pth[#1]{#2}}

\newcommand{\lSym}[2]{%
   \left(
       \ds {#1}\divides{#2}%
   \right) %
}
\newcommand{\jSym}[2]{\left\llbracket%
       \ds {#1}\divides{#2}%
       \right\rrbracket%
}
\newcommand{\Line}{\ell}

\newcommand{\LevelChar}{\mathbb{L}}
\newcommand{\AtMostKLevelChar}{\LevelChar_{\leq k}}
\newcommand{\AtMostKLevel}[1]{\LevelChar_{\leq k} \pth{ #1}}

\newcommand{\LineSet}{\mathsf{L}}%
\newcommand{\ArrX}[1]{\EuScript{A}\pth{#1}}%

\providecommand{\Matousek}{Matou{\v s}ek\xspace}
\providecommand{\Erdos}{Erd{\H o}s\xspace}
\providecommand{\Lovasz}{Lov\'asz\xspace}

\newcommand{\ProblemC}[1]{{\textsf{\textcolor{red}{#1}\index{problem!#1}}}}
\DefineNamedColor{named}{OliveGreen} {cmyk}{0.64,0,0.95,0.40}
\providecommand{\ComplexityClass}[1]{{{\textcolor[named]{OliveGreen}{%
            \textsc{#1}}}}}
\providecommand{\NPHard}{{\ComplexityClass{NP-Hard}}%
   \index{NP!hard}\xspace}
\providecommand{\NPComplete}{\ComplexityClass{NP-Complete}%
   \index{NP!complete}\xspace}
\providecommand{\APXHard}{\ComplexityClass{APX-Hard}%
   \index{APX!Hard}\xspace}
\providecommand{\NPC}{\ComplexityClass{NPC}%
   \index{NP!complete}\xspace}
\providecommand{\NP}{\ComplexityClass{NP}%
   \index{NP}\xspace}
\providecommand{\POLYT}{\ComplexityClass{P}\xspace}

\newcommand{\envFrame}[1]{\begin{center}\fbox{\begin{minipage}[c]{0.75\textwidth}{#1}\end{minipage}}\end{center}}

\newcommand{\mysqrt}[1]{\ts\!\!\sqrt{#1}}
\newcommand{\SegSet}{\mathsf{S}}
\newcommand{\ArrVD}[1]{\Arr^{|} \pth{#1}}
\newcommand{\RIC}{\textsf{RIC}\index{random incremental
      construction}\xspace}
\newcommand{\ArrB}{\EuScript{B}}

\newcommand{\TreeOpt}{\EuScript{T}_\mathrm{opt}}
\newcommand{\TreeOptA}{\EuScript{T}_\mathrm{opt}'}

\newcommand{\DepthExt}[2]{\mathrm{depth}_{#1}\pth{#2}}
\newcommand{\depthX}[1]{\mathrm{depth}\pth{#1}}

\newcommand{\TreeA}{\EuScript{U}}
\newcommand{\TreeB}{\EuScript{V}}

\newcommand{\Trap}{\sigma}
\newcommand{\TrapA}{\sigma'}
\newcommand{\TrapB}{\sigma''}

\newcommand{\SampleA}{\mathsf{R}'}

\newcommand{\AllRegions}{\mathcal{T}}
\newcommand{\RegionsSet}[1]{\mathcal{F}\pth{#1}}
\newcommand{\RegionsHeavy}[2]{\mathcal{F}_{\!\geq#1}\pth{#2}}
\newcommand{\NumOfRegionsZero}[1]{\ExChar \!f\pth{#1}}
\newcommand{\NumOfRegions}[2]{\ExChar\! f_{\!\geq #1}\pth[]{#2}}
\newcommand{\ObjSet}{\mathsf{S}}

\newcommand{\WeightX}[2][\!]{\omega \pth[#1]{#2}}
\newcommand{\Edge}{\mathsf{e}}

\newcommand{\Face}{\mathsf{f}}

\newcommand{\DefSet}[1]{D\pth{#1}}
\newcommand{\KillSet}[1]{K\pth{#1}}

\newcommand{\CSProbFull}[4]{
   \,\rho_{\!\!\stackrel{\rule[-.005cm]{0cm}{0.1cm}}{#1,
       #4}}\pth{\MakeSBig #2,#3}}

\newcommand{\disk}{D}
\newcommand{\diskA}{D'}

\newcommand{\IntervalSet}{\widehat{\mathcal{I}}}

\newcommand{\GroundSet}{\textsf{X}}
\newcommand{\FGroundSet}{\mathsf{x}}
\newcommand{\RangeSet}{{\mathcal{R}}}
\newcommand{\RangeSetA}{\mathcal{R}'}
\newcommand{\MeasureChar}{\overline{m}}
\newcommand{\Measure}[1]{\MeasureChar\pth{#1}}
\newcommand{\sMeasure}[1]{\overline{s}\pth{#1}}

\newcommand{\SetA}{S}%
\newcommand{\SetB}{T}%
\newcommand{\SetC}{N}%

\newcommand{\VCProj}[2]{#1_{|#2}}
\newcommand{\VCDim}{\mathrm{dim}_{\VC}}
\newcommand{\Dim}{\delta}
\newcommand{\DimA}{\delta'}
\newcommand{\PntSetB}{{{\mathsf{Q}}}}
\newcommand{\PntSetC}{{{\mathsf{R}}}}
\newcommand{\Growth}[2]{{\EuScript{G}}_{#1}\pth{#2}}
\newcommand{\ShatterFunc}[1]{\pi_{#1}}
\newcommand{\Distribution}{\EuScript{D}}
\newcommand{\Error}{\EuScript{E}}

\newcommand{\bd}{{\partial}}%
\newcommand{\eps}{{\varepsilon}}%

\newcommand{\clmlab}[1]{\label{claim:#1}}
\newcommand{\clmref}[1]{\HLink{Claim}{claim:#1}}

\newcommand{\CHX}[1]{{\mathcal{CH}}\pth{#1}}

\newcommand{\ts}{\hspace{0.6pt}}

\def\qedsymbol{\myqedsymbol}%

\newcommand{\qedsign}{\hfill{\hfill\myqedsymbol}}
\newcommand{\qedsignBig}{\hfill{\hfill\rule{4mm}{4mm}}}
\newcommand{\QED}{{\hfill{\hfill\qedsymbol}}}

\newcommand{\num}{\alpha}

\newcommand{\NEW}[1]{#1}

\newcommand{\ThreeFiguresExt}[9]{%
   \begin{tabular}{ccc}
     \begin{minipage}{0.31\linewidth}
         \begin{center}
             {\IncludeGraphics[#1]{#2}}
         \end{center}
     \end{minipage}
     &
       \begin{minipage}{0.31\linewidth}
           \centerline{\IncludeGraphics[#4]{#5}}
       \end{minipage}
     &
       \begin{minipage}{0.31\linewidth}
           \centerline{\IncludeGraphics[#7]{#8}}
       \end{minipage}
     \\
     ~\\[-0.2cm]
     {#3} & {#6} & {#9}
   \end{tabular}}

\newcommand{\IncludeGraphics}[2][]{%
   \typeout{}%
   \typeout{Graphics: #2}%
   \typeout{\ includegraphics[#1]{#2}}%
   \includegraphics[#1]{#2}
   \typeout{}%
}

%    -----------------------------------------------------------------
%    -----------------------------------------------------------------
%    -----------------------------------------------------------------
%    Handling references
%    -----------------------------------------------------------------
%    -----------------------------------------------------------------
%    -----------------------------------------------------------------
\newcommand{\HLinkShort}[2]{\hyperref[#2]{#1\ref*{#2}}}
\newcommand{\HLink}[2]{\hyperref[#2]{#1~\ref*{#2}}}
\newcommand{\HLinkPage}[2]{\hyperref[#2]{#1~\ref*{#2}%
      $_\text{p\pageref{#2}}$}}

\newcommand{\HLinkSuffix}[3]{\hyperref[#2]{#1\ref*{#2}{#3}}}
\newcommand{\HLinkPageSuffix}[3]{\hyperref[#2]{#1\ref*{#2}%
      #3$_\text{p\pageref{#2}}$}}

\newcommand{\figlab}[1]{\label{fig:#1}}
\newcommand{\figref}[1]{\HLink{Figure}{fig:#1}}

\newcommand{\exelab}[1]{\label{exer:#1}}
\newcommand{\exeref}[1]{\HLink{Exercise}{exer:#1}}

\providecommand{\deflab}[1]{\label{def:#1}}
\newcommand{\defref}[1]{\HLink{Definition}{def:#1}}
\newcommand{\defrefY}[2]{\hyperref[def:#2]{#1}}
\newcommand{\defrefpage}[1]{\HLinkPage{Definition}{def:#1}}

\newcommand{\etal}{\textit{et~al.}\xspace}

\newcommand{\obslab}[1]{\label{observation:#1}}
\newcommand{\obsref}[1]{\HLink{Observation}{observation:#1}}
\newcommand{\obsrefY}[2]{\hyperref[observation:#2]{#1}}

\newcommand{\exmlab}[1]{\label{example:#1}}
\newcommand{\exmref}[1]{\HLink{Example}{example:#1}}

\newcommand{\remlab}[1]{\label{rem:#1}}
\newcommand{\remref}[1]{\HLink{Remark}{rem:#1}}%

\newcommand{\chaplab}[1]{\label{chap:#1}}

\newcommand{\itemlab}[1]{\label{item:#1}}
\newcommand{\itemref}[1]{\HLinkSuffix{}{item:#1}{}}

\newcommand{\tedlab}[1]{\label{tedium:#1}}
\newcommand{\tedref}[1]{\HLink{Tedium}{tedium:#1}}%

\newcommand{\factlab}[1]{\label{fact:#1}}%

\newcommand{\alglab}[1]{\label{algorithm:#1}}

\newcommand{\net}{\mathcal{N}}%


\newcommand{\lemlab}[1]{\label{lemma:#1}}
\newcommand{\lemref}[1]{\HLink{Lemma}{lemma:#1}}%
\newcommand{\lemrefY}[2]{\hyperref[lemma:#2]{#1}}
\newcommand{\lemrefpage}[1]{\HLinkPage{Lemma}{lemma:#1}}
\newcommand{\lemrefshort}[1]{\HLinkShort{L}{lemma:#1}}

\newcommand{\problab}[1]{\label{prob:#1}}%
\newcommand{\probref}[1]{\HLink{Problem}{prob:#1}}%

\newcommand{\eqlab}[1]{\label{equation:#1}}
\newcommand{\Eqrefbare}[1]{\HLinkSuffix{(}{equation:#1}{)}}
\newcommand{\Eqref}[1]{\HLinkSuffix{Eq.~(}{equation:#1}{)}}
\newcommand{\Eqrefpage}[1]{\HLinkPageSuffix{Eq.~(}{equation:#1}{)}}

\newcommand{\seclab}[1]{\label{sec:#1}}
\newcommand{\secref}[1]{\HLink{Section}{sec:#1}}
\newcommand{\secrefY}[2]{\hyperref[sec:#2]{#1}}
\newcommand{\secrefpage}[1]{\HLinkPage{Section}{sec:#1}}

\newcommand{\thmlab}[1]{{\label{theo:#1}}}
\newcommand{\thmref}[1]{\HLink{Theorem}{theo:#1}}
\newcommand{\thmrefY}[2]{\hyperref[theo:#2]{#1}}
\newcommand{\thmrefpage}[1]{\HLinkPage{Theorem}{theo:#1}}

\newcommand{\corlab}[1]{\label{cor:#1}}
\newcommand{\corref}[1]{\HLink{Corollary}{cor:#1}}%
\newcommand{\correfY}[2]{\hyperref[cor:#2]{#1}}

\newcommand{\Yend}{Y_{\mathrm{end}}}%


%    -----------------------------------------------------------------
%    -----------------------------------------------------------------
%    -----------------------------------------------------------------
%    Handling references
%    -----------------------------------------------------------------
%    -----------------------------------------------------------------
%    -----------------------------------------------------------------

\newcommand{\ceil}[1]{\left\lceil {#1} \right\rceil}
\newcommand{\floor}[1]{\left\lfloor {#1} \right\rfloor}
\newcommand{\brc}[1]{\left\{ {#1} \right\}}
\newcommand{\pth}[2][\!]{\mleft({#2}\mright)}%
\newcommand{\brk}[1]{\mleft[{#1}\mright]}%
\newcommand{\pbrc}[2][\!\!]{#1\left[ {#2} \bigr. \right]}
\newcommand{\pbrcS}[1]{\left[ {#1} \right]}

\newcommand{\dotprod}[2]{{#1}\cdot{#2}}
\newcommand{\DotProd}[2]{\permut{{#1},{#2}}}
\newcommand{\permut}[1]{\left\langle {#1} \right\rangle}

\newcommand{\MakeBig}{}
\newcommand{\MakesBig}{}
\newcommand{\MakeSBig}{}
\newcommand{\MakeVBig}{}

\newcommand{\DistChar}{\mathsf{d}}%
\newcommand{\DistSetY}[2]{\DistChar\pth{#1,#2}}
\newcommand{\distG}[2]{\DistChar_\Graph\pth{#1, #2}}
\newcommand{\DistMtrCharX}[1]{\DistChar_{#1}}

\newcommand{\dY}[2]{\DistChar\pth{#1, #2}}%
\newcommand{\dGY}[2]{\DistChar_\Graph\pth{#1, #2}}%

\newcommand{\cset}{\mathcal{C}}%

\newcommand{\MTR}{\mathcal{M}}
\newcommand{\DistMtr}[2]{\DistChar\pth{#1,#2}}

\newcommand{\dist}[1]{\left\| {#1} \right\|}
\newcommand{\norm}[1]{\left\| {#1} \right\|}

\providecommand{\ds}{\displaystyle}
\newcommand{\sep}[1]{\,\left|\, {#1} \bigr.\right.}%
\newcommand{\sepw}[1]{\left|\, {#1} \right.}%

\newcommand{\ProbExt}[2]{\mathop{\mathbf{Pr}}_{#1}\pbrcx{#2}}
\newcommand{\ProbLTR}{\mathbb{P}}%

\newcommand{\Prob}[1]{\mathop{\ProbLTR} \mleft[ #1 \mright]}%
\newcommand{\ProbCond}[2]{\mathop{\ProbLTR}\!\left[%
       #1 \;\middle\vert\; #2 \right]}
\newcommand{\ProbY}[2]{\ProbLTR^{}_{\!#1}\pbrc{#2}}

\newcommand{\ExChar}{\mathbb{E}}%
\newcommand{\ExSym}{\mathop{\ExChar}}%
\newcommand{\Ex}[2][\!]{\ExSym#1\pbrcx{#2}}
\newcommand{\ExExt}[2]{\ExSym_{#1}\!\pbrcx{#2}}
\newcommand{\ExExtS}[2]{\ExSym_{\substack{#1}}\!\pbrcx{#2}}
\newcommand{\ExCond}[2]{\ExSym\!\left[%
       #1 \;\middle\vert\; #2 \right]}

\newcommand{\VarSymb}{\mathop{\mathbb{V}}}
\newcommand{\Var}[1]{\mathop{\mathbb{V}}\!\pbrcx{#1}}

\newcommand{\pbrcx}[1]{\left[ {#1} \right]}%
\newcommand{\nfrac}[2]{{#1}/{#2}}
\newcommand{\cardin}[1]{\left| {#1} \right|}%

\newcommand\CCancel[2][black]{{\renewcommand\CancelColor{\color{#1}}\cancel{#2}}}

\newcommand{\IncGraphPage}[4][]{%
   \IncludeGraphics[page=#4,#1]{\File{#2/#3}}%
}

\newenvironment{enumEx}{
   \begin{compactenumA}
    }{\end{compactenumA}}

\newcommand{\cl}{{\mathrm{cl}}}

\usepackage[amsmath,thmmarks]{ntheorem}
\theoremseparator{.}%

\newtheorem{theorem:unnumbered}{Theorem}[section]
\newtheorem{theorem}{Theorem}[section]

\newtheorem{observation}[theorem]{Observation}
\newtheorem{conjecture}[theorem]{Conjecture}
\newtheorem{lemma}[theorem]{Lemma}
\newtheorem{corollary}[theorem]{Corollary}
\newtheorem{claim}[theorem]{Claim}
\newtheorem{assumption}[theorem]{Assumption}
\newtheorem{tedium}[theorem]{Tedium}

\theoremstyle{plain}%
\theoremheaderfont{\sffamily}
\theorembodyfont{\upshape}%

\newtheorem{defn}[theorem]{Definition}
\newtheorem{definition}[theorem]{Definition}

\newtheorem*{remark_unnumbered}[theorem:unnumbered]{Remark}%
\newtheorem*{remarks}[theorem]{Remarks}%
\newtheorem{remark}[theorem]{Remark}%
\newtheorem{fact}[theorem]{Fact}%
\newtheorem{example}[theorem]{Example}
\newtheorem{exercise}[theorem]{Exercise}
\newtheorem{problem}[theorem]{Problem}
\newtheorem{algorithm}[theorem]{Algorithm}%
\newtheorem{xca}[theorem]{Exercise}
\newtheorem{openproblem}[theorem]{Open Problem}
\newtheorem{exercise_h}[theorem]{Exercise}

% Proof environments
\theoremheaderfont{\em}%
\theorembodyfont{\upshape}%
\theoremstyle{nonumberplain}
\theoremseparator{}
\theoremsymbol{\myqedsymbol}
\newtheorem{proof}{Proof:}
\newtheorem{proofextin}{Proof}
\newtheorem{altproof}{Alternative proof:}

\theoremsymbol{}
\newtheorem{proofnoqed}{Proof:}

\theoremsymbol{\myqedsymbol}
\theorempreskip{0pt}%
\theorempostskip{0pt}%
\newtheorem{proofof}{Proof of\!}%

\newenvironment{proofExt}[1]{{\emph{#1:}}}{\hfill{\hfill\myqedsymbol}}
\newenvironment{proofnoproof}{{}}{{\hfill{\hfill\myqedsymbol}}}

%%%%%%%%%%%%%%%%%%%%%%%%%%%%%%%%%%%%%%%%%%%%%%%%%%%%%%%%%%%%%%%%%%%%
%%%%%%%%%%%%%%%%%%%%%%%%%%%%%%%%%%%%%%%%%%%%%%%%%%%%%%%%%%%%%%%%%%%%
% Macros to embed figures in text

\newcommand{\EBRY}[2]{#1_{\langle#2\rangle}}%

\newcommand{\NiceFrac}[2]{%
   {\raisebox{0.2em}{\ensuremath{#1}}\!\!\!}%
   {\rotatebox[origin=c]{-20}{\Big/}}% Rotated /
   {\! \raisebox{-0.2em}{\ensuremath{#2}} }}%

% end - Macros to embed figures in text
%%%%%%%%%%%%%%%%%%%%%%%%%%%%%%%%%%%%%%%%%%%%%%%%%%%%%%%%%%%%%%%%%%%% 
%%%%%%%%%%%%%%%%%%%%%%%%%%%%%%%%%%%%%%%%%%%%%%%%%%%%%%%%%%%%%%%%%%%%

%%%%%%%%%%%%%%%%%%%%%%%%%%%%%%%%%%%%%%%%%%%%%%%%%%%%%%%%%%%%%%%%%%
% Color in math
%\providecommand{\Mh}[1]{{#1}}%
\providecommand{\IfColorMath}[2]{#2}%

\IfFileExists{.latex_printer_friendly}{\def\GenPrinterVer{1}}{}%

\ifx\GenPrinterVer\undefined
   \IfFileExists{.latex_color}{\def\GenColorMath{1}}{}
\else
   \renewcommand{\IfPrinterVer}[2]{#1}%
\fi

\ifx\GenColorMath\undefined
\else
\renewcommand{\IfColorMath}[2]{#1}%
%\renewcommand{\Mh}[1]{{\textcolor{red}{#1}}}%
\fi
% Color in math - end
%%%%%%%%%%%%%%%%%%%%%%%%%%%%%%%%%%%%%%%%%%%%%%%%%%%%%%%%%%%%%%%%%%

\newcommand{\HSetY}[2]{H\pth{#1, #2}}%
\newcommand{\CHSetY}[2]{C\pth{#1, #2}}%
\newcommand{\Holder}{H\"older\xspace}
\newcommand{\OldOm}{{\mathcal{O}}}
\newcommand{\NewOm}{{\mathcal{N}}}

%%%%%%%%%%%%%%%%%%%%%%%%%%%%%%%%%%%%%%%%%%%%%%%%%%%%%%%%%%%%%%%%%%%%%%%
%---------------------------------------------------------------
\newcommand{\IntRange}[1]{\left\llbracket #1 \right\rrbracket}
\newcommand{\IRX}[1]{\left\llbracket #1 \right\rrbracket}%
\newcommand{\IRY}[2]{\left\llbracket #1:#2 \right\rrbracket}

%---------------------------------------------------------------
%%%%%%%%%%%%%%%%%%%%%%%%%%%%%%%%%%%%%%%%%%%%%%%%%%%%%%%%%%%%%%%%%%%%%%%

\newcommand{\obj}{f}%
\newcommand{\Wg}{\omega}%
\newcommand{\TotalW}{W}%

\newcommand{\rankY}[2]{{#1}_{ \permut{#2}}}%%
\newcommand{\areaX}[1]{\mathrm{area}\pth{#1}}%%

\newcommand{\mybinom}[2]%
{\Bigl(\begin{array}{@{}c@{}}#1\\#2\end{array}\Bigr)}

\providecommand{\Set}[2]{\left\{ #1 \;\middle\vert\; #2 \right\}}
\newcommand*{\biga}[1]{\vcenter{\hbox{\scalebox{1.5}{\ensuremath#1}}}}
\newcommand{\CNF}{\TermI{CNF}\xspace}%
\newcommand{\Tree}{T}%

\DeclareSymbolFont{sfletters}{OML}{cmbrm}{m}{it}

\DeclareMathSymbol{\salpha}{\mathord}{sfletters}{"0B}

%%%%%%%%%%%%%%%%%%%%%%%%%%%%%%%%%%%%%%%%%%%%%%%%%%%%%%%%%%%%%%%%%%%%

\providecommand{\ProblemInd}[2]{\textsf{\textcolor{red}{#1}\index{Problem!#2}}}
\newcommand{\NPProblemInd}[4]{
   \begin{center}
      \vspace{0.6cm}
       \fbox{
          \begin{minipage}{0.9\linewidth}
              \vspace{-0.6cm} \hspace{-0.35cm}
              {\large \bfseries \ProblemInd{#1}{#2}}
              \vspace{+0.2cm}%

              \noindent{\textcolor{blue}{\bfseries \textsf{Instance}}:} {#3}
              
              \noindent{\textcolor{blue}{\bfseries \textsf{Question}:}} {#4}
          \end{minipage}
       }
   \end{center}
   \smallskip
   }

\newcommand{\NPProblemOpt}[4]{\NPProblemInd{#1}{#2}{#3}{#4}}

%%%%%%%%%%%%%%%%%%%%%%%%%%%%%%%%%%%%%%%%%%%%%%%%%%%%%%%%%%%%%%%%%%%%

\newcommand{\Hastad}{H{\aa}stad\xspace}

\newcommand{\ANNbr}{\Term{ANN}\index{approximate nearest neighbor}\xspace}
\newcommand{\LSH}{\TermI{L{S}H}\xspace}

\newcommand{\DS}{\EuScript{D}} % DataStructure
\newcommand{\DSC}{\mathsf{D}} % DataStructure

\newcommand{\distH}[2]{\DistChar_{H}\pth{#1, #2}}
\newcommand{\distHY}[2]{\DistChar_{H}\pth{#1, #2}}

\newcommand{\DSAprxNearNbr}{\ensuremath{\DataStructure{{\approx%
         \text{Near{}}}}}%
   \index{near neighbor!data-structure!approximate}%
}
\newcommand{\DataStructure}[1]{\ensuremath{\DS_{#1}}}
\newcommand{\DistNN}[2]{\DistChar\pth{#1,#2}}
\newcommand{\pntc}{{p}}%

\newcommand{\seq}{{M}}%
\newcommand{\cTimes}{\beta}%

\newcommand{\DSTimes}{L}%
\newcommand{\rad}{r}%

\newcommand{\HypercubeX}[1]{\Hypercube^{#1}}
\newcommand{\kFunctions}[2]{\mathrm{combine}\pth{#1,#2}}
\newcommand{\WEqualChar}{\Bumpeq}
\newcommand{\WEqual}[2]{{#1}\WEqualChar{#2}}
\newcommand{\FamilyAmp}{\FamilyC_{\WEqualChar}}
\newcommand{\PSPACE}{\ComplexityClass{PSPACE}\xspace}%

\newcommand{\PS}{P}%

\newcommand{\PSetFar}{\ensuremath{\PntSet^{}_{\geq}\xspace}}%
\newcommand{\DA}{D}%
\newcommand{\uniqX}[1]{\mathrm{uniq}\pth{#1}}%
\newcommand{\diffX}[1]{\mathop{}\!\mathrm{d}#1}

\newcommand{\Frq}{F}%
\newcommand{\frq}{f}%
\newcommand{\STRM}{\mathcal{S}}%
\newcommand{\strm}{{s}}%

\newcommand{\Calinescu}{C\u{a}linescu\xspace}


\providecommand{\Turan}{Tur\'an\xspace}
\newcommand{\fGrowX}[1]{G^{}_{\Dim}\pth{#1}}
\newcommand{\CompressFunc}{\mathtt{Compress}}
\newcommand{\RSA}{\Term{RSA}\xspace}
\newcommand{\MST}{\Term{MST}\xspace}
\providecommand{\CClassPP}{\CClass{P{}P}\xspace}
\providecommand{\ID}{\mathrm{id}}

\newcommand{\EncodeFunc}{\text{\ttfamily{Enc}}}
\newcommand{\DecodeFunc}{\text{\ttfamily{Dec}}}
\newcommand{\discX}[1]{\mathrm{disc}\pth{#1}}%

\newcommand{\propertyC}{\psi}%
\newcommand{\propertyX}[1]{\propertyC\pth{#1}}%
\newcommand{\numC}{\Gamma}%

\newcommand{\degreeX}[1]{\mathrm{degree}\pth{#1}}

\newcommand{\ind}{m}

\providecommand{\nequiv}{\not\equiv}
\providecommand{\sqcupplus}{\sqcup^+}

\newcommand{\AlgorithmI}[1]{{\protect\color{RedViolet}{#1}}}

\usepackage{marginnote}
\usepackage{ragged2e} 
% Short, subtle snark macro
\setlength{\marginparwidth}{1.5in}
\setlength{\marginparsep}{10pt}
\newcommand{\snark}[1]{%
   \marginnote{\raggedright\scriptsize\sffamily\color{gray}#1}%
}

\newcommand{\binsY}[2]{\sqcup_{\geq #1}\pth{#2}}
\newcommand{\ballsY}[2]{\text{\Football}_{\geq #1}\pth{#2}}

\newcommand{\BLX}[1]{{{\sqcupplus{}}_{\geq #1}}}

\MakePerPage{footnote} % Resets footnote counter on every page

\allowdisplaybreaks % Allows page breaks inside multiline math environments

\let\Mh\undefined


% 1. Create custom page style for chapter start pages
\makepagestyle{chapfirstpage}

\makefootrule{chapfirstpage}{\textwidth}{0.4pt}{1ex}

% 2. Define footers: {Left}{Center}{Right}
% For two-sided documents (recto/verso):
\makeoddfoot{chapfirstpage}{\footnotesize Sariel Har-Peled}{\thepage}{\today}
\makeevenfoot{chapfirstpage}{\thepage}{Sariel Har-Peled}{}

% (Leave headers empty on chapter start pages)
\makeoddhead{chapfirstpage}{}{}{}
\makeevenhead{chapfirstpage}{}{}{}

% 3. Alias 'chapter' style to use your new style
\aliaspagestyle{chapter}{chapfirstpage}

\newcommand{\pdf}{\textsf{pdf}\xspace}
\newcommand{\cdf}{\textsf{cdf}\xspace}

% -----------------------------------------------------------------------------
% Bibliography & Document Begin
%-----------------------------------------------------------------------------
% book.bib is a relative symlink to shared/book.bib in each build directory.
\addbibresource{./book.bib}%
\makeindex
\begin{document}
}{\input{prefix_latex}}%
  \fi
\fi
\makeatother
%
\ChapterNumPage{57}{552}

\Chapter{Algorithmic Version of \Lovasz Local Lemma}
\ChapterPrefix{}
\ChapterGen{}



% Leave room for the two-line title above the standalone epigraph.
\StandAloneMode{\vspace*{1.6cm}}%
\SaveIndent%
\QuoteExt{%
   \RestoreIndent%
   As for me, they took away my toy merchant, wishing with him to
   banish all toys from the world.%
}%
{The tin drum}%
{Gunter Grass}


\newcommand{\Vars}{\EuScript{P}}%
\newcommand{\BEvent}{{B}}%
\newcommand{\BEventA}{{C}}%
\newcommand{\BEventSet}{\mathcal{B}}%
\newcommand{\SetVars}{S}%
\newcommand{\vblX}[1]{\mathrm{v{b}l}\pth{#1}}%
\newcommand{\NbrX}[1]{\Gamma\pth{#1}}%
\newcommand{\NbrCX}[1]{\Gamma_{\!\!+}\pth{#1}}%
\newcommand{\ELogChar}{L}%
\newcommand{\ELogX}[1]{\ELogChar\pth{#1}}%
\newcommand{\WTree}{\tau}%
\newcommand{\WTreeX}[1]{\WTree\pth{#1}}%
\newcommand{\LabelX}[1]{\left[ #1 \right]}%
\newcommand{\nodeDepthX}[1]{\mathrm{depth}\pth{#1}}%
\newcommand{\WSum}{W}%
\newcommand{\NResX}[1]{N_{#1}}%
\newcommand{\FixAlg}{\Algorithm{Fix}\xspace}%
\newcommand{\SolveAlg}{\Algorithm{SolveSAT}\xspace}%
\newcommand{\MTAlg}{\Algorithm{MTResample}\xspace}%

\section{Introduction}

Color the vertices $v_1,\ldots,v_n$ of a cycle independently and
uniformly from a palette of $q$ colors. Let $B_i$ be the bad event
that $v_i$ and $v_{i+1}$ have the same color, with indices taken modulo
$n$. Whenever $B_i$ occurs, give its two endpoints fresh independent
colors from the same palette. Which other edges can change from
properly colored to monochromatic in this step?

Only the two edges on either side can change. Thus, if we draw one node
for each bad event and join events that share a vertex color, we get
another cycle, as in \figref{mt:cycle}. This is the dependency graph.

\begin{figure}[htb]
    \centering
    \includegraphics[width=.72\linewidth]{\File{figs/cycle-coloring.pdf}}
    \caption{The red edge $v_1v_2$ gives the violated event $B_1$.
       Its neighbors in the dependency graph are $B_5$ and $B_2$.}
    \figlab{mt:cycle}
\end{figure}

The \Lovasz local lemma (LLL) shows that if bad events are sufficiently
unlikely relative to their neighborhoods, then with positive probability
none occurs. That probability can still be exponentially small. The
cycle suggests a different algorithm: pick any bad event that occurs,
resample its variables, and repeat. How many such repairs should we expect?

We avoid stating the \Lovasz local lemma explicitly here, because we
derive an algorithmic version of it here.

\subsection{The setting}

Let $\Vars$ be a finite collection of mutually independent random
variables. We are interested in (bad) events that are determined by
these variables. Specifically, for an event $\BEvent$, there is a set
of variables $\SetVars \subseteq \Vars$ that determine if $\BEvent$
happens. The minimal such set of variables is denoted by
$\vblX{\BEvent}$ (note that it is unique). We assume that
$\vblX{\BEvent}$ is available, and that one can check efficiently,
given an assignment to the variables, whether $\BEvent$ happens.

Consider a finite family $\BEventSet$ of such (bad) events. The
\emphOnly{dependency graph}\index{dependency graph}
$\Graph = \Graph(\BEventSet)$ has the events of $\BEventSet$ as
vertices, and two distinct events $\BEvent_1,\BEvent_2 \in \BEventSet$
are connected by an edge $\iff$
$\vblX{\BEvent_1} \cap \vblX{\BEvent_2} \neq \emptyset$. Let
$\NbrX{\BEvent}$ be the set of neighbors of $\BEvent$ in $\Graph$, and
let
\begin{equation*}
    \NbrCX{\BEvent} = \brc{\BEvent} \cup \NbrX{\BEvent}.
\end{equation*}
The events
outside $\NbrCX{\BEvent}$ are jointly determined by variables disjoint
from $\vblX{\BEvent}$. Hence $\BEvent$ is independent of the entire
collection of those events, so $\Graph$ is a dependency
graph in the sense used by the \Lovasz local lemma.

A specific assignment to the variables
\emphOnly{violates}\index{violates} $\BEvent$ if it causes $\BEvent$
to happen (remember, it is a bad event).

\paragraph*{Task.}
Compute an assignment to the variables of $\Vars$, such that none of
the bad events of $\BEventSet$ happens.

\paragraph*{Algorithm.}
Initially, the algorithm assigns the variables of $\Vars$ random
values (each according to its own distribution). As long as there is
a violated event $\BEvent \in \BEventSet$, the algorithm resamples all
the variables of $\vblX{\BEvent}$ (independently according to their
own distributions) --- this is a
\emphOnly{resampling}\index{resampling} of $\BEvent$. The algorithm
repeats this till no event is violated. See \figref{m:t:alg}.

\begin{figure}[h]
    \centerline{%
       \AlgorithmBox{%
          \MTAlg{}($\Vars, \BEventSet$) \+\\
          \For each $P \in \Vars$ \Do\+\\
          $v_P \leftarrow$ a random sample of $P$ \-\\
          \While $\exists \BEvent \in \BEventSet$ violated by
          $\pth{v_P}_{P \in \Vars}$ \Do\+\\
          pick (any) such violated $\BEvent$\\
          \For each $P \in \vblX{\BEvent}$ \Do\+\\
          $v_P \leftarrow$ a fresh random sample of $P$ \-\-\\
          \Return $\pth{v_P}_{P \in \Vars}$%
       }%
    }
    \caption{The resampling algorithm of Moser and Tardos.}
    \figlab{m:t:alg}
\end{figure}

\paragraph*{Finer details.}
We fix some arbitrary strategy (randomized or deterministic) of how to
pick the next violated event. The analysis below works for any such
strategy, which can even be adversarial (as long as it does not peek
into the future random values).

\bigskip

Our purpose is to prove the following.

\begin{theorem}[Moser and Tardos \cite{mt-cpgll-10}]
    \thmlab{m:t}%
    %
    Let $\Vars$ be a finite set of independent random variables, and
    let $\BEventSet$ be a finite set of (bad) events determined by
    these variables. If there is an assignment
    $x:\BEventSet \rightarrow (0,1)$, such that
    \begin{equation}
        \forall \BEvent \in \BEventSet%
        \qquad %
        \Prob{\BEvent} \leq x(\BEvent) \prod_{\BEventA \in
           \NbrX{\BEvent}} \bigl( 1 - x(\BEventA) \bigr),
        \eqlab{l:l:l:cond}
    \end{equation}
    then \MTAlg resamples any event $\BEvent \in \BEventSet$, in
    expectation, at most $x(\BEvent)/ \bigl(1-x(\BEvent) \bigr)$
    times. Overall, the expected number of resampling steps is at
    most 
    \[
        \sum_{\BEvent\in\BEventSet} \frac{x(\BEvent)}{ 1-x(\BEvent) }.
    \]
     In particular, there exists an assignment for the variables of
    $\Vars$ such that no event in $\BEventSet$ happens.
\end{theorem}

Note that the last part is the (variable version of the) \Lovasz
local lemma itself, so the above gives a new proof of the LLL. Before
proving the general result, we demonstrate the main idea on the
special (but important) case of $k$-SAT, where the proof is
especially short.



%-----------------------------------------------------------------------
\section{Warm-up: entropy compression}
\seclab{entropy:compression}

Let $F$ be a $k$-CNF formula with $n$ variables and $m$ clauses, where
every clause has exactly $k$ distinct variables. For a clause $C$,
let $\NbrCX{C}$ be the set of clauses sharing a variable with $C$
(including $C$ itself). We assume that
$\cardin{\NbrCX{C}} \leq 2^{k-3}$ for all clauses $C$. In particular,
one can specify a clause of $\NbrCX{C}$, given $C$, using $k-3$ bits.

A clause $C$ is \emphOnly{violated}\index{violated} by an assignment if
all its literals are false. The key (trivial) observation is the
following.

\begin{observation}
    \obslab{unique:violating}%
    %
    There is exactly one assignment to the $k$ variables of a clause
    $C$ that violates it. As such, if we know that $C$ is violated,
    then we know the values of its $k$ variables.
\end{observation}

Consider the recursive algorithm of \figref{fix}.

\begin{figure}[h]
    \centerline{%
       \AlgorithmBox{%
          \SolveAlg{}($F$) \+\\
          Assign all variables independent fair bits\\
          \For $C = C_1, \ldots, C_m$ \Do\+\\
          \If $C$ is violated \Then \FixAlg{}($C$) \-\\
          \Return the current assignment \-\\
          \\
          \FixAlg{}($C$) \+\\
          Assign the $k$ variables of $C$ fresh independent fair bits\\
          \While some clause $D \in \NbrCX{C}$ is violated \Do\+\\
          \FixAlg{}($D$) \quad\CodeComment{// Pick the first such $D$}
       }%
    }
    \caption{Moser's algorithm for $k$-SAT.}
    \figlab{fix}
\end{figure}

\begin{lemma}
    \lemlab{fix:invariant}%
    %
    If a call \FixAlg{}($C$) returns, then $C$ is satisfied, and every
    clause that was satisfied before the call is still satisfied. In
    particular, if \SolveAlg terminates, it outputs a satisfying
    assignment.
\end{lemma}

\begin{proof}
    Induct on the number of calls in the completed call's subtree.
    Each recursive call has a smaller subtree. The call returns only
    when all clauses of $\NbrCX{C}$ are satisfied, and $C$ is one of
    them. A clause $D' \notin \NbrCX{C}$ does not share variables
    with $C$, so the resampling of $C$ does not change it, and by
    induction, the recursive calls do not change a satisfied clause
    into a violated one.

    As for \SolveAlg, after the iteration handling $C_i$, the clauses
    $C_1, \ldots, C_i$ are all satisfied.
\end{proof}

\begin{theorem}
    \thmlab{moser:k:sat}%
    %
    The probability that \SolveAlg performs at least $t$ calls to
    \FixAlg is at most $2^{m+1-t}$. In particular, the expected number
    of calls to \FixAlg is at most $m+2$, and $F$ is satisfiable.
\end{theorem}

\begin{proof}
    Think about the algorithm as reading its random bits from a string
    $R$ of $n + tk$ random bits: the first $n$ bits are the initial
    assignment, and each call to \FixAlg reads the next $k$ bits. We
    run the algorithm on $R$, and stop it as soon as it either
    terminates, or the $t$\th call to \FixAlg has resampled its
    variables. Let
    $\mathcal{R}$ be the set of strings for which the algorithm is
    stopped because of the second reason --- the strings where at least
    $t$ calls are made.

    For $R \in \mathcal{R}$ we record the following
    \emphOnly{log}\index{log}:
    \begin{compactenumi}
        \item An $m$-bit vector indicating which clauses of $F$ were
        called from the top level loop of \SolveAlg.

        \item The recursive calls and returns, in execution order:
        a recursive call \FixAlg{}($D$) from inside \FixAlg{}($C$) is
        written as the bit $1$ followed by the $k-3$ bits specifying
        $D$ as a member of $\NbrCX{C}$, and a return from a call is
        written as the bit $0$.

        \item The assignment $A_t$ of the variables at the time the
        algorithm was stopped.
    \end{compactenumi}
    There are at most $t$ calls, and at most $t$ returns, so the
    second part has at most $t(k-2) + t = t(k-1)$ bits. As such, the
    number of different possible logs is at most
    $2^m \cdot 2^{t(k-1)+1} \cdot 2^n$ (the $+1$ is because the second
    part has variable length).

    \smallskip

    \noindent\textbf{Claim:} The log determines $R$.

    Start with the first marked top-level clause. A $1$ identifies
    a child of the current clause, and a $0$ returns to its parent.
    Whenever the stack empties, start the next marked top-level clause.
    This reconstructs the sequence
    $C^{(1)}, C^{(2)}, \ldots, C^{(t)}$ of clauses that \FixAlg was
    called on (in order). Now, just before the $t$\th call, the clause
    $C^{(t)}$ was violated, and thus, by
    \obsref{unique:violating}, we know the values of its variables
    just before the call. Since this call only changed these
    variables, we recover $A_{t-1}$, the assignment before the $t$\th
    call. Furthermore, the $k$ random bits read by the $t$\th call are
    the values of the variables of $C^{(t)}$ in $A_t$. Repeating this
    backward, we recover all the bits read by all the calls, and the
    initial assignment $A_0$. Namely, we recovered $R$.

    \smallskip

    The claim implies that the mapping from $\mathcal{R}$ to logs is
    one-to-one, and thus
    $\cardin{\mathcal{R}} \leq 2^{n + m + t(k-1) + 1}$. As such,
    \begin{equation*}
        \Prob{\text{at least $t$ calls}}%
        =%
        \frac{\cardin{\mathcal{R}}}{2^{n+tk}}%
        \leq%
        \frac{2^{n + m + t(k-1) + 1}}{2^{n+tk}}%
        =%
        2^{m+1-t}.
    \end{equation*}
    The expected number of calls is
    $\sum_{t \geq 1} \Prob{\text{at least $t$ calls}} \leq (m+1) +
    \sum_{j \geq 1} 2^{-j} = m+2$. In particular, the algorithm
    terminates with probability one, and by \lemref{fix:invariant} its
    output is a satisfying assignment.
\end{proof}

\begin{remark}
    The intuition is that every call to \FixAlg consumes $k$ fresh
    random bits, but the log grows only by $\approx k-1$ bits per
    call. A random string can not be compressed (with any decent
    probability), so the algorithm must stop quickly. This is known
    as \emphOnly{entropy compression}\index{entropy compression}.

    The argument used the specific structure of $k$-SAT in a crucial
    way --- that a violated clause reveals its variables. It also only
    gives the symmetric version of the LLL, with a weaker constant
    ($2^{k-3}$ instead of $\approx 2^k/e$). For the general case, we
    need a different way to account for the history of the algorithm.
\end{remark}


%-----------------------------------------------------------------------
\section{Analysis of the general algorithm}

\subsection{The resampling table}

It would be convenient to fix all the randomness used by \MTAlg in
advance. For every variable $P \in \Vars$, let
$P^{(0)}, P^{(1)}, P^{(2)}, \ldots$ be an infinite sequence of
independent samples of $P$, with independent sequences for different
variables. This is the
\emphOnly{resampling table}\index{resampling table}. The algorithm
uses $P^{(0)}$ as the initial value of $P$, and the $j$\th time it
resamples $P$, it sets it to $P^{(j)}$. Clearly, the algorithm
behaves exactly as before. Once the table is fixed (and the random
choices of the selection strategy), the execution of the algorithm is
deterministic.

Let $\ELogX{i} \in \BEventSet$ be the event that was resampled in the
$i$\th iteration of the algorithm, for $i \geq 1$. The sequence
$\ELogX{1}, \ELogX{2}, \ldots$ is the \emphOnly{log}\index{log} of the
execution. It is finite if the algorithm terminates.

By construction, the value of a variable $P$ at the time of the
$i$\th resampling (just before it is performed) is $P^{(a)}$, where
\begin{equation}
    a = \cardin{\brc{ j < i \sep{ P \in \vblX{\ELogX{j}}}}}.
    \eqlab{table:index}
\end{equation}

\subsection{Witness trees}

\begin{defn}
    A \emphOnly{witness tree}\index{witness tree} $\WTree$ is a finite
    rooted tree, where each node $v$ is labeled by an event
    $\LabelX{v} \in \BEventSet$, such that if $u$ is a child of $v$,
    then $\LabelX{u} \in \NbrCX{\LabelX{v}}$. The witness tree is
    \emphOnly{proper}\index{witness tree!proper} if the children of
    every node have distinct labels.
\end{defn}

We identify witness trees up to root- and label-preserving isomorphism.
The order of siblings is immaterial.

\begin{figure}[htb]
    \centering
    \includegraphics[width=.83\linewidth]{\File{figs/witness-tree.pdf}}
    \caption{Subscripts mark resampling times. The backward scan skips $D$,
       which uses only $Z$. At time~2, $A$ shares $X$ with both $C_5$ and
       $A_4$, so it attaches to the deeper node $A_4$.}
    \figlab{mt:witness:example}
\end{figure}

\paragraph*{The witness tree of a resampling.}
Construct $\WTreeX{t}$ by scanning backward from time $t$
(\figref{mt:witness:example}).
\begin{compactenumI}
    \smallskip%
    \item Create a root labeled $\ELogX{t}$ with time $t$.

    \smallskip%
    \item For $i=t-1,\ldots,1$, find a deepest node $v$ with
    $\ELogX{i}\in\NbrCX{\LabelX{v}}$, breaking ties by a fixed rule.
    Give $v$ a child labeled $\ELogX{i}$ with time $i$.
    If no such $v$ exists, skip $i$.
\end{compactenumI}
    \smallskip%
A witness tree $\WTree$
\emphOnly{appears in the log}\index{witness tree!appears in log} if
$\WTree=\WTreeX{t}$ for some $t$, after erasing the times.

\begin{lemma}
    \lemlab{deeper}%
    %
    Let $u$ and $w$ be two nodes of $\WTreeX{t}$, associated with times
    $i<j$, such that $\ELogX{i} \in \NbrCX{\ELogX{j}}$. Then
    $\nodeDepthX{u} > \nodeDepthX{w}$. In particular:
    \begin{compactenumi}
        \item Two distinct nodes at the same depth have distinct
        labels that share no variable. Hence $\WTreeX{t}$ is proper.

        \item If $w \in \WTreeX{t}$ is associated with time $j$, and
        $i < j$ is such that $\ELogX{i} \in \NbrCX{\ELogX{j}}$, then $i$
        is associated with some node of $\WTreeX{t}$.
    \end{compactenumi}
\end{lemma}

\begin{proof}
    The nodes are added to the tree in decreasing order of time. When
    the algorithm handles time $i$, the node $w$ is already in the
    tree, and $\ELogX{i} \in \NbrCX{\LabelX{w}}$. As such, $i$ is
    added to the tree (this implies (ii)), as a child of a node of
    depth at least $\nodeDepthX{w}$. Hence
    $\nodeDepthX{u} \geq \nodeDepthX{w}+1$.

    (i) If $u, w$ are at the same depth and their labels are equal
    or share a variable, then the labels are in each other's
    $\NbrCX{\cdot}$, which contradicts the above. Siblings are at the
    same depth, so $\WTreeX{t}$ is proper.
\end{proof}

\begin{lemma}
    \lemlab{distinct:trees}%
    %
    For any event $\BEvent$, let $\NResX{\BEvent}$ be the number of
    times $\BEvent$ is resampled by the algorithm. Then
    $\NResX{\BEvent}$ is equal to the number of distinct proper
    witness trees with root label $\BEvent$ that appear in the log.
\end{lemma}

\begin{proof}
    Let $t_1 < t_2 < \cdots$ be the times at which $\BEvent$ was
    resampled. Since $\BEvent \in \NbrCX{\BEvent}$, by
    \lemref{deeper} (ii) applied to the root, every time
    $t_\ell < t_r$ is associated with a node of $\WTreeX{t_r}$. Thus,
    $\WTreeX{t_r}$ has exactly $r$ nodes labeled $\BEvent$. As such,
    the trees $\WTreeX{t_1}, \WTreeX{t_2}, \ldots$ are all distinct.
\end{proof}

\subsection{The probability that a witness tree appears}

Recall that a node $v$ of a fixed witness tree has a bad event
$\BEvent = \LabelX{v}$ as its label. Then $\Prob{\Bigl.\LabelX{v}}$
is the probability that $\BEvent$ occurs when each variable in
$\vblX{\BEvent}$ is drawn independently from its original distribution.

%%%
\begin{lemma}[Coupling lemma]
    \lemlab{coupling}%
    %
    Let $\WTree$ be a fixed proper witness tree. Then
    \begin{math}
        \ds \Prob{\WTree \text{ appears in the log}}%
        \leq%
        \prod_{v \in \WTree} \Prob{\Bigl.\LabelX{v}}.
    \end{math}
\end{lemma}

\begin{proof}
    Consider the following \emphOnly{$\WTree$-check} procedure, which
    uses the same resampling table as the algorithm. Order the nodes
    of $\WTree$ by decreasing depth (breaking ties arbitrarily), as
    $v_1, \ldots, v_s$. For each variable $P$, maintain a counter
    $c(P)$, initially zero. Now, for $\ell = 1, \ldots, s$:
    \begin{compactitem}
        \item Check that $\LabelX{v_\ell}$ is violated by the assignment
        $P \mapsto P^{(c(P))}$, for $P \in \vblX{\LabelX{v_\ell}}$.

        \item Increase $c(P)$ by one, for all
        $P \in \vblX{\LabelX{v_\ell}}$.
    \end{compactitem}
    The $\WTree$-check \emphOnly{passes} if all these checks succeed.

    \bigskip

    \noindent\textbf{Probability of passing.} Every check inspects
    only table entries that no previous check inspected (since the
    counters only go up). As such, the checks are determined by
    disjoint sets of independent random values, and they are
    independent. The $\ell$\th check succeeds with probability
    exactly $\Prob{\Bigl.\LabelX{v_\ell}}$. Thus, the $\WTree$-check passes
    with probability $\prod_{v \in \WTree} \Prob{\Bigl.\LabelX{v}}$.

    \bigskip

    \noindent\textbf{If $\WTree$ appears in the log, then the
       $\WTree$-check passes.} So, assume $\WTree = \WTreeX{t}$. Let
    $v$ be a node of $\WTree$, associated with time $i$, and let
    $P \in \vblX{\LabelX{v}}$. When the check handles $v$, the
    counter $c(P)$ is the number of nodes $w$ handled before $v$, such
    that $P \in \vblX{\LabelX{w}}$. We claim that these nodes are
    exactly the nodes associated with the times
    $\brc{j < i \sep{ P \in \vblX{\ELogX{j}}}}$:
    \begin{compactitem}
        \smallskip%
        \item For every time $j < i$ with
        $P \in \vblX{\ELogX{j}}$, we have
        $\ELogX{j} \in \NbrCX{\ELogX{i}}$, so by \lemref{deeper},
        $j$ is associated with a node of $\WTree$ which is deeper than
        $v$, and is thus handled before $v$.

        \smallskip%
        \item Conversely, let $w$ be a node handled before $v$, with
        $P \in \vblX{\LabelX{w}}$. The labels of $v$ and $w$ share
        the variable $P$, so by \lemref{deeper} (i), $w$ is not at
        the same depth as $v$, and it is thus strictly deeper. By
        \lemref{deeper}, the time $j$ associated with $w$ can not be
        larger than $i$ (otherwise, $v$ would be deeper than $w$).
        Hence $j < i$.
    \end{compactitem}
    \smallskip%
    As such, by \Eqref{table:index}, the check of $v$ inspects
    exactly the values the variables of $\vblX{\LabelX{v}}$ had at time
    $i$ in the execution of the algorithm. But $\ELogX{i}$ was
    violated at time $i$ (that is why it was resampled), and the check
    succeeds.

    \smallskip

    We conclude that
    $\Prob{\WTree \text{ appears}} \leq \Prob{\WTree\text{-check
          passes}} = \prod_{v \in \WTree} \Prob{\LabelX{v}}$.
\end{proof}

\subsection{Summing over all witness trees}

\begin{lemma}
    \lemlab{tree:sum}%
    %
    Assume the conditions of \thmref{m:t} hold. For an event
    $\BEvent \in \BEventSet$, let
    \begin{equation*}
        \WSum(\BEvent)%
        =%
        \sum_{\WTree} \prod_{v \in \WTree} \Prob{\Bigl.\LabelX{v}},
    \end{equation*}
    where the summation is over all proper witness trees $\WTree$ with
    root label $\BEvent$. Then
    $\WSum(\BEvent) \leq x(\BEvent)/\pth{1 - x(\BEvent)}$.
\end{lemma}

\begin{proof}
    Let $y(\BEvent) = \frac{x(\BEvent)}{1 - x(\BEvent)}$, and observe
    that $1 + y(\BEvent) = \frac{1}{1 - x(\BEvent)}$.

    For $h \geq 0$, let $\WSum_h(\BEvent)$ be the same sum, restricted
    to trees of depth at most $h$. We prove by induction on $h$ that
    $\WSum_h(\BEvent) \leq y(\BEvent)$. Since $\WSum(\BEvent) =
    \lim_{h \to \infty} \WSum_h(\BEvent)$, this implies the claim.

    For the base case, set $\WSum_{-1}(\BEvent)=0$, the sum over
    the empty family of trees. The claimed bound holds at $h=-1$.

    For $h \geq 0$, a proper witness tree of depth at most $h$ with root
    label $\BEvent$ is specified by a set
    $S \subseteq \NbrCX{\BEvent}$ of labels of the children of the root,
    and for each $\BEventA \in S$, a proper witness tree of depth at
    most $h-1$ with root label $\BEventA$. As such
    \begin{align*}
        \WSum_h(\BEvent)%
        &=%
        \Prob{\BEvent} \sum_{S \subseteq \NbrCX{\BEvent}} \;
        \prod_{\BEventA \in S} \WSum_{h-1}(\BEventA)%
        =%
        \Prob{\BEvent} \prod_{\BEventA \in \NbrCX{\BEvent}} \pth{1 +
           \WSum_{h-1}(\BEventA)}%
        \\&%
        \leq%
        \Prob{\BEvent} \prod_{\BEventA \in \NbrCX{\BEvent}} \pth{1 +
           y(\BEventA)}%
        =%
        \frac{\Prob{\BEvent}}{\prod_{\BEventA \in \NbrCX{\BEvent}} \pth{1 -
              x(\BEventA)}}%
        \\&%
        \leq%
        \frac{x(\BEvent) \prod_{\BEventA \in \NbrX{\BEvent}} \pth{1 -
              x(\BEventA)}}{\pth{1 - x(\BEvent)} \prod_{\BEventA \in
              \NbrX{\BEvent}} \pth{1 - x(\BEventA)}}%
        =%
        y(\BEvent),
    \end{align*}
    by induction, and by \Eqref{l:l:l:cond}.
\end{proof}

\begin{remark}
    The second equality above uses the identity
    \begin{equation*}
        \sum_{S \subseteq U} \prod_{u \in S} a_u = \prod_{u \in U} (1 +
        a_u).
    \end{equation*}
    The last step is the only place where the LLL condition
    \Eqref{l:l:l:cond} is used, and it ``explains'' the strange form
    of this condition: it is exactly what is needed to make the
    induction go through.
\end{remark}

\subsection{Putting things together}

\begin{proof}[Proof of \thmref{m:t}]
    By \lemref{distinct:trees}, \lemref{coupling}, and
    \lemref{tree:sum}, we have
    \begin{equation*}
        \Ex{\NResX{\BEvent}}%
        =%
        \sum_{\WTree} \Prob{ \WTree \text{ appears in the log}}%
        \leq%
        \sum_{\WTree} \prod_{v \in \WTree} \Prob{\LabelX{v}}%
        =%
        \WSum(\BEvent)%
        \leq%
        \frac{x(\BEvent)}{1 - x(\BEvent)},
    \end{equation*}
    where the summation is over all proper witness trees with root
    label $\BEvent$. The total number of resamplings is
    $\sum_{\BEvent} \NResX{\BEvent}$, and the bound on its
    expectation follows by linearity of expectation. In particular,
    the number of resamplings is finite with probability one. Namely,
    the algorithm terminates with probability one, and when it
    terminates, it outputs an assignment that violates no event of
    $\BEventSet$.
\end{proof}

\begin{corollary}
    \corlab{m:t:symmetric}%
    %
    If every event of $\BEventSet$ has probability at most $p$, and
    shares variables with at most $d \geq 1$ other events of
    $\BEventSet$, and $e p (d+1) \leq 1$, then \MTAlg performs, in
    expectation, at most $\cardin{\BEventSet}/d$ resamplings.
\end{corollary}

\begin{proof}
    Set $x(\BEvent) = 1/(d+1)$ for all $\BEvent$. As in the
    non-constructive proof, we have
    $x(\BEvent) \prod_{\BEventA \in \NbrX{\BEvent}} \pth{1 -
       x(\BEventA)} \geq \frac{1}{d+1}\pth{1 - \frac{1}{d+1}}^d >
    \frac{1}{e(d+1)} \geq p$. Now
    $x(\BEvent)/\pth{1-x(\BEvent)} = 1/d$, and the claim follows from
    \thmref{m:t}.
\end{proof}

\begin{example}
    Extend the cycle example to a simple $\Delta$-regular graph with $m$
    edges and $\Delta\ge 2$. Its dependency graph is the line graph:
    each edge gives one bad event, and two events are neighbors when their
    edges share an endpoint.
    For a fixed edge $uv$, once the color of $u$ is exposed, $v$ matches
    it with probability $p=1/q$. Exactly $\Delta-1$ other edges use the
    color of $u$, and another $\Delta-1$ use the color of $v$. Thus the
    variable-overlap dependency graph has degree $d=2(\Delta-1)$.
    By \corref{m:t:symmetric}, if $q\ge e(2\Delta-1)$, the algorithm
    finds a proper coloring after at most $m/[2(\Delta-1)]$ edge
    resamplings in expectation.
\end{example}


%-----------------------------------------------------------------------
\section{Bibliographical notes}

The local lemma is due to Erd\H{o}s and \Lovasz \cite{el-prthr-75}.
The first algorithmic version is due to Beck \cite{b-aallll-91}, who
used a two-phase approach (freeze variables in dangerous clauses, and solve the
remaining small components by brute force), with significant loss in
the parameters --- this is the algorithm described in the chapter on
the probabilistic method. Moser \cite{m-cplll-09} discovered the
resampling algorithm for $k$-SAT and analyzed it via entropy
compression; the presentation in \secref{entropy:compression} follows
expositions by Fortnow and Tao of Moser's argument. Moser and Tardos
\cite{mt-cpgll-10} extended this to the general (asymmetric) LLL,
using witness trees. They also gave a parallel version, and a
derandomization in some cases.

The original paper of Moser and Tardos bounds the sum over witness
trees (\lemref{tree:sum}) by bounding each tree's weight by
$x(\BEvent)/(1-x(\BEvent))$ times the probability that a Galton--Watson
process rooted at $\BEvent$ produces that tree. The
direct induction on the depth used here avoids this, and is shorter.

Kolipaka and Szegedy \cite{ks-mtml-11} showed that the same algorithm
works all the way up to Shearer's bound (the best possible condition
on the dependency graph), and introduced the analysis via stable set
sequences. Harvey and Vondr\'ak \cite{hv-aplll-20} extended the
algorithm beyond independent variables, using
\emphOnly{resampling oracles}\index{resampling oracle}; see also
Achlioptas and Iliopoulos \cite{ai-rwfpo-16}. Haeupler, Saha, and
Srinivasan \cite{hss-ncall-11} bounded the probability of an event
in the output using its original probability and the bad events
sharing its variables. With slack in the LLL condition, they also
gave efficient algorithms for applications with exponentially many events. Giotis \etal
\cite{gkpt-aeclll-17} analyzed a resampling algorithm for acyclic
edge coloring. Knuth's treatment
\cite{k-taocp4b-22} develops the subject via the theory of traces.


%-----------------------------------------------------------------------
\section*{Credits}

This chapter was written by Gemini (November 2026), expanding on an
initial draft by Sariel Har-Peled. OpenAI Astra 6.2 reviewed the proofs
and illustrated the cycle coloring and witness-tree construction.


\ChapterEnd{}
